%% Generated by Sphinx.
\def\sphinxdocclass{jupyterBook}
\documentclass[letterpaper,10pt,english]{jupyterBook}
\ifdefined\pdfpxdimen
   \let\sphinxpxdimen\pdfpxdimen\else\newdimen\sphinxpxdimen
\fi \sphinxpxdimen=.75bp\relax
\ifdefined\pdfimageresolution
    \pdfimageresolution= \numexpr \dimexpr1in\relax/\sphinxpxdimen\relax
\fi
%% let collapsible pdf bookmarks panel have high depth per default
\PassOptionsToPackage{bookmarksdepth=5}{hyperref}
%% turn off hyperref patch of \index as sphinx.xdy xindy module takes care of
%% suitable \hyperpage mark-up, working around hyperref-xindy incompatibility
\PassOptionsToPackage{hyperindex=false}{hyperref}
%% memoir class requires extra handling
\makeatletter\@ifclassloaded{memoir}
{\ifdefined\memhyperindexfalse\memhyperindexfalse\fi}{}\makeatother

\PassOptionsToPackage{booktabs}{sphinx}
\PassOptionsToPackage{colorrows}{sphinx}

\PassOptionsToPackage{warn}{textcomp}

\catcode`^^^^00a0\active\protected\def^^^^00a0{\leavevmode\nobreak\ }
\usepackage{cmap}
\usepackage{fontspec}
\defaultfontfeatures[\rmfamily,\sffamily,\ttfamily]{}
\usepackage{amsmath,amssymb,amstext}
\usepackage{polyglossia}
\setmainlanguage{english}



\setmainfont{FreeSerif}[
  Extension      = .otf,
  UprightFont    = *,
  ItalicFont     = *Italic,
  BoldFont       = *Bold,
  BoldItalicFont = *BoldItalic
]
\setsansfont{FreeSans}[
  Extension      = .otf,
  UprightFont    = *,
  ItalicFont     = *Oblique,
  BoldFont       = *Bold,
  BoldItalicFont = *BoldOblique,
]
\setmonofont{FreeMono}[
  Extension      = .otf,
  UprightFont    = *,
  ItalicFont     = *Oblique,
  BoldFont       = *Bold,
  BoldItalicFont = *BoldOblique,
]



\usepackage[Bjarne]{fncychap}
\usepackage[,numfigreset=1,mathnumfig]{sphinx}

\fvset{fontsize=\small}
\usepackage{geometry}


% Include hyperref last.
\usepackage{hyperref}
% Fix anchor placement for figures with captions.
\usepackage{hypcap}% it must be loaded after hyperref.
% Set up styles of URL: it should be placed after hyperref.
\urlstyle{same}


\usepackage{sphinxmessages}



        % Start of preamble defined in sphinx-jupyterbook-latex %
         \usepackage[Latin,Greek]{ucharclasses}
        \usepackage{unicode-math}
        % fixing title of the toc
        \addto\captionsenglish{\renewcommand{\contentsname}{Contents}}
        \hypersetup{
            pdfencoding=auto,
            psdextra
        }
        % End of preamble defined in sphinx-jupyterbook-latex %
        

\title{BIPN 145 Lab Manual}
\date{May 10, 2026}
\release{}
\author{Ashley Juavinett}
\newcommand{\sphinxlogo}{\vbox{}}
\renewcommand{\releasename}{}
\makeindex
\begin{document}

\pagestyle{empty}
\sphinxmaketitle
\pagestyle{plain}
\sphinxtableofcontents
\pagestyle{normal}
\phantomsection\label{\detokenize{Cover::doc}}


\sphinxstepscope


\chapter{Welcome to BIPN 145}
\label{\detokenize{Introduction:welcome-to-bipn-145}}\label{\detokenize{Introduction::doc}}
\sphinxAtStartPar
This is a course about how we study the nervous system.
You’ll be asked to think, plan, and write like a neuroscientist.

\sphinxAtStartPar
This lab manual describes the set of experiments
you’ll be conducting over the course of the quarter.


\section{This course is for everyone.}
\label{\detokenize{Introduction:this-course-is-for-everyone}}
\sphinxAtStartPar
Neuroscience is a diverse discipline, with folks from many different academic and personal backgrounds. We’ll be working together to create an equitable and inclusive environment of mutual respect, in which we all feel comfortable to share our moments of confusion, ask questions, and challenge our understanding.

\sphinxAtStartPar
Everyone should be able to succeed in this course.
If you do not feel that is the case, please let us know.

\sphinxstepscope


\chapter{Modeling Neural Membranes}
\label{\detokenize{ModelingNeuralMembranes:modeling-neural-membranes}}\label{\detokenize{ModelingNeuralMembranes::doc}}
\sphinxAtStartPar
\sphinxstylestrong{In this lab, we’ll use an online membrane simulator and a breadboard to model the properties of neural membranes and passive potentials.}


\section{Neuromembrane}
\label{\detokenize{ModelingNeuralMembranes:neuromembrane}}
\sphinxAtStartPar
Neuromembrane was made by a team of University of Alberta researchers and a company known as Atmist. Many thanks to Christelle Sabatier, Michael Wright, and Declan Ali for sharing similar lesson plans. Their ideas are integrated here.


\bigskip\hrule\bigskip



\subsection{Part I: Resting membrane potential}
\label{\detokenize{ModelingNeuralMembranes:part-i-resting-membrane-potential}}
\sphinxAtStartPar
First, we’ll explore the properties of the membrane that lead to a resting membrane potential. By the end of this part, you will be able to:
\begin{itemize}
\item {} 
\sphinxAtStartPar
Explain how a negative resting membrane potential is maintained

\item {} 
\sphinxAtStartPar
Identify the factors that determine and maintain resting membrane potential

\item {} 
\sphinxAtStartPar
Use the Goldman\sphinxhyphen{}Hodgkin\sphinxhyphen{}Katz to calculate the membrane potential given the concentrations and permeabilities of K+ and Na+

\item {} 
\sphinxAtStartPar
Generate simulation data to illustrate the dependence of the resting membrane potential on K+

\end{itemize}


\subsection{Steps}
\label{\detokenize{ModelingNeuralMembranes:steps}}\begin{enumerate}
\sphinxsetlistlabels{\arabic}{enumi}{enumii}{}{.}%
\item {} 
\sphinxAtStartPar
Go to \sphinxurl{https://neuromembrane.biology.ualberta.ca/} and log in as a guest.

\item {} 
\sphinxAtStartPar
To open the resting membrane potential simulation, click on the top left to open a menu of stimulations. Choose “Resting Potential Simulation.”

\sphinxAtStartPar
On the right hand side, you’ll see a membrane. The area below the membrane is the inside of the cell. The area above the membrane is the outside of the cell.

\item {} 
\sphinxAtStartPar
This simulation will open by default without any leak channels or pumps. Let’s first add a couple of leak channels to the membrane and observe what happens. Click the plus sign next to “LEAK CHANNELS” to add a Na+ and K+ channel to the membrane.

\end{enumerate}

\sphinxAtStartPar
These single channels represent many, many channels in the membrane that are allowing either Na+ or K+ to cross. By default, these will open with a relative permeability of 5:100 (Na:K). Note also that the concentration settings are also set for you — at first, there is equal {[}Na+{]} and {[}K+{]} inside and outside of the membrane.
\begin{enumerate}
\sphinxsetlistlabels{\arabic}{enumi}{enumii}{}{.}%
\setcounter{enumi}{3}
\item {} 
\sphinxAtStartPar
First, we’ll observe what happens when we have more K+ permeability, but equal ion concentrations. Click on “CREATE SIMULATION” to open the stimulation page. \sphinxstylestrong{Before running the simulation, write a 1\sphinxhyphen{}2 sentence prediction for what will happen in Table 1.}

\item {} 
\sphinxAtStartPar
To run Simulation \#1, click the play button. You can also change the speed so that it will run faster. \sphinxstylestrong{What happens to the voltage over time? Why is the voltage across the membrane this value? Write your observations in Table 1.}

\item {} 
\sphinxAtStartPar
Go back to the Settings page by choosing “Back to settings” in the simulation play menu.

\item {} 
\sphinxAtStartPar
In order to modify the concentration across the membrane, we’ll need a Na+/K+ pump.
Click the plus button next to Na+/K+ pump to add one to the membrane.

\item {} 
\sphinxAtStartPar
Notice that the CONCENTRATION SETTINGS have now changed as well.  Change these concentrations back to the previous settings — where all of the concentrations are 50 mM. \sphinxstylestrong{Before running Simulation \#2, write your 1\sphinxhyphen{}2 sentence prediction in Table 1.}

\item {} 
\sphinxAtStartPar
Run Simulation \#2 and observe what happens to the membrane potential. \sphinxstylestrong{Has adding a Na+/K+ pump changed anything about the resting membrane potential? Why or why not? Add your observations to the table.} For this simulation, note that we’re just using the Na/K pump to be able to change the concentration across the membrane. The Na/K is not operational (yet).

\item {} 
\sphinxAtStartPar
Go back to the settings page. In real neurons, ion concentrations are not equivalent. Change the concentrations of K+ and Na+ back to their default values to reflect more biological values across the membrane. \sphinxstylestrong{Write your prediction for Simulation \#3 in the table.}

\item {} 
\sphinxAtStartPar
Run the simulation again and observe what happens to the membrane potential. \sphinxstylestrong{What happens to the voltage over time? Why is the voltage across the membrane this value? Write your observations in Table 1.}

\end{enumerate}

\sphinxAtStartPar
\sphinxstylestrong{Table 1. Simulation predictions \& observations}


\begin{savenotes}\sphinxattablestart
\sphinxthistablewithglobalstyle
\centering
\begin{tabulary}{\linewidth}[t]{TTTTTT}
\sphinxtoprule
\sphinxstyletheadfamily 
\sphinxAtStartPar
Sim \#
&\sphinxstyletheadfamily 
\sphinxAtStartPar
Types of Channels
&\sphinxstyletheadfamily 
\sphinxAtStartPar
{[}Ion{]}
&\sphinxstyletheadfamily 
\sphinxAtStartPar
PNa/K
&\sphinxstyletheadfamily 
\sphinxAtStartPar
Prediction
&\sphinxstyletheadfamily 
\sphinxAtStartPar
Observation
\\
\sphinxmidrule
\sphinxtableatstartofbodyhook
\sphinxAtStartPar
1
&
\sphinxAtStartPar
Leak channels only
&
\sphinxAtStartPar
All 50 mM
&
\sphinxAtStartPar
5 Na+; 100 K+
&
\sphinxAtStartPar

&
\sphinxAtStartPar

\\
\sphinxhline
\sphinxAtStartPar
2
&
\sphinxAtStartPar
Leak channels \& Na+/K+ pump
&
\sphinxAtStartPar
All 50 mM
&
\sphinxAtStartPar
5 Na+; 100 K+
&
\sphinxAtStartPar

&
\sphinxAtStartPar

\\
\sphinxhline
\sphinxAtStartPar
3
&
\sphinxAtStartPar
Leak channels \& Na+/K+ pump
&
\sphinxAtStartPar
Na+ high outside; K+ high inside
&
\sphinxAtStartPar
5 Na+; 100 K+
&
\sphinxAtStartPar

&
\sphinxAtStartPar

\\
\sphinxhline
\sphinxAtStartPar
4
&
\sphinxAtStartPar
Leak channels \& Na+/K+ pump
&
\sphinxAtStartPar
Na+ high outside; K+ high inside
&
\sphinxAtStartPar
100 Na+; 5 K+
&
\sphinxAtStartPar

&
\sphinxAtStartPar

\\
\sphinxhline
\sphinxAtStartPar
5
&
\sphinxAtStartPar
Leak channels \& Na+/K+ pump
&
\sphinxAtStartPar
Na+ high outside; K+ high inside
&
\sphinxAtStartPar
100 Na+; 100 K+
&
\sphinxAtStartPar

&
\sphinxAtStartPar

\\
\sphinxbottomrule
\end{tabulary}
\sphinxtableafterendhook\par
\sphinxattableend\end{savenotes}
\begin{enumerate}
\sphinxsetlistlabels{\arabic}{enumi}{enumii}{}{.}%
\setcounter{enumi}{11}
\item {} 
\sphinxAtStartPar
Using the Goldman\sphinxhyphen{}Hodgkin\sphinxhyphen{}Katz (GHK) equation for Na+ and K+ only, calculate what the voltage across the membrane should be for simulation \#3 and show your work. Does this match the voltage in the simulation? \sphinxstylestrong{Answer Q1 on the quiz. Report your answer in mV.}

\end{enumerate}

\sphinxAtStartPar
\sphinxstylestrong{Here are a few hints to help you:}
\begin{itemize}
\item {} 
\sphinxAtStartPar
P is relative permeability; other constants can be found on the lecture slides.

\item {} 
\sphinxAtStartPar
To simplify, you may determine a constant for the first portion of the equation (RT/F) and use that for future calculations.

\item {} 
\sphinxAtStartPar
You can use “Insert > Equation” in Google Docs or Microsoft Word to show your work, or simply jot it on paper, take a picture, and paste it here.

\item {} 
\sphinxAtStartPar
Google search has a built\sphinxhyphen{}in scientific calculator! You can bring it up by typing any formula into the search bar. For example:

\sphinxAtStartPar
\sphinxincludegraphics{{c016895e880266a7736d72bba73ec0b0ed180146}.png}

\end{itemize}

\begin{sphinxadmonition}{note}{Q1. Calculation of GHK}

\sphinxAtStartPar
Using the GHK equation for Na⁺ and K⁺ only, calculate the voltage across the membrane for simulation \#3 and show your work. Does this match the voltage in the simulation? Report your answer in mV.
\end{sphinxadmonition}
\begin{enumerate}
\sphinxsetlistlabels{\arabic}{enumi}{enumii}{}{.}%
\setcounter{enumi}{12}
\item {} 
\sphinxAtStartPar
Finally, let’s check to see how much ion permeability matters. Go back to the settings page and reverse the Na+ and K+ permeabilities.

\item {} 
\sphinxAtStartPar
Write your prediction for Simulation \#4 in Table 1.

\item {} 
\sphinxAtStartPar
Run Simulation \#4 and observe what happens. Record your observations in Table 1 and respond to the question below. \sphinxstylestrong{Answer Q2 on the quiz. Report your answer in mV.}

\end{enumerate}

\begin{sphinxadmonition}{note}{Q2.}

\sphinxAtStartPar
With the high permeability to Na⁺ and low permeability to K⁺, what is the resting membrane potential? Report your value in mV out to one decimal point.
\end{sphinxadmonition}
\begin{enumerate}
\sphinxsetlistlabels{\arabic}{enumi}{enumii}{}{.}%
\setcounter{enumi}{15}
\item {} 
\sphinxAtStartPar
On the settings page, set both the Na+ and K+ permeabilities to 100. Write your prediction in Table 1.

\item {} 
\sphinxAtStartPar
Run the simulation to see what happens, record your observations, and respond to the question below. \sphinxstylestrong{Answer Q3 on the quiz.}

\end{enumerate}

\begin{sphinxadmonition}{note}{Q3.}

\sphinxAtStartPar
Which of these simulations best models a typical neural membrane?
\end{sphinxadmonition}
\begin{enumerate}
\sphinxsetlistlabels{\arabic}{enumi}{enumii}{}{.}%
\setcounter{enumi}{17}
\item {} 
\sphinxAtStartPar
Next, we’ll check to see whether or not the membrane potential depends on {[}K+{]}out. Reset the permeabilities and ion concentrations to the default values.

\item {} 
\sphinxAtStartPar
Systematically change {[}K+{]}out to fill in the table below. Hint: if you hover over the chart on the left, you’ll be able to see the precise voltage. \sphinxstylestrong{Answer Q4 on the quiz.}

\end{enumerate}

\sphinxAtStartPar
\sphinxstylestrong{Table 2. Dependence of Vrest on {[}K+{]}out}


\begin{savenotes}\sphinxattablestart
\sphinxthistablewithglobalstyle
\centering
\begin{tabulary}{\linewidth}[t]{TT}
\sphinxtoprule
\sphinxstyletheadfamily 
\sphinxAtStartPar
{[}K+{]}out (mM)
&\sphinxstyletheadfamily 
\sphinxAtStartPar
Vrest (mV)
\\
\sphinxmidrule
\sphinxtableatstartofbodyhook
\sphinxAtStartPar
4
&
\sphinxAtStartPar

\\
\sphinxhline
\sphinxAtStartPar
10
&
\sphinxAtStartPar

\\
\sphinxhline
\sphinxAtStartPar
20
&
\sphinxAtStartPar

\\
\sphinxhline
\sphinxAtStartPar
50
&
\sphinxAtStartPar

\\
\sphinxhline
\sphinxAtStartPar
200
&
\sphinxAtStartPar

\\
\sphinxbottomrule
\end{tabulary}
\sphinxtableafterendhook\par
\sphinxattableend\end{savenotes}
\begin{enumerate}
\sphinxsetlistlabels{\arabic}{enumi}{enumii}{}{.}%
\setcounter{enumi}{19}
\item {} 
\sphinxAtStartPar
Click the “Toggle Circuit Diagram” button on the top right hand corner to overlay circuit components. \sphinxstylestrong{Answer Q5 on the quiz}.  You can also fill out the table below for your own notes.

\end{enumerate}

\begin{sphinxadmonition}{note}{Q5.}

\sphinxAtStartPar
In a few words, describe what each of the following are modeling in the neuron. The first one has been done for you.


\begin{savenotes}\sphinxattablestart
\sphinxthistablewithglobalstyle
\centering
\begin{tabulary}{\linewidth}[t]{TT}
\sphinxtoprule
\sphinxstyletheadfamily 
\sphinxAtStartPar
Component
&\sphinxstyletheadfamily 
\sphinxAtStartPar
What it models
\\
\sphinxmidrule
\sphinxtableatstartofbodyhook
\sphinxAtStartPar
\sphinxstylestrong{I\_Na}
&
\sphinxAtStartPar
sodium current across the membrane
\\
\sphinxhline
\sphinxAtStartPar
\sphinxstylestrong{g\_Na}
&
\sphinxAtStartPar

\\
\sphinxhline
\sphinxAtStartPar
\sphinxstylestrong{E\_Na}
&
\sphinxAtStartPar

\\
\sphinxhline
\sphinxAtStartPar
\sphinxstylestrong{I\_K}
&
\sphinxAtStartPar

\\
\sphinxhline
\sphinxAtStartPar
\sphinxstylestrong{g\_K}
&
\sphinxAtStartPar

\\
\sphinxhline
\sphinxAtStartPar
\sphinxstylestrong{E\_K}
&
\sphinxAtStartPar

\\
\sphinxbottomrule
\end{tabulary}
\sphinxtableafterendhook\par
\sphinxattableend\end{savenotes}
\end{sphinxadmonition}


\bigskip\hrule\bigskip



\subsection{Part II: Passive membrane simulation}
\label{\detokenize{ModelingNeuralMembranes:part-ii-passive-membrane-simulation}}
\sphinxAtStartPar
In this part, you’ll set up a simulation that models one portion of an axon or dendrite. By the end of this part, you will be able to:
\begin{itemize}
\item {} 
\sphinxAtStartPar
Describe how voltage passively spreads through an axon or dendrite

\item {} 
\sphinxAtStartPar
Understand length and time constants

\item {} 
\sphinxAtStartPar
Determine how the passive spread of voltage is affected by diameter and membrane capacitance

\end{itemize}


\subsection{Steps}
\label{\detokenize{ModelingNeuralMembranes:id1}}\begin{enumerate}
\sphinxsetlistlabels{\arabic}{enumi}{enumii}{}{.}%
\item {} 
\sphinxAtStartPar
Go to  \sphinxurl{https://neuromembrane.biology.ualberta.ca/} (if you’re not already there) and open up the “Passive Membrane Simulation” (or Cable Theory Simulation)  in the menu in the top left corner.

\item {} 
\sphinxAtStartPar
Add an external recording electrode that is two length constants () away from your stimulating electrode by clicking the \sphinxstylestrong{+} button and changing the value after the second recording site. \sphinxstylestrong{Answer Q6 on the quiz.}

\end{enumerate}

\begin{sphinxadmonition}{note}{Q6.}

\sphinxAtStartPar
Before running the simulation or changing any parameters, describe the setup of this “experiment” in 1\sphinxhyphen{}2 sentences.
\end{sphinxadmonition}
\begin{enumerate}
\sphinxsetlistlabels{\arabic}{enumi}{enumii}{}{.}%
\setcounter{enumi}{2}
\item {} 
\sphinxAtStartPar
Click the “Toggle Circuit Diagram” button on the top right hand corner to overlay circuit components. \sphinxstylestrong{Answer Q7 on the quiz.}

\end{enumerate}

\begin{sphinxadmonition}{note}{Q7.}

\sphinxAtStartPar
Describe what this circuit diagram is showing us. Be sure to define each of the different components (r\_m, r\_i, c\_m).
\end{sphinxadmonition}
\begin{enumerate}
\sphinxsetlistlabels{\arabic}{enumi}{enumii}{}{.}%
\setcounter{enumi}{3}
\item {} 
\sphinxAtStartPar
Re\sphinxhyphen{}set the cable settings and click on “CREATE SIMULATION” to run our first iteration of this passive current injection. The current will be injected starting at 10 ms, and will be injected for 30 ms total. Observe the top plot to fill out row one in Table 3. For your “observations,” take a close look at the Voltage vs. Time plot. \sphinxstylestrong{Why is our current injection this shape when recorded across the membrane?}

\end{enumerate}

\sphinxAtStartPar
\sphinxstylestrong{Table 3. Current injection results}


\begin{savenotes}\sphinxattablestart
\sphinxthistablewithglobalstyle
\centering
\begin{tabulary}{\linewidth}[t]{TTTTTTTT}
\sphinxtoprule
\sphinxstyletheadfamily 
\sphinxAtStartPar
Sim \#
&\sphinxstyletheadfamily 
\sphinxAtStartPar
Diameter
&\sphinxstyletheadfamily 
\sphinxAtStartPar
Cm
&\sphinxstyletheadfamily 
\sphinxAtStartPar
Peak voltage at 2
&\sphinxstyletheadfamily 
\sphinxAtStartPar
Distance to 2 electrode
&\sphinxstyletheadfamily 
\sphinxAtStartPar
Voltage at 0 electrode @ 28 ms
&\sphinxstyletheadfamily 
\sphinxAtStartPar
Voltage at 2 electrode @ 28 ms
&\sphinxstyletheadfamily 
\sphinxAtStartPar
Observations
\\
\sphinxmidrule
\sphinxtableatstartofbodyhook
\sphinxAtStartPar
1
&
\sphinxAtStartPar
5 µm
&
\sphinxAtStartPar
1
&
\sphinxAtStartPar

&
\sphinxAtStartPar

&
\sphinxAtStartPar

&
\sphinxAtStartPar

&
\sphinxAtStartPar

\\
\sphinxhline
\sphinxAtStartPar
2
&
\sphinxAtStartPar
2 µm
&
\sphinxAtStartPar
1
&
\sphinxAtStartPar

&
\sphinxAtStartPar

&
\sphinxAtStartPar

&
\sphinxAtStartPar

&
\sphinxAtStartPar

\\
\sphinxhline
\sphinxAtStartPar
3
&
\sphinxAtStartPar
2 µm
&
\sphinxAtStartPar
2
&
\sphinxAtStartPar

&
\sphinxAtStartPar

&
\sphinxAtStartPar

&
\sphinxAtStartPar

&
\sphinxAtStartPar

\\
\sphinxbottomrule
\end{tabulary}
\sphinxtableafterendhook\par
\sphinxattableend\end{savenotes}
\begin{enumerate}
\sphinxsetlistlabels{\arabic}{enumi}{enumii}{}{.}%
\setcounter{enumi}{4}
\item {} 
\sphinxAtStartPar
The dorsal root ganglion axon of a giant blue whale has a diameter of \textasciitilde{}2 µm. In the CABLE SETTINGS on the left, change the diameter of the simulated axon to model this axon.

\end{enumerate}

\begin{sphinxadmonition}{note}{Reflection}

\sphinxAtStartPar
With this smaller diameter of 2 µm, do other parameters of the membrane (bottom of the CABLE SETTINGS window) change? If so, why?
\end{sphinxadmonition}
\begin{enumerate}
\sphinxsetlistlabels{\arabic}{enumi}{enumii}{}{.}%
\setcounter{enumi}{5}
\item {} 
\sphinxAtStartPar
Run the simulation with the smaller diameter and observe the Voltage vs Time plot again. Hint: you can use the Autoscale button (\sphinxincludegraphics{{7284b663657c51e31baaf9522713c4c9af0bfa83}.png}) if the trace goes off the plot. Fill out Table 3 for this 2 µm cable simulation.

\item {} 
\sphinxAtStartPar
Double the capacitance of the membrane, from 1 to 2 μF/cm2. Describe what happens to the time constant and the resulting voltage over time plot with more capacitance in Table 3.

\end{enumerate}

\sphinxAtStartPar
Optional: Later in the course, you may want to use the Neuromembrane simulator to model \sphinxhref{https://docs.google.com/document/d/14yS7gj4IyZvIWO3KI5tqYkXIkytr7Y8ES7v4SbK9quM/edit?tab=t.0\#heading=h.z78xt7uiejq}{action potential}.


\bigskip\hrule\bigskip



\section{RC Circuits}
\label{\detokenize{ModelingNeuralMembranes:rc-circuits}}
\sphinxAtStartPar
Let’s build a circuit that models the passive membrane simulations you performed above.


\subsection{Building an RC Circuit}
\label{\detokenize{ModelingNeuralMembranes:building-an-rc-circuit}}
\sphinxAtStartPar
Read pages 4\sphinxhyphen{}7 of the RC circuit lab exercise in your lab manual. Then, build the circuit below. We’ll use the OUTPUT of your Powerlab to send a pulse of current through the circuit.
\begin{enumerate}
\sphinxsetlistlabels{\arabic}{enumi}{enumii}{}{.}%
\item {} 
\sphinxAtStartPar
It helps to start with the output ends in the power rails and work your way around. The \sphinxstylestrong{red} wire from the PowerLab is the positive terminal.

\sphinxAtStartPar
\sphinxstylestrong{CAUTION}: Do not connect the + and \sphinxhyphen{} terminals. This will short the circuit, heat, and melt the plastic. Check that your circuit is set up properly before turning on the pulse from your computer.

\item {} 
\sphinxAtStartPar
Attach single BNC adapters to the + and \sphinxhyphen{} outputs of the PowerLab.
\begin{enumerate}
\sphinxsetlistlabels{\arabic}{enumii}{enumiii}{}{.}%
\item {} 
\sphinxAtStartPar
Alternatively, connect a \sphinxstylestrong{BNC to double banana cable adapter} to the + output of your PowerLab.

\end{enumerate}

\item {} 
\sphinxAtStartPar
Attach a red banana cable to the + terminal and a black banana cable to the \sphinxhyphen{} terminal of your PowerLab.

\item {} 
\sphinxAtStartPar
Attach the banana cables to red and black jumper wires using alligator clips, and insert the black and red jumper wires to the power rails, as shown in the circuit.

\sphinxAtStartPar
\sphinxstylestrong{NOTE:} Ensure that all circuit components are properly inserted into the breadboard. The wires should make contact with the metal strips inside the breadboard.

\item {} 
\sphinxAtStartPar
Add additional jumper wires, a resistor, and a capacitor, as shown in the circuit.

\item {} 
\sphinxAtStartPar
We’ll use the PowerLab as our voltmeter, to record voltage in our circuit. Plug a \sphinxstylestrong{BNC to DIN8 adapter} and \sphinxstylestrong{double banana adapter} into Input 1 of the PowerLab.

\item {} 
\sphinxAtStartPar
Connect the ground/reference (\sphinxstylestrong{black}, \sphinxhyphen{}) side of the double banana adapter to the circuit ground with alligator clips.

\sphinxAtStartPar
\sphinxstylestrong{NOTE:} The “ground” in this circuit is really going to be our reference electrode. It should be connected to the black side of Input 1.

\item {} 
\sphinxAtStartPar
Connect the “hot” pin (\sphinxstylestrong{red,} +) of the double banana adaptor to the “record here” point on the circuit diagram.

\item {} 
\sphinxAtStartPar
Open LabChart with the \sphinxstylestrong{Circuit Lab} settings (\sphinxurl{http://bit.ly/labchart})

\item {} 
\sphinxAtStartPar
Set up the Stimulator to stimulate with 1 V for \sphinxstylestrong{1 second}. Make sure the range of your recording channel is 1 V.  Switch the stimulator to “on”.

\item {} 
\sphinxAtStartPar
Open Scope View and press \sphinxstylestrong{>Start}. You should see a waveform on your screen.\\
\sphinxstylestrong{NOTE:} You may need to Autoscale the channel (right click) to be able to see the entire trace.

\end{enumerate}

\sphinxAtStartPar
\sphinxstylestrong{The time constant should equal 𝛕 = RC. What \sphinxstyleemphasis{should} the time constant be (in seconds), given the resistors and capacitors in the circuit?} Mind your units. \sphinxstylestrong{Answer Q8 on the quiz.}

\sphinxAtStartPar
𝛕 = \_\_\_\_\_\_\_\_\_\_\_\_\_\_\_

\sphinxAtStartPar
We can also measure the time constant by measuring how long it takes the curve to rise or decay. One time constant (𝛕) is defined by when the circuit rises to \sphinxstylestrong{63.2\%} of its total charge (alternately, we could see how long it takes to decay 1/e, about 37\%).

\sphinxAtStartPar
To empirically measure your time constant, follow these steps.
\begin{enumerate}
\sphinxsetlistlabels{\arabic}{enumi}{enumii}{}{.}%
\item {} 
\sphinxAtStartPar
Make sure your voltage is very close to zero at baseline.

\item {} 
\sphinxAtStartPar
Find the maximum voltage (Vmax) that your circuit reaches.

\item {} 
\sphinxAtStartPar
Calculate 0.632 * Vmax

\item {} 
\sphinxAtStartPar
Using the marker tool (right click anywhere in the window and choose “Clamp to trace”) find the point in your trace that is equal to the value calculated in Step 3.

\item {} 
\sphinxAtStartPar
Move your cursor to the start of the stimulus. In the upper right corner of the window, you’ll see something that says Δt. This is the change in time since the mark on the curve and your time constant.

\end{enumerate}

\sphinxAtStartPar
\sphinxstylestrong{What is your recorded time constant with 1 V stimulation, a 100 kΩ resistor, and a 1 uF capacitor?}

\sphinxAtStartPar
𝛕 = \_\_\_\_\_\_\_\_\_\_\_\_\_\_\_

\sphinxAtStartPar
\sphinxstylestrong{How close is the recorded time constant to the value you calculated with 𝛕 = RC?}

\sphinxAtStartPar
In engineering, 5𝛕 is consisted the time it takes for a capacitor to fully charge (technically, its at 99\% of full charge). The time constant is symmetrical for charging as well as losing charge. \sphinxstylestrong{How long would it take this patch of membrane to return to baseline, then? What does that mean for \sphinxstyleemphasis{computation} happening in this membrane?}


\subsection{Record and plot your time constant with different resistors}
\label{\detokenize{ModelingNeuralMembranes:record-and-plot-your-time-constant-with-different-resistors}}
\sphinxAtStartPar
Next, we’ll see how the time constant of the circuit changes in response to the resistance strength. Use the following resistors in your circuit, and use the steps above to determine the time constant for each.


\begin{savenotes}\sphinxattablestart
\sphinxthistablewithglobalstyle
\centering
\begin{tabulary}{\linewidth}[t]{TT}
\sphinxtoprule
\sphinxstyletheadfamily 
\sphinxAtStartPar
Resistance
&\sphinxstyletheadfamily 
\sphinxAtStartPar
Time Constant
\\
\sphinxmidrule
\sphinxtableatstartofbodyhook
\sphinxAtStartPar
3 kΩ
&
\sphinxAtStartPar

\\
\sphinxhline
\sphinxAtStartPar
15 kΩ
&
\sphinxAtStartPar

\\
\sphinxhline
\sphinxAtStartPar
33 kΩ
&
\sphinxAtStartPar

\\
\sphinxhline
\sphinxAtStartPar
100  kΩ
&
\sphinxAtStartPar

\\
\sphinxhline
\sphinxAtStartPar
220 kΩ
&
\sphinxAtStartPar

\\
\sphinxhline
\sphinxAtStartPar
330 kΩ
&
\sphinxAtStartPar

\\
\sphinxbottomrule
\end{tabulary}
\sphinxtableafterendhook\par
\sphinxattableend\end{savenotes}

\sphinxAtStartPar
\sphinxstylestrong{Plot the time constant against the resistance.  Answer Q9a and Q9b on the quiz.}

\sphinxAtStartPar
The plot should have:
\begin{itemize}
\item {} 
\sphinxAtStartPar
A line fit to it, with a description of what the trendline is (you can either show the equation or say in the figure caption what kind of trendline is plotted with the R2 value.

\item {} 
\sphinxAtStartPar
Labeled axes \& units

\item {} 
\sphinxAtStartPar
At least six data points

\end{itemize}

\sphinxAtStartPar
Plotting Tips:
\begin{itemize}
\item {} 
\sphinxAtStartPar
Python: Use the \sphinxhref{https://bipn145.github.io/Python/PlotScatter.html}{Plot Scatter} notebook

\item {} 
\sphinxAtStartPar
Excel:
\begin{itemize}
\item {} 
\sphinxAtStartPar
Insert your plot as a scatter plot (this will tell Excel to treat your x\sphinxhyphen{}axis values (resistance) as actual values, not categories.

\item {} 
\sphinxAtStartPar
Add a trendline by right clicking on the data, and choosing “Add Trendline”. You should add a trendline for the trend you’re expecting (we talked about this in lecture).

\item {} 
\sphinxAtStartPar
Make sure that none of the markers for your points are cut off. If they are, change your axes so that they fit.

\end{itemize}

\end{itemize}


\subsection{Check your time constant against the literature}
\label{\detokenize{ModelingNeuralMembranes:check-your-time-constant-against-the-literature}}
\sphinxAtStartPar
Let’s circle back to biology. Are the time constants you measured similar to what you’d find in a real neuron? Check this website to look at a collection of examples from the literature: \sphinxurl{https://neuroelectro.org/ephys\_prop/index.html}

\sphinxAtStartPar
Choose a paper from the example above and look for the reported time constant. Is the time constant in your circuit similar? If not, why might it be different?


\bigskip\hrule\bigskip



\subsection{Troubleshooting}
\label{\detokenize{ModelingNeuralMembranes:troubleshooting}}

\begin{savenotes}\sphinxattablestart
\sphinxthistablewithglobalstyle
\centering
\begin{tabulary}{\linewidth}[t]{TTT}
\sphinxtoprule
\sphinxstyletheadfamily 
\sphinxAtStartPar
Observation
&\sphinxstyletheadfamily 
\sphinxAtStartPar
Likely Issue
&\sphinxstyletheadfamily 
\sphinxAtStartPar
Possible Solution
\\
\sphinxmidrule
\sphinxtableatstartofbodyhook
\sphinxAtStartPar
There isn’t any voltage in your circuit
&
\sphinxAtStartPar
The power terminals are not connected properly Your stimulator is not on
&
\sphinxAtStartPar
Make sure everything is connected correctly on your breadboard. Consult \sphinxurl{https://wiring.org.co/learning/tutorials/breadboard/index.html}  Check that the stimulator panel is switched on, and that you pressed Start
\\
\sphinxhline
\sphinxAtStartPar
Your signal is really noisy
&
\sphinxAtStartPar
You’re not grounded. Your pulse is not being properly sent.
&
\sphinxAtStartPar
Double check your circuit is connected correctly
\\
\sphinxhline
\sphinxAtStartPar
You don’t see a fin\sphinxhyphen{}shaped waveform
&
\sphinxAtStartPar
Your Scope/Chart View is not scaled properly.  Your capacitor is not properly connected or broken.
&
\sphinxAtStartPar
Try autoscaling Check the capacitor placement or replace it.
\\
\sphinxhline
\sphinxAtStartPar
You’ve got “rabbit ears” on your waveform
&
\sphinxAtStartPar
Your voltage is reversed somewhere
&
\sphinxAtStartPar
Check the power terminals are connected in the right direction
\\
\sphinxbottomrule
\end{tabulary}
\sphinxtableafterendhook\par
\sphinxattableend\end{savenotes}

\sphinxstepscope


\chapter{String Nervous Systems}
\label{\detokenize{StringNervousSystems:string-nervous-systems}}\label{\detokenize{StringNervousSystems::doc}}

\section{Background \& Goals}
\label{\detokenize{StringNervousSystems:background-goals}}

\subsection{Why are we recording from a string?}
\label{\detokenize{StringNervousSystems:why-are-we-recording-from-a-string}}
\sphinxAtStartPar
The purpose of this lab is to get familiar with the recording equipment and software that we’ll be using in this course.

\sphinxAtStartPar
First, we’ll learn how you can configure LabChart 8 to acquire data and send out pulses of electricity. Then, you’ll record the \sphinxstylestrong{stimulus artifact} (the passively conducted signal from the stimulus) from a saline\sphinxhyphen{}soaked piece of string. The string mimics some of the properties of the earthworm that you’ll observe in a future experiment. This will give you the opportunity to learn how to use the stimulating and recording devices on the PowerLab, as well as reduce noise and stimulus artifacts from your recordings, before you actually have to do all of this in an experiment.

\sphinxAtStartPar
Asking questions during this lab is a great idea, and will help you succeed in future experiments.


\subsection{Pre\sphinxhyphen{}Lab}
\label{\detokenize{StringNervousSystems:pre-lab}}
\sphinxAtStartPar
Before the lab, watch this video and read the protocol below:




\bigskip\hrule\bigskip


\sphinxAtStartPar
\sphinxstylestrong{Today you will:}
\begin{itemize}
\item {} 
\sphinxAtStartPar
Configure a recording experiment in LabChart 8.

\item {} 
\sphinxAtStartPar
Determine how to record signals from a string.

\item {} 
\sphinxAtStartPar
Identify sources of electrical noise in your rig.

\end{itemize}

\sphinxAtStartPar
\sphinxstyleemphasis{This lab was adapted from Experiment \#1 from the BIPN 105 \& BIPN 145 Lab Manuals, as well as AD Instruments.}


\bigskip\hrule\bigskip



\section{Lab Protocol}
\label{\detokenize{StringNervousSystems:lab-protocol}}

\begin{savenotes}\sphinxattablestart
\sphinxthistablewithglobalstyle
\centering
\begin{tabulary}{\linewidth}[t]{TTT}
\sphinxtoprule
\sphinxstyletheadfamily 
\sphinxAtStartPar
\sphinxstylestrong{Supplies}
&\sphinxstyletheadfamily 
\sphinxAtStartPar
\sphinxstylestrong{Cables \& connectors}
&\sphinxstyletheadfamily 
\sphinxAtStartPar
\sphinxstylestrong{Solutions}
\\
\sphinxmidrule
\sphinxtableatstartofbodyhook
\sphinxAtStartPar
String
&
\sphinxAtStartPar
BNC to single banana adapter (2)
&
\sphinxAtStartPar
Saline
\\
\sphinxhline
\sphinxAtStartPar
Dissection tray
&
\sphinxAtStartPar
BNC to double banana adapter
&
\sphinxAtStartPar

\\
\sphinxhline
\sphinxAtStartPar
PowerLab 26T, connected to computer
&
\sphinxAtStartPar
BNC to DIN8 adapter
&
\sphinxAtStartPar

\\
\sphinxhline
\sphinxAtStartPar
BioAmp
&
\sphinxAtStartPar
Alligator clip adapters (2)
&
\sphinxAtStartPar

\\
\sphinxhline
\sphinxAtStartPar
Sewing pins (2)
&
\sphinxAtStartPar
Banana plug cords (2)
&
\sphinxAtStartPar

\\
\sphinxhline
\sphinxAtStartPar
Faraday Cage
&
\sphinxAtStartPar

&
\sphinxAtStartPar

\\
\sphinxhline
\sphinxAtStartPar
Large Petri Dish
&
\sphinxAtStartPar

&
\sphinxAtStartPar

\\
\sphinxhline
\sphinxAtStartPar
Eyedropper
&
\sphinxAtStartPar

&
\sphinxAtStartPar

\\
\sphinxbottomrule
\end{tabulary}
\sphinxtableafterendhook\par
\sphinxattableend\end{savenotes}


\section{I. Setting up LabChart 8}
\label{\detokenize{StringNervousSystems:i-setting-up-labchart-8}}
\sphinxAtStartPar
LabChart 8 is the software we’ll be using to record electrophysiology data. Before performing our experiment, we need to set up LabChart to \sphinxstylestrong{acquire} and \sphinxstylestrong{display} the data as we want.
\begin{enumerate}
\sphinxsetlistlabels{\arabic}{enumi}{enumii}{}{.}%
\item {} 
\sphinxAtStartPar
Turn on the PowerLab 26T.

\item {} 
\sphinxAtStartPar
Open LabChart 8. You should see that the PowerLab is connected with a green check mark.

\item {} 
\sphinxAtStartPar
Open a new experiment by choosing the \sphinxstylestrong{New} button (bottom left\sphinxhyphen{}hand corner).

\item {} 
\sphinxAtStartPar
First, we need to decide how many channels to record. Go to \sphinxstylestrong{Setup > Channel Settings}.
\begin{itemize}
\item {} 
\sphinxAtStartPar
At the bottom of the window, change the number of channels to \sphinxstylestrong{3}, and make sure that they are all on (box on the left is checked).

\item {} 
\sphinxAtStartPar
Rename the first channel \sphinxstylestrong{“Stimulus”} and the second channel \sphinxstylestrong{“Raw Recording”}.

\item {} 
\sphinxAtStartPar
Set the range for Channel 1 (your “Stimulus” channel) to \sphinxstylestrong{5 V}.

\end{itemize}
\begin{quote}

\sphinxAtStartPar
\sphinxstylestrong{Note:} The range determines the range of values the PowerLab will record, and when using the BioAmp (later in this lab), it will actually change the amplification.
\end{quote}
\begin{itemize}
\item {} 
\sphinxAtStartPar
Set the range for Channel 2 (your “Raw Recording” channel) to \sphinxstylestrong{100 mV}.

\item {} 
\sphinxAtStartPar
Make the sampling rate for Channel 1 \sphinxstylestrong{“40k/s”} and ensure that \sphinxstyleemphasis{“Same sampling rate for all channels”} is checked.

\item {} 
\sphinxAtStartPar
Make sure that the far\sphinxhyphen{}right column says \sphinxstylestrong{“No calculation”} for all three channels.

\end{itemize}

\end{enumerate}


\section{II. Finding sources of electrical noise}
\label{\detokenize{StringNervousSystems:ii-finding-sources-of-electrical-noise}}
\sphinxAtStartPar
One of the biggest issues in electrophysiology is noise from overhead lights and other electronics in the room. It’ll help to get an understanding of where noise in your rig is coming from.
\begin{enumerate}
\sphinxsetlistlabels{\arabic}{enumi}{enumii}{}{.}%
\item {} 
\sphinxAtStartPar
Connect a DIN8\sphinxhyphen{}BNC adapter and a double 3\sphinxhyphen{}way BNC adapter to \sphinxstylestrong{Input 2} of the PowerLab. Plug in a red banana cable to the positive (red) side of the connector.

\item {} 
\sphinxAtStartPar
Hold the lead (the end of the banana cable) in your hand.

\item {} 
\sphinxAtStartPar
In the Channel Settings window, click on \sphinxstylestrong{Input Amplifier} for Channel 2. The Input Amplifier window will show the voltage that you’re recording. This allows for precise setting of the input range for a recording channel and provides filtering options.

\item {} 
\sphinxAtStartPar
The signal at your rig may be variable. To adjust the sensitivity of the channel, choose an appropriate range setting from the \sphinxstylestrong{Range} drop\sphinxhyphen{}down list in the Input Amplifier dialog.
\begin{quote}

\sphinxAtStartPar
\sphinxstylestrong{Note:} The number displayed in the range menu indicates the maximum input voltage currently selected. Notice how as you decrease the range the vertical scale changes and the small rhythmic deflections that appear on the signal trace increase in amplitude.
\end{quote}

\item {} 
\sphinxAtStartPar
Choose a range such that your signal fills about ⅔ of the window.

\item {} 
\sphinxAtStartPar
Observe the change in the recorded voltage with the electrode lead held overhead (near the lights) and then near the ground (away from the lights).

\item {} 
\sphinxAtStartPar
Try the lead inside of your Faraday cage. Is the signal more or less noisy?
\begin{quote}

\sphinxAtStartPar
\sphinxstylestrong{Tip:} Is your Faraday cage actually grounded? How could you ground it?
\end{quote}

\end{enumerate}

\begin{sphinxadmonition}{note}{Questions for reflection}
\begin{itemize}
\item {} 
\sphinxAtStartPar
How does the signal change as you move the electrode lead around?

\item {} 
\sphinxAtStartPar
What do these sources of noise mean for our recording configuration, and our attempts to record electrical activity?

\end{itemize}
\end{sphinxadmonition}


\section{III. Recording stimulator outputs}
\label{\detokenize{StringNervousSystems:iii-recording-stimulator-outputs}}
\sphinxAtStartPar
We’ll use our PowerLab as a stimulator that can send voltage pulses into our specimen. We can change the timing, shape, amplitude, and repetition of these stimuli via the Stimulator window in LabChart. We’ll directly record these stimuli by sending the output of the PowerLab into an input.
\begin{enumerate}
\sphinxsetlistlabels{\arabic}{enumi}{enumii}{}{.}%
\item {} 
\sphinxAtStartPar
Change your view to \sphinxstylestrong{Scope view}, which looks like the sweeps of an oscilloscope. The advantage of Scope View is that it allows the overlay and averaging of your data.

\item {} 
\sphinxAtStartPar
In Scope view, set the \sphinxstylestrong{Duration} to \sphinxstylestrong{20 ms}. This changes how long each sweep of data is.

\item {} 
\sphinxAtStartPar
Remove the DIN8\sphinxhyphen{}BNC adapter \& cable from Input 2.

\item {} 
\sphinxAtStartPar
On the PowerLab, connect the \sphinxstylestrong{(+) Output} to \sphinxstylestrong{Input 1} using the BNC to DIN8 cable.

\item {} 
\sphinxAtStartPar
Go to \sphinxstylestrong{Setup > Stimulator} to open the Stimulator window. This is the window where you can modify the shape, timing, and amplitude of the stimulus output.

\item {} 
\sphinxAtStartPar
Make sure that the \sphinxstylestrong{Waveform Name} is set to \sphinxstylestrong{Pulse}.

\item {} 
\sphinxAtStartPar
Change the range of the Pulse Height to \sphinxstylestrong{\sphinxhyphen{}5 to 5 V} by clicking the ruler button next to the Pulse Height option.

\item {} 
\sphinxAtStartPar
Configure the following by clicking and then editing the values in the window:


\begin{savenotes}\sphinxattablestart
\sphinxthistablewithglobalstyle
\centering
\begin{tabulary}{\linewidth}[t]{TT}
\sphinxtoprule
\sphinxstyletheadfamily 
\sphinxAtStartPar
Parameter
&\sphinxstyletheadfamily 
\sphinxAtStartPar
Value
\\
\sphinxmidrule
\sphinxtableatstartofbodyhook
\sphinxAtStartPar
Start Delay
&
\sphinxAtStartPar
1 ms
\\
\sphinxhline
\sphinxAtStartPar
Pulse Width
&
\sphinxAtStartPar
0.1 ms
\\
\sphinxhline
\sphinxAtStartPar
Pulse Height
&
\sphinxAtStartPar
0.150 V
\\
\sphinxbottomrule
\end{tabulary}
\sphinxtableafterendhook\par
\sphinxattableend\end{savenotes}

\item {} 
\sphinxAtStartPar
Close the Stimulator window.

\item {} 
\sphinxAtStartPar
Go to \sphinxstylestrong{Setup > Stimulator Panel} to open a smaller window where you can directly modify the stimulus. This window can stay open while you’re recording.

\item {} 
\sphinxAtStartPar
Press \sphinxstylestrong{> Start} in the main LabChart GUI.

\item {} 
\sphinxAtStartPar
Increase the stimulator voltage (Pulse Height) and watch the level change in the window.

\item {} 
\sphinxAtStartPar
Try a negative stimulus voltage.

\item {} 
\sphinxAtStartPar
Open the Stimulator window again and vary the delay, duration, and waveform shape, and see what happens. If your recording is cut off, how can you fix it? \sphinxstyleemphasis{(Hint: see Visualization tips in Meet Your PowerLab.)}

\end{enumerate}

\begin{sphinxadmonition}{note}{Questions for reflection}
\begin{itemize}
\item {} 
\sphinxAtStartPar
Why would you want to record the stimulus output?

\item {} 
\sphinxAtStartPar
Why would you want to delay the stimulus output?

\item {} 
\sphinxAtStartPar
Why would you want to send different shapes of waveforms?

\end{itemize}
\end{sphinxadmonition}


\bigskip\hrule\bigskip



\section{IV. Recording the stimulus artifact}
\label{\detokenize{StringNervousSystems:iv-recording-the-stimulus-artifact}}
\sphinxAtStartPar
Lastly, we’ll add an amplifier (the BioAmp) to our circuit, and use this to record from a string soaked in saline. The string doesn’t have a nervous system, but we’ll add some voltage with our stimulator and observe the change in voltage from another point on the string.


\subsection{Set up your string stimulation experiment}
\label{\detokenize{StringNervousSystems:set-up-your-string-stimulation-experiment}}
\sphinxAtStartPar
Your setup for the string stimulation will look like this:

\sphinxAtStartPar
\sphinxincludegraphics{{powerlab_front}.png}
\begin{enumerate}
\sphinxsetlistlabels{\arabic}{enumi}{enumii}{}{.}%
\item {} 
\sphinxAtStartPar
Remove all of the cords from the PowerLab.

\item {} 
\sphinxAtStartPar
Connect your two single BNC to banana cable adapters into the \sphinxstylestrong{(+)} and \sphinxstylestrong{(\sphinxhyphen{})} Outputs of the PowerLab.
\begin{quote}

\sphinxAtStartPar
\sphinxstylestrong{Note:} The color of these does not matter, but it can help to be consistent — red for the anode (+ Output) and black for the cathode (\sphinxhyphen{} Output).
\end{quote}

\item {} 
\sphinxAtStartPar
Using banana cables, connect the \sphinxstylestrong{(+) Output} to your \sphinxstylestrong{anode} stimulating electrode (the sewing pin) and the \sphinxstylestrong{(\sphinxhyphen{})  Output} to the \sphinxstylestrong{cathode} stimulating electrode.

\item {} 
\sphinxAtStartPar
Connect the BioAmp to the PowerLab.

\item {} 
\sphinxAtStartPar
Connect the (+) recording, (\sphinxhyphen{}) recording, and ground needle electrodes to your string and into the corresponding spots on the BioAmp.
\begin{quote}

\sphinxAtStartPar
\sphinxstylestrong{Note:} The colors on the BioAmp defy convention — consult the diagram above for clarity.
\end{quote}

\item {} 
\sphinxAtStartPar
Change the configuration of your recording so that you are now using the BioAmp:
\begin{itemize}
\item {} 
\sphinxAtStartPar
Go into \sphinxstylestrong{Setup > Channel Settings}.

\item {} 
\sphinxAtStartPar
Make sure that \sphinxstylestrong{BioAmp} is listed under the Input Amplifier column for Channel 3.

\item {} 
\sphinxAtStartPar
\sphinxstylestrong{Uncheck} Channel 2 — we don’t need to record from it anymore.

\item {} 
\sphinxAtStartPar
Change the name of Channel 3 to \sphinxstylestrong{“BioAmp Recording.”}

\end{itemize}

\item {} 
\sphinxAtStartPar
We’ve removed our ability to directly record the stimulus output, so we need another way to know when the stimulator sends outputs. Open the Stimulator Settings and set the \sphinxstylestrong{“Marker Channel”} to Ch1 (your “Stimulus” channel). With this setting, LabChart will put a small tick on Channel 1 to show you when the stimulus was sent.
\begin{quote}

\sphinxAtStartPar
\sphinxstylestrong{Note:} This marker tells you \sphinxstyleemphasis{when} the stimulus was sent, but does not indicate the height or length of the stimulus. This means you need to label your pages with your stimulation parameters.
\end{quote}

\end{enumerate}


\subsection{Stimulate your string!}
\label{\detokenize{StringNervousSystems:stimulate-your-string}}
\noindent{\hspace*{\fill}\sphinxincludegraphics[width=150\sphinxpxdimen]{{string_electrodes}.png}\hspace*{\fill}}
\begin{enumerate}
\sphinxsetlistlabels{\arabic}{enumi}{enumii}{}{.}%
\item {} 
\sphinxAtStartPar
Place your string in saline in the small petri dish, and give it a few minutes to soak. If it’s not completely soaked, you will not get a stimulus artifact.
\begin{quote}

\sphinxAtStartPar
\sphinxstylestrong{Note:} If your string dries out during the experiment, you can use an eyedropper to add saline to it.
\end{quote}

\item {} 
\sphinxAtStartPar
Lay the string out on the dissection tray, with the stimulating and recording electrodes laid out as in the diagram above.

\item {} 
\sphinxAtStartPar
Set the stimulator voltage to \sphinxstylestrong{1 V}, and press \sphinxstylestrong{> Start}. You should see your stimulus recording on Channel 1, as you did before, but you should also now have a recording on Channel 3 (BioAmp Recording).

\item {} 
\sphinxAtStartPar
Change the scaling if you cannot see the entire deflection.

\item {} 
\sphinxAtStartPar
If you do not see anything, increase the voltage (no more than 5 V).
\begin{quote}

\sphinxAtStartPar
\sphinxstylestrong{Note:} This recorded deflection is called the \sphinxstylestrong{stimulus artifact}. This artifact can exceed the range of our amplifier. Recovery from the artifact should be rapid, however, so that the baseline is reached within one millisecond or so after the stimulus pulse.
\end{quote}

\end{enumerate}

\begin{sphinxadmonition}{note}{Question for reflection}

\sphinxAtStartPar
Is the amplitude of the stimulus artifact the same as the amplitude of the stimulus? Why or why not?
\end{sphinxadmonition}
\begin{enumerate}
\sphinxsetlistlabels{\arabic}{enumi}{enumii}{}{.}%
\setcounter{enumi}{5}
\item {} 
\sphinxAtStartPar
Increase the stimulator voltage, and stimulate again. What happens?
\begin{quote}

\sphinxAtStartPar
\sphinxstylestrong{Note:} Change the name of Scope pages, or make comments on the recording, to indicate what you’ve done during that run of the stimulus. You need to know the metadata for the experiment once you’re ready to export.
\end{quote}

\item {} 
\sphinxAtStartPar
Change the stimulus to a \sphinxstylestrong{negative voltage}, and stimulate again. Does the stimulus artifact change?

\item {} 
\sphinxAtStartPar
Choose one of your pages where you have a stimulus artifact with a clear peak. Using the \sphinxstylestrong{Marker tool}, measure the latency from the stimulus to the beginning of your artifact. How much of a delay is there? How much of a delay \sphinxstyleemphasis{should} there be? (How fast does an electrical current flow?)

\item {} 
\sphinxAtStartPar
Change the location of the stimulating electrodes. How does this change the shape of the artifact?

\item {} 
\sphinxAtStartPar
Remove the ground electrode. How does this change your recording?

\end{enumerate}


\subsection{Save your file \& export your data}
\label{\detokenize{StringNervousSystems:save-your-file-export-your-data}}
\sphinxAtStartPar
LabChart buffers files to disk in case of a power failure or computer crash, but it’s wise to save work frequently. Saving files in LabChart is the same as saving any file on your personal computer.

\sphinxAtStartPar
Please backup your LabChart data files to the Google Drive you set up at the start of the quarter.
\begin{enumerate}
\sphinxsetlistlabels{\arabic}{enumi}{enumii}{}{.}%
\item {} 
\sphinxAtStartPar
Save the \sphinxstylestrong{settings file} for your experiment by going to \sphinxstylestrong{File > Save As Settings File} in LabChart. Save it in a safe place (not locally on the computer). You will need this settings file for the earthworm experiment.

\item {} 
\sphinxAtStartPar
Save the \sphinxstylestrong{data file} for your experiment by going to \sphinxstylestrong{File > Save As} and save as a LabChart Data File (\sphinxstyleemphasis{.adicht}).

\item {} 
\sphinxAtStartPar
Go to the Scope page for \sphinxstylestrong{two of your favorite string stimulation experiments} that you can compare and contrast — for example, two different voltages, or with two different distances between the stimulator and recording electrodes.

\item {} 
\sphinxAtStartPar
Follow \sphinxhref{https://bipn145.github.io/LabChart/ExportingLabChart.html}{these instructions for exporting your data}.

\item {} 
\sphinxAtStartPar
Plot your recordings (separately or on the same axes — be sure to label two different recordings if overlaid). Make sure your axes are labeled with units. You can create these plots in \sphinxhref{https://bipn145.github.io/LabChart/ImportingPython.html}{Python} or \sphinxhref{https://bipn145.github.io/LabChart/ImportingExcel.html}{Excel}.

\item {} 
\sphinxAtStartPar
Copy your plots into a Word or Google document.

\item {} 
\sphinxAtStartPar
See your instructor’s Canvas assignment for detailed directions about what else you need to submit.

\end{enumerate}


\bigskip\hrule\bigskip



\subsection{Troubleshooting}
\label{\detokenize{StringNervousSystems:troubleshooting}}

\begin{savenotes}\sphinxattablestart
\sphinxthistablewithglobalstyle
\centering
\begin{tabulary}{\linewidth}[t]{TTT}
\sphinxtoprule
\sphinxstyletheadfamily 
\sphinxAtStartPar
Observation
&\sphinxstyletheadfamily 
\sphinxAtStartPar
Likely issue
&\sphinxstyletheadfamily 
\sphinxAtStartPar
Possible solution
\\
\sphinxmidrule
\sphinxtableatstartofbodyhook
\sphinxAtStartPar
You can see the stimulus marker on Channel 1, but no stimulus artifact on Channel 3
&
\sphinxAtStartPar
Your string may not be soaked in saline well enough
&
\sphinxAtStartPar
Add saline to the string
\\
\sphinxhline
\sphinxAtStartPar
The recording trace is cut off
&
\sphinxAtStartPar
Your visualization needs to be re\sphinxhyphen{}scaled, or your range is too small
&
\sphinxAtStartPar
Right\sphinxhyphen{}click on channel and choose Auto Scale (see Visualization tips above); change the range on your recording channel
\\
\sphinxhline
\sphinxAtStartPar
Your string recording trace has clear 60 Hz noise
&
\sphinxAtStartPar
You’re not grounded, or there are ground loops
&
\sphinxAtStartPar
Make sure your Faraday cage and ground electrode are both grounded to the back of the PowerLab. If this doesn’t help, consult an instructor.
\\
\sphinxbottomrule
\end{tabulary}
\sphinxtableafterendhook\par
\sphinxattableend\end{savenotes}

\sphinxstepscope


\chapter{Earthworm Electrophysiology}
\label{\detokenize{Earthworm:earthworm-electrophysiology}}\label{\detokenize{Earthworm::doc}}

\section{Background \& Goals}
\label{\detokenize{Earthworm:background-goals}}

\subsection{Earthworm essentials}
\label{\detokenize{Earthworm:earthworm-essentials}}
\sphinxAtStartPar
Earthworms are a member of the Annelida phylum, which includes segmented worms such as the ragworm and our soon\sphinxhyphen{}to\sphinxhyphen{}be\sphinxhyphen{}friend, the leech. We’ll be working with earthworms in the Lumbricidae family.

\sphinxAtStartPar
Earthworms are laid out with a mouth at one end, and an anus at the other. In the middle is the clitellum, a reproductive gland that can secrete a material that helps them form a cocoon. The clitellum is only visible in older worms that are ready to reproduce.

\sphinxAtStartPar
Earthworms have a closed nervous system, with one main blood vessel running down the length of their body. They breathe through their skin, so it is important to keep it moist. They crawl using muscles that are located just underneath the epidermis.

\sphinxAtStartPar



\subsection{Earthworm nervous system}
\label{\detokenize{Earthworm:earthworm-nervous-system}}
\sphinxAtStartPar
Earthworms have a pretty simple brain: it’s actually a fused pair of nerve ganglia, mostly located in their third segment. The earthworm has one main nerve cord on its ventral (bottom) side with three fibers running through it: one medial giant fiber, and two lateral giant fibers.

\sphinxAtStartPar
These fibers are big bundles of axons. They serve as high\sphinxhyphen{}speed motor pathways that enable the worm to send movement signals throughout its body for quick, stereotyped reflexes.

\sphinxAtStartPar
Each giant fiber is made up of large individual cells (one per segment) that are electrically coupled to each other through gap junctions. This results in the rapid conduction of action potentials from cell to cell so that each giant fiber behaves as though it were a single axon. There are cross\sphinxhyphen{}connections between the medial and lateral fibers, so they often function as one unit.

\sphinxAtStartPar
\sphinxstylestrong{These fibers differ in important ways.} The medial giant fiber transmits information from skin cells in the anterior (mouth) end of the worm, and the lateral giant fibers transmit information from the skin cells of the posterior end of the worm.

\sphinxAtStartPar
\sphinxstylestrong{It’s quite easy to listen in to the activity of these large fibers, even from the skin of the earthworm.} Importantly, these fibers fire action potentials very rarely. This type of nervous system activity is often called sparse coding. That means that we should be able to elicit one action potential and cleanly record it.

\sphinxAtStartPar

\sphinxstyleemphasis{Images from \sphinxhref{https://docs.backyardbrains.com/retired/experiments/speed}{Backyard Brains}, © 2009–2026, used under a Creative Commons Attribution\sphinxhyphen{}ShareAlike license.}


\subsection{Properties of action potentials}
\label{\detokenize{Earthworm:properties-of-action-potentials}}
\sphinxAtStartPar
In this lab, we’ll determine the \sphinxstylestrong{threshold}, \sphinxstylestrong{rheobase \& chronaxie}, \sphinxstylestrong{refractory period}, and \sphinxstylestrong{conduction velocity} for the earthworm giant fiber system. We’re recording compound action potentials (CAPs) from many axons in these fiber bundles, but the underlying principles are the same.

\sphinxAtStartPar
The \sphinxstylestrong{threshold} is defined as the minimum stimulus voltage that elicits an action potential 50\% of the time. We’ll determine this with short stimuli. Stimulus voltages that don’t elicit an action potential are called subthreshold; stimulus voltages that do elicit an action potential are called suprathreshold.

\sphinxAtStartPar
In some neurons, we can elicit an action potential with a lower stimulus voltage if the stimulus length is longer. A longer stimulus length enables charge to build up that can ultimately elicit an action potential. However, there is a minimum stimulus voltage needed, even at infinite stimulus durations. The \sphinxstylestrong{rheobase} is defined as the minimum stimulus voltage that can elicit an action potential with very long durations of the stimulus. The stimulus duration at 2× the rheobase is called the \sphinxstylestrong{chronaxie}. Chronaxie is a useful measure of the excitability of a nerve — the most excitable nerves have the smallest chronaxie. We can plot this data in a strength duration curve:

\sphinxAtStartPar
\sphinxincludegraphics{{bahringbauer2014}.png}

\sphinxAtStartPar
The \sphinxstylestrong{absolute refractory period} is the minimum amount of time needed for a neuron to fire a second action potential. This is caused by the inactivation of sodium channels after an action potential — it takes time for them to close, to then be reopened by a depolarizing stimulus. Since we’re recording from many axons, your measured refractory period will be more variable than for a single axon. We’ll define the absolute refractory period as the time when the amplitude of the second CAP is 30\% of the first.

\sphinxAtStartPar
After an action potential, the membrane of the axon is also hyperpolarized, due to the slowness of K⁺ channels closing. So, there is a period of time where the neuron requires more voltage to fire an action potential. This period is called the \sphinxstylestrong{relative refractory period}. In this lab, we’ll identify it as the period where the amplitude of the second CAP is 30–90\% of the first.

\sphinxAtStartPar
\sphinxstylestrong{In these experiments you will:}
\begin{itemize}
\item {} 
\sphinxAtStartPar
Stimulate the ventral nerve cord of the earthworm to elicit an action potential

\item {} 
\sphinxAtStartPar
Determine the voltage that is needed to recruit the lateral giant fibers

\item {} 
\sphinxAtStartPar
Estimate the absolute and relative refractory period between action potentials

\item {} 
\sphinxAtStartPar
Calculate the speed of an action potential

\item {} 
\sphinxAtStartPar
Propose your own hypothesis about conduction velocity

\item {} 
\sphinxAtStartPar
Conduct your experiment to test your hypothesis

\end{itemize}


\subsection{References}
\label{\detokenize{Earthworm:references}}\begin{itemize}
\item {} 
\sphinxAtStartPar
Bähring, R. \& Bauer C.K. (2014) Easy method to examine single nerve fiber excitability and conduction parameters using intact nonanesthetized earthworms. \sphinxstyleemphasis{Advances in Physiology Education} 38(3): 253–264.

\item {} 
\sphinxAtStartPar
Bullock, T.H. (1945) Functional organization of the giant fibre system of Lumbricus. \sphinxstyleemphasis{J Neurophysiol} 8:54–71.

\item {} 
\sphinxAtStartPar
Drewes CD, Landa KB, McFall JL (1978) Giant nerve fibre activity in intact, freely moving earthworms. \sphinxstyleemphasis{J Exp Biol} 72:217\sphinxhyphen{}227.

\item {} 
\sphinxAtStartPar
Kladt et al. (2010) Teaching basic neurophysiology in intact earthworms. \sphinxstyleemphasis{JUNE} 9(1):20\sphinxhyphen{}35.

\item {} 
\sphinxAtStartPar
Rudell \& Fox (1991) The Propagation potential: An Axonal response with implications for scalp\sphinxhyphen{}recorded EEG. \sphinxstyleemphasis{Biophysical Journal} 60: 556\sphinxhyphen{}567.

\item {} 
\sphinxAtStartPar
Shannon et al. (2014) Portable conduction velocity experiments using earthworms for the college and high school neuroscience teaching laboratory. \sphinxstyleemphasis{Advances in Physiology Education} 38(1): 62\sphinxhyphen{}70.

\end{itemize}


\bigskip\hrule\bigskip



\section{Lab Protocol}
\label{\detokenize{Earthworm:lab-protocol}}
\sphinxAtStartPar
There are multiple days dedicated to these earthworm experiments. You are encouraged to complete as many of the experiments below while you have a functioning earthworm.


\begin{savenotes}\sphinxattablestart
\sphinxthistablewithglobalstyle
\centering
\begin{tabulary}{\linewidth}[t]{TTT}
\sphinxtoprule
\sphinxstyletheadfamily 
\sphinxAtStartPar
\sphinxstylestrong{Supplies}
&\sphinxstyletheadfamily 
\sphinxAtStartPar
\sphinxstylestrong{Cables \& connectors}
&\sphinxstyletheadfamily 
\sphinxAtStartPar
\sphinxstylestrong{Solutions}
\\
\sphinxmidrule
\sphinxtableatstartofbodyhook
\sphinxAtStartPar
Earthworm
&
\sphinxAtStartPar
BNC to single banana adapter (2)
&
\sphinxAtStartPar
Earthworm saline
\\
\sphinxhline
\sphinxAtStartPar
Dissection tray
&
\sphinxAtStartPar
Banana plug cords (2)
&
\sphinxAtStartPar
Earthworm anesthetic (10\% ethanol in Ringers)
\\
\sphinxhline
\sphinxAtStartPar
PowerLab 26T, connected to computer
&
\sphinxAtStartPar
Alligator clip adapters (2)
&
\sphinxAtStartPar

\\
\sphinxhline
\sphinxAtStartPar
BioAmp
&
\sphinxAtStartPar
Needle electrodes
&
\sphinxAtStartPar

\\
\sphinxhline
\sphinxAtStartPar
Sewing pins (3)
&
\sphinxAtStartPar

&
\sphinxAtStartPar

\\
\sphinxhline
\sphinxAtStartPar
Ruler
&
\sphinxAtStartPar

&
\sphinxAtStartPar

\\
\sphinxhline
\sphinxAtStartPar
Thermometer
&
\sphinxAtStartPar

&
\sphinxAtStartPar

\\
\sphinxhline
\sphinxAtStartPar
Faraday Cage
&
\sphinxAtStartPar

&
\sphinxAtStartPar

\\
\sphinxhline
\sphinxAtStartPar
Large Petri Dish
&
\sphinxAtStartPar

&
\sphinxAtStartPar

\\
\sphinxhline
\sphinxAtStartPar
Eyedropper
&
\sphinxAtStartPar

&
\sphinxAtStartPar

\\
\sphinxbottomrule
\end{tabulary}
\sphinxtableafterendhook\par
\sphinxattableend\end{savenotes}


\bigskip\hrule\bigskip



\subsection{I. Preparation}
\label{\detokenize{Earthworm:i-preparation}}

\subsubsection{Setting up the PowerLab 26T}
\label{\detokenize{Earthworm:setting-up-the-powerlab-26t}}
\sphinxAtStartPar
Our first step is making sure our PowerLab \& Amplifier are wired up correctly.

\sphinxAtStartPar
Make sure the PowerLab is turned off.
\begin{enumerate}
\sphinxsetlistlabels{\arabic}{enumi}{enumii}{}{.}%
\item {} 
\sphinxAtStartPar
\sphinxstylestrong{Connect a single BNC adapter to the (+) Output} of the PowerLab. Plug a red banana plug cable into the adapter.

\item {} 
\sphinxAtStartPar
\sphinxstylestrong{Connect a single BNC to banana adapter to the (−) Output.} Plug a black banana plug cable into the adapter.

\item {} 
\sphinxAtStartPar
Make sure that the BioAmp is connected to the PowerLab, with three different needle electrodes plugged into it.

\item {} 
\sphinxAtStartPar
Once you’re confident everything is connected correctly, turn on the PowerLab.

\end{enumerate}

\sphinxAtStartPar
\sphinxincludegraphics{{powerlab_front}.png}
\sphinxstyleemphasis{Image courtesy of ADInstruments ©}


\subsubsection{Setting up LabChart 8}
\label{\detokenize{Earthworm:setting-up-labchart-8}}\begin{enumerate}
\sphinxsetlistlabels{\arabic}{enumi}{enumii}{}{.}%
\item {} 
\sphinxAtStartPar
\sphinxstylestrong{Open LabChart 8 with the settings file you saved during the string experiment.} Most settings are saved within this settings file. However, it’s always good to ensure that you have the correct settings before beginning.
\begin{quote}

\sphinxAtStartPar
\sphinxstylestrong{Note:} If for some reason you cannot access your settings file, you can find one at \sphinxurl{http://bit.ly/labchart}. Use the “Earthworm bioAmp Settings.”
\end{quote}

\item {} 
\sphinxAtStartPar
\sphinxstylestrong{Go to Setup > Channel Settings} and confirm that you have three recording channels, and that the third (“BioAmp Recording”) is connected to the BioAmp.

\item {} 
\sphinxAtStartPar
\sphinxstylestrong{Go to Setup > Stimulator} and change the Stimulator settings as follows. Make sure your Marker Channel is Channel 1.


\begin{savenotes}\sphinxattablestart
\sphinxthistablewithglobalstyle
\centering
\begin{tabulary}{\linewidth}[t]{TT}
\sphinxtoprule
\sphinxstyletheadfamily 
\sphinxAtStartPar
Parameter
&\sphinxstyletheadfamily 
\sphinxAtStartPar
Value
\\
\sphinxmidrule
\sphinxtableatstartofbodyhook
\sphinxAtStartPar
Start Delay
&
\sphinxAtStartPar
0.1 ms
\\
\sphinxhline
\sphinxAtStartPar
Repeats
&
\sphinxAtStartPar
1
\\
\sphinxhline
\sphinxAtStartPar
Max Repeat Rate
&
\sphinxAtStartPar
1 Hz
\\
\sphinxhline
\sphinxAtStartPar
Pulse Height
&
\sphinxAtStartPar
0.2 V
\\
\sphinxhline
\sphinxAtStartPar
Pulse Width
&
\sphinxAtStartPar
0.2 ms
\\
\sphinxbottomrule
\end{tabulary}
\sphinxtableafterendhook\par
\sphinxattableend\end{savenotes}

\end{enumerate}


\subsubsection{Setting up your earthworm}
\label{\detokenize{Earthworm:setting-up-your-earthworm}}
\sphinxAtStartPar
\sphinxincludegraphics{{string_electrodes}.png}
\begin{quote}

\sphinxAtStartPar
\sphinxstylestrong{Note:} Throughout every step of this procedure, be careful not to stretch the earthworm. This will damage the nerve and your experiment will not work.
\end{quote}
\begin{enumerate}
\sphinxsetlistlabels{\arabic}{enumi}{enumii}{}{.}%
\item {} 
\sphinxAtStartPar
Obtain a large, lively earthworm. Rinse it with saline if it is dirty.

\item {} 
\sphinxAtStartPar
\sphinxstylestrong{Anesthetize the earthworm} by placing it in the petri dish with 10\% ethanol (in Earthworm Ringer’s solution) for 5 minutes.
\begin{quote}

\sphinxAtStartPar
\sphinxstylestrong{Note:} Do not leave your earthworm unattended in anesthesia. With too little anesthesia, the earthworm will move around during the experiment, and the resulting muscle electrical activity (electromyogram) will drown out the small neural electrical signals you are interested in. Too much anesthesia and the nerves will not fire.

\sphinxAtStartPar
\sphinxstylestrong{Note \#2:} Your earthworm may still react to pressure after 5 minutes — you can proceed with the experiment as long as you are able to handle the worm and pin it down.
\end{quote}

\item {} 
\sphinxAtStartPar
\sphinxstylestrong{When your earthworm is anesthetized enough to pin it down,} remove the earthworm from the ethanol and place it on the dissection tray with the dorsal (darker side) up.

\item {} 
\sphinxAtStartPar
Pin the earthworm to the tray using one sewing pin on each end of the worm.

\item {} 
\sphinxAtStartPar
Place the ruler down next to your worm.

\item {} 
\sphinxAtStartPar
\sphinxstylestrong{Insert a third sewing pin into the anterior end of the earthworm,} \textasciitilde{}1 cm from the first pin. Connect the anode \& cathode stimulating electrodes to the first two pins.
\begin{quote}

\sphinxAtStartPar
\sphinxstylestrong{Note:} The anterior end is the end nearest the clitellum, a smooth, non\sphinxhyphen{}segmented band.
\end{quote}

\item {} 
\sphinxAtStartPar
To keep the worm anesthetized, apply a few drops of the earthworm anesthetic along the length of the earthworm with the eyedropper. Soak up any excess fluid.
\begin{quote}

\sphinxAtStartPar
\sphinxstylestrong{Note:} Make sure there isn’t a large puddle of solution — too much may interfere with your recording.
\end{quote}

\item {} 
\sphinxAtStartPar
Insert the three recording needle electrodes into the worm (same configuration as the string experiment!).
\begin{quote}

\sphinxAtStartPar
\sphinxstylestrong{Note:} Make sure no wire is touching any other wire. Be careful when positioning the wires to avoid damaging the worm.

\sphinxAtStartPar
\sphinxstylestrong{Note \#2:} It helps to place your recording electrodes as lateral as possible, while still holding down the worm.
\end{quote}

\end{enumerate}


\bigskip\hrule\bigskip



\subsection{II. Determining the threshold of your fibers ★}
\label{\detokenize{Earthworm:ii-determining-the-threshold-of-your-fibers}}
\sphinxAtStartPar
In this experiment, you will send a stimulus into the earthworm while recording from it to determine the threshold voltage of the medial and lateral giant fibers.


\subsubsection{Determine the threshold to activate the Medial Giant Fiber}
\label{\detokenize{Earthworm:determine-the-threshold-to-activate-the-medial-giant-fiber}}\begin{enumerate}
\sphinxsetlistlabels{\arabic}{enumi}{enumii}{}{.}%
\item {} 
\sphinxAtStartPar
\sphinxstylestrong{Open your Input Amplifier (BioAmp)} for your Action Potential channel to make sure that your signal looks reasonable. You shouldn’t see 60 Hz noise, but it should look like a physiological signal. (If you have noise, see Troubleshooting.)

\item {} 
\sphinxAtStartPar
\sphinxstylestrong{Open the Stimulator Panel} by going to the Setup menu. Your stimulator should start with:
\begin{itemize}
\item {} 
\sphinxAtStartPar
Max Repeat Rate: 1 Hz

\item {} 
\sphinxAtStartPar
Pulse Height: 0.2 V

\item {} 
\sphinxAtStartPar
Pulse Width: 0.2 ms

\end{itemize}
\begin{quote}

\sphinxAtStartPar
\sphinxstylestrong{Note:} Leave the Stimulator Panel open — we’ll need it throughout the experiment.
\end{quote}

\item {} 
\sphinxAtStartPar
\sphinxstylestrong{Press > Start.} LabChart should be in Scope View and display one 20 ms sweep every time you click Start. For each 20 ms sweep, a stimulus pulse is generated through the PowerLab Stimulator outputs.

\item {} 
\sphinxAtStartPar
\sphinxstylestrong{Examine the recording} — you should see a quick deflection very shortly after your stimulus pulse. The deflection just after the start of the sweep is your stimulus artifact, caused by spread of the stimulus voltage to the recording electrodes.
\begin{quote}

\sphinxAtStartPar
\sphinxstylestrong{Note:} The stimulus artifact can be very large in amplitude, and can exceed the range of the recording amplifier. It is okay if it is cut off.
\end{quote}

\item {} 
\sphinxAtStartPar
\sphinxstylestrong{Change the Pulse Height to 1 V.}

\item {} 
\sphinxAtStartPar
Each time you change the stimulus, click Start and a new page will appear to the left of the Scope View Window. This is the Page Explorer. \sphinxstylestrong{Label the individual pages} of your recording as you go along.

\item {} 
\sphinxAtStartPar
\sphinxstylestrong{Increase the Pulse Height in 0.1 V steps} until you see an inflection after the stimulus artifact. This means you’ve successfully elicited a response from the medial giant fiber! Note the stimulus voltage as an estimate of the threshold stimulus strength.
\begin{quote}

\sphinxAtStartPar
\sphinxstylestrong{Note:} If you can’t see anything, right\sphinxhyphen{}click on the channel to change the scaling (or use the − and + buttons on the bottom left of the channel).

\sphinxAtStartPar
\sphinxstylestrong{Note \#2:} If no response is seen with a stimulus of 3.5 V or greater, contact your instructor for assistance.
\end{quote}

\item {} 
\sphinxAtStartPar
\sphinxstylestrong{Find a more accurate threshold} by first reducing the stimulus by 0.2 V (or more) to diminish your response. Then, increase the stimulus by 0.05 V between each sweep. You should report a threshold (Table 1) to the hundredth decimal point (e.g. 1.25 V).

\item {} 
\sphinxAtStartPar
Save your data, but do not close the file.

\end{enumerate}


\subsubsection{Determine the threshold to recruit the Lateral Giant Fibers}
\label{\detokenize{Earthworm:determine-the-threshold-to-recruit-the-lateral-giant-fibers}}
\sphinxAtStartPar
In the next stage, you will attempt to elicit an action potential in the lateral giant fibers.
\begin{quote}

\sphinxAtStartPar
\sphinxstylestrong{Note:} Remember to keep the earthworm anesthetized!
\end{quote}
\begin{enumerate}
\sphinxsetlistlabels{\arabic}{enumi}{enumii}{}{.}%
\item {} 
\sphinxAtStartPar
Set the Pulse Height to your threshold for the MGF.

\item {} 
\sphinxAtStartPar
Keep increasing the stimulus in 0.1 V steps until you see a second inflection in your recorded trace.

\item {} 
\sphinxAtStartPar
Once you see a second inflection, follow the same procedure as above to determine the threshold where you can elicit an action potential from the lateral giant fiber about 50\% of the time. Record the stimulus threshold in Table 1 below.

\item {} 
\sphinxAtStartPar
Save your data, but do not close the file.
\begin{quote}

\sphinxAtStartPar
\sphinxstylestrong{Note:} In some worms, the median and lateral fiber responses may not be separated — they will appear to have the same threshold. Other worms may not show the second response, possibly because the lateral giant fibers are damaged.
\end{quote}

\item {} 
\sphinxAtStartPar
Save your LabChart file.

\end{enumerate}

\sphinxAtStartPar
\sphinxstylestrong{Table 1.} MGF and LGF compound action potential (CAP) properties.


\begin{savenotes}\sphinxattablestart
\sphinxthistablewithglobalstyle
\centering
\begin{tabulary}{\linewidth}[t]{TTTT}
\sphinxtoprule

\sphinxAtStartPar

&\sphinxstyletheadfamily 
\sphinxAtStartPar
Threshold
&\sphinxstyletheadfamily 
\sphinxAtStartPar
Amplitude
&\sphinxstyletheadfamily 
\sphinxAtStartPar
Latency
\\
\sphinxmidrule
\sphinxtableatstartofbodyhook
\sphinxAtStartPar
\sphinxstylestrong{MGF}
&
\sphinxAtStartPar

&
\sphinxAtStartPar

&
\sphinxAtStartPar

\\
\sphinxhline
\sphinxAtStartPar
\sphinxstylestrong{LGF}
&
\sphinxAtStartPar

&
\sphinxAtStartPar

&
\sphinxAtStartPar

\\
\sphinxbottomrule
\end{tabulary}
\sphinxtableafterendhook\par
\sphinxattableend\end{savenotes}

\begin{sphinxadmonition}{note}{Questions for reflection}
\begin{itemize}
\item {} 
\sphinxAtStartPar
Did the shape or amplitude of the CAPs change with more stimulation voltage? What does this tell us about the action potential that is firing?

\item {} 
\sphinxAtStartPar
Which fiber fires an action potential sooner (with a shorter latency)?

\item {} 
\sphinxAtStartPar
Why would the MGF and LGF latencies differ?

\end{itemize}
\end{sphinxadmonition}


\bigskip\hrule\bigskip



\subsection{III. Calculate the rheobase ★}
\label{\detokenize{Earthworm:iii-calculate-the-rheobase}}
\sphinxAtStartPar
In Part III, we’ll change the stimulus amplitude and duration to calculate the rheobase and chronaxie of the earthworm medial giant fiber (read the Background if you forget what these are).

\sphinxAtStartPar
In the previous experiment, we gave a stimulus for 0.2 ms. Now, we’ll modify the stimulus duration and amplitude to determine the minimum stimulus requirement, regardless of duration.
\begin{enumerate}
\sphinxsetlistlabels{\arabic}{enumi}{enumii}{}{.}%
\item {} 
\sphinxAtStartPar
Write your MGF threshold in the column for 0.20 ms stimulus duration in Table 2.

\item {} 
\sphinxAtStartPar
Change your Pulse Height to the threshold for your MGF.

\item {} 
\sphinxAtStartPar
Systematically change the duration of your stimulus (Pulse Width) to the values in the table below. As necessary, change the stimulus voltage to successfully get a CAP from your MGF.

\item {} 
\sphinxAtStartPar
For each duration, record the stimulus voltage (Pulse Height) values that successfully elicit a CAP in Table 2.
\begin{quote}

\sphinxAtStartPar
\sphinxstylestrong{Note:} Be sure to label your pages in Scope view accordingly. You don’t need to save the pages that do not elicit an action potential.
\end{quote}

\end{enumerate}

\sphinxAtStartPar
\sphinxstylestrong{Table 2.} Data for strength\sphinxhyphen{}duration curve.


\begin{savenotes}\sphinxattablestart
\sphinxthistablewithglobalstyle
\centering
\begin{tabulary}{\linewidth}[t]{TTTTTTTT}
\sphinxtoprule
\sphinxstyletheadfamily 
\sphinxAtStartPar
Duration (ms)
&\sphinxstyletheadfamily 
\sphinxAtStartPar
0.06
&\sphinxstyletheadfamily 
\sphinxAtStartPar
0.08
&\sphinxstyletheadfamily 
\sphinxAtStartPar
0.10
&\sphinxstyletheadfamily 
\sphinxAtStartPar
0.20
&\sphinxstyletheadfamily 
\sphinxAtStartPar
0.40
&\sphinxstyletheadfamily 
\sphinxAtStartPar
0.60
&\sphinxstyletheadfamily 
\sphinxAtStartPar
1.00
\\
\sphinxmidrule
\sphinxtableatstartofbodyhook
\sphinxAtStartPar
Stimulus Amplitude (V)
&
\sphinxAtStartPar

&
\sphinxAtStartPar

&
\sphinxAtStartPar

&
\sphinxAtStartPar

&
\sphinxAtStartPar

&
\sphinxAtStartPar

&
\sphinxAtStartPar

\\
\sphinxbottomrule
\end{tabulary}
\sphinxtableafterendhook\par
\sphinxattableend\end{savenotes}


\bigskip\hrule\bigskip



\subsection{IV. Determine the refractory period}
\label{\detokenize{Earthworm:iv-determine-the-refractory-period}}
\sphinxAtStartPar
In the next step, you’ll determine the shortest time between stimuli where you can elicit two different CAPs from the MGF. In other words, we’ll determine the relative and absolute refractory period of these fibers.
\begin{enumerate}
\sphinxsetlistlabels{\arabic}{enumi}{enumii}{}{.}%
\item {} 
\sphinxAtStartPar
Set your pulse height so that you can comfortably elicit a response from the MGF.

\item {} 
\sphinxAtStartPar
\sphinxstylestrong{In Setup > Stimulator, change Repeats to 2} in order to play multiple pulses.

\item {} 
\sphinxAtStartPar
We also need to change the time between each stimulus. PowerLab thinks of this interval as Hz. To start, we want 15 ms between each stimulus.
\begin{quote}

\sphinxAtStartPar
\sphinxstylestrong{Note:} Remember that Hertz (Hz) means events per second. So, 1 Hz = 1 event per second, and 10 Hz = 10 events per second. 1 ÷ frequency (in Hz) = time between stimuli (in seconds).
\end{quote}

\end{enumerate}

\begin{sphinxadmonition}{note}{Questions for reflection}
\begin{itemize}
\item {} 
\sphinxAtStartPar
How many milliseconds would there be between each stimulus at 10 Hz stimulation?

\item {} 
\sphinxAtStartPar
For 15 ms between each stimulus, what should the frequency of stimulation (in Hz) be?

\end{itemize}
\end{sphinxadmonition}
\begin{enumerate}
\sphinxsetlistlabels{\arabic}{enumi}{enumii}{}{.}%
\setcounter{enumi}{3}
\item {} 
\sphinxAtStartPar
Fill out this chart with the corresponding Hz values for each interval length:

\end{enumerate}


\begin{savenotes}\sphinxattablestart
\sphinxthistablewithglobalstyle
\centering
\begin{tabulary}{\linewidth}[t]{TTTTTTTTTTTTTTTT}
\sphinxtoprule
\sphinxstyletheadfamily 
\sphinxAtStartPar
Interval (ms)
&\sphinxstyletheadfamily 
\sphinxAtStartPar
15
&\sphinxstyletheadfamily 
\sphinxAtStartPar
14
&\sphinxstyletheadfamily 
\sphinxAtStartPar
13
&\sphinxstyletheadfamily 
\sphinxAtStartPar
12
&\sphinxstyletheadfamily 
\sphinxAtStartPar
11
&\sphinxstyletheadfamily 
\sphinxAtStartPar
10
&\sphinxstyletheadfamily 
\sphinxAtStartPar
9
&\sphinxstyletheadfamily 
\sphinxAtStartPar
8
&\sphinxstyletheadfamily 
\sphinxAtStartPar
7
&\sphinxstyletheadfamily 
\sphinxAtStartPar
6
&\sphinxstyletheadfamily 
\sphinxAtStartPar
5
&\sphinxstyletheadfamily 
\sphinxAtStartPar
4
&\sphinxstyletheadfamily 
\sphinxAtStartPar
3
&\sphinxstyletheadfamily 
\sphinxAtStartPar
2
&\sphinxstyletheadfamily 
\sphinxAtStartPar
1
\\
\sphinxmidrule
\sphinxtableatstartofbodyhook
\sphinxAtStartPar
Freq. (Hz)
&
\sphinxAtStartPar

&
\sphinxAtStartPar

&
\sphinxAtStartPar

&
\sphinxAtStartPar

&
\sphinxAtStartPar

&
\sphinxAtStartPar

&
\sphinxAtStartPar

&
\sphinxAtStartPar

&
\sphinxAtStartPar

&
\sphinxAtStartPar

&
\sphinxAtStartPar

&
\sphinxAtStartPar

&
\sphinxAtStartPar

&
\sphinxAtStartPar

&
\sphinxAtStartPar

\\
\sphinxbottomrule
\end{tabulary}
\sphinxtableafterendhook\par
\sphinxattableend\end{savenotes}
\begin{enumerate}
\sphinxsetlistlabels{\arabic}{enumi}{enumii}{}{.}%
\setcounter{enumi}{4}
\item {} 
\sphinxAtStartPar
\sphinxstylestrong{Change your Max Repeat Rate} to the equivalent frequency for 15 ms intervals.

\item {} 
\sphinxAtStartPar
Change your Pulse Width to 0.2 ms and your Pulse Height to the threshold for your MGF.

\item {} 
\sphinxAtStartPar
\sphinxstylestrong{Press > Start} to record a trial with two stimuli, separated by 15 ms.

\item {} 
\sphinxAtStartPar
Right\sphinxhyphen{}click and choose \sphinxstylestrong{Add Marker > Clamp to Trace} to place a marker at the moment the stimulus occurs. Move your cursor to the peak of the action potential. Look at the ΔV on the right\sphinxhyphen{}hand side to see the change in voltage from that point.

\item {} 
\sphinxAtStartPar
\sphinxstylestrong{Measure the amplitude of your first and second action potentials.} Determine how strong the second was relative to the first (e.g. if AP \#2 amplitude was half as high as \#1, it would be 50\%). Record your response in row one of Table 3.

\item {} 
\sphinxAtStartPar
Decrease the duration between stimuli by 1 ms by changing the Repeat Rate to the Hz for a 14 ms interval.

\item {} 
\sphinxAtStartPar
\sphinxstylestrong{Record the amplitudes and corresponding percentages for each stimulus interval in Table 3.}

\item {} 
\sphinxAtStartPar
At some point, AP \#2 will diminish below 30\% of AP \#1. We can consider this our absolute refractory period.
\begin{quote}

\sphinxAtStartPar
\sphinxstylestrong{Note:} Depending on the latency of your MGF CAP, it may be difficult to see a CAP with a stimulus interval of 1 or 2 ms. Reduce the stimulus interval to as short as you can before the MGF CAP is cut off by the second stimulus. If the shortest stimulus interval is >3 ms, you should move your pins closer together.
\end{quote}

\item {} 
\sphinxAtStartPar
Save your LabChart file.
\begin{quote}

\sphinxAtStartPar
\sphinxstylestrong{Note:} If you have data that doesn’t seem quite right, you should repeat your experiment. You must be upfront about your number of repeats in your lab report, comprehensively reporting all of the data that you collected.
\end{quote}

\end{enumerate}

\sphinxAtStartPar
\sphinxstylestrong{Table 3.} Amplitudes of CAPs elicited by stimuli with decreasing stimulus intervals.


\begin{savenotes}\sphinxattablestart
\sphinxthistablewithglobalstyle
\centering
\begin{tabulary}{\linewidth}[t]{TTTT}
\sphinxtoprule
\sphinxstyletheadfamily 
\sphinxAtStartPar
Stimulus Interval (ms)
&\sphinxstyletheadfamily 
\sphinxAtStartPar
\#1 CAP Amplitude
&\sphinxstyletheadfamily 
\sphinxAtStartPar
\#2 CAP Amplitude
&\sphinxstyletheadfamily 
\sphinxAtStartPar
\% \#2/\#1
\\
\sphinxmidrule
\sphinxtableatstartofbodyhook
\sphinxAtStartPar
15
&
\sphinxAtStartPar

&
\sphinxAtStartPar

&
\sphinxAtStartPar

\\
\sphinxhline
\sphinxAtStartPar
14
&
\sphinxAtStartPar

&
\sphinxAtStartPar

&
\sphinxAtStartPar

\\
\sphinxhline
\sphinxAtStartPar
13
&
\sphinxAtStartPar

&
\sphinxAtStartPar

&
\sphinxAtStartPar

\\
\sphinxhline
\sphinxAtStartPar
12
&
\sphinxAtStartPar

&
\sphinxAtStartPar

&
\sphinxAtStartPar

\\
\sphinxhline
\sphinxAtStartPar
11
&
\sphinxAtStartPar

&
\sphinxAtStartPar

&
\sphinxAtStartPar

\\
\sphinxhline
\sphinxAtStartPar
10
&
\sphinxAtStartPar

&
\sphinxAtStartPar

&
\sphinxAtStartPar

\\
\sphinxhline
\sphinxAtStartPar
9
&
\sphinxAtStartPar

&
\sphinxAtStartPar

&
\sphinxAtStartPar

\\
\sphinxhline
\sphinxAtStartPar
8
&
\sphinxAtStartPar

&
\sphinxAtStartPar

&
\sphinxAtStartPar

\\
\sphinxhline
\sphinxAtStartPar
7
&
\sphinxAtStartPar

&
\sphinxAtStartPar

&
\sphinxAtStartPar

\\
\sphinxhline
\sphinxAtStartPar
6
&
\sphinxAtStartPar

&
\sphinxAtStartPar

&
\sphinxAtStartPar

\\
\sphinxhline
\sphinxAtStartPar
5
&
\sphinxAtStartPar

&
\sphinxAtStartPar

&
\sphinxAtStartPar

\\
\sphinxhline
\sphinxAtStartPar
4
&
\sphinxAtStartPar

&
\sphinxAtStartPar

&
\sphinxAtStartPar

\\
\sphinxhline
\sphinxAtStartPar
3
&
\sphinxAtStartPar

&
\sphinxAtStartPar

&
\sphinxAtStartPar

\\
\sphinxhline
\sphinxAtStartPar
2
&
\sphinxAtStartPar

&
\sphinxAtStartPar

&
\sphinxAtStartPar

\\
\sphinxhline
\sphinxAtStartPar
1
&
\sphinxAtStartPar

&
\sphinxAtStartPar

&
\sphinxAtStartPar

\\
\sphinxbottomrule
\end{tabulary}
\sphinxtableafterendhook\par
\sphinxattableend\end{savenotes}

\begin{sphinxadmonition}{note}{Questions for reflection}
\begin{itemize}
\item {} 
\sphinxAtStartPar
What was the shortest duration between stimuli where the second action potential was less than 90\% but more than 30\% of the original CAP amplitude?

\item {} 
\sphinxAtStartPar
What was the shortest duration between stimuli where the second CAP was diminished below 30\%?

\item {} 
\sphinxAtStartPar
Which of these values is your absolute refractory period, and which is the relative refractory period?

\end{itemize}
\end{sphinxadmonition}


\bigskip\hrule\bigskip



\subsection{V. Compute the conduction velocity ★}
\label{\detokenize{Earthworm:v-compute-the-conduction-velocity}}
\sphinxAtStartPar
\sphinxstyleemphasis{(If you’re doing this on a different day, follow the same steps for preparing LabChart \& your earthworm as you did in the previous earthworm experiment.)}

\sphinxAtStartPar
\sphinxstylestrong{How far is your first recording pin (−) from your cathode (−) stimulation pin? Be sure to include units:} \_\_\_\_\_\_\_\_\_\_\_\_\_\_\_\_\_\_\_\_\_\_
\begin{enumerate}
\sphinxsetlistlabels{\arabic}{enumi}{enumii}{}{.}%
\item {} 
\sphinxAtStartPar
\sphinxstylestrong{In PowerLab, press Start.} Confirm that you have two recording channels.

\item {} 
\sphinxAtStartPar
Follow the same protocol as in the previous experiments to determine the threshold required to elicit an action potential from your medial giant fiber.

\end{enumerate}

\sphinxAtStartPar
\sphinxstylestrong{What is the stimulus voltage needed to elicit an MGF action potential from this worm?} \_\_\_\_\_\_\_\_\_\_\_\_\_\_\_\_\_\_\_\_\_\_


\subsubsection{Determine the MGF latency to spike}
\label{\detokenize{Earthworm:determine-the-mgf-latency-to-spike}}\begin{enumerate}
\sphinxsetlistlabels{\arabic}{enumi}{enumii}{}{.}%
\item {} 
\sphinxAtStartPar
Using the Marker tool, measure the latency (in seconds) from the peak of the stimulus artifact to your elicited action potential.

\item {} 
\sphinxAtStartPar
Repeat this 5 times and record your data in Table 4.

\end{enumerate}

\sphinxAtStartPar
\sphinxstylestrong{Table 4.} MGF conduction velocity data.


\begin{savenotes}\sphinxattablestart
\sphinxthistablewithglobalstyle
\centering
\begin{tabulary}{\linewidth}[t]{TTTTT}
\sphinxtoprule
\sphinxstyletheadfamily 
\sphinxAtStartPar
Trial
&\sphinxstyletheadfamily 
\sphinxAtStartPar
MGF Baseline Latency (s)
&\sphinxstyletheadfamily 
\sphinxAtStartPar
Speed (m/s)
&\sphinxstyletheadfamily 
\sphinxAtStartPar
MGF Experimental Latency (s)
&\sphinxstyletheadfamily 
\sphinxAtStartPar
Speed (m/s)
\\
\sphinxmidrule
\sphinxtableatstartofbodyhook
\sphinxAtStartPar
1
&
\sphinxAtStartPar

&
\sphinxAtStartPar

&
\sphinxAtStartPar

&
\sphinxAtStartPar

\\
\sphinxhline
\sphinxAtStartPar
2
&
\sphinxAtStartPar

&
\sphinxAtStartPar

&
\sphinxAtStartPar

&
\sphinxAtStartPar

\\
\sphinxhline
\sphinxAtStartPar
3
&
\sphinxAtStartPar

&
\sphinxAtStartPar

&
\sphinxAtStartPar

&
\sphinxAtStartPar

\\
\sphinxhline
\sphinxAtStartPar
4
&
\sphinxAtStartPar

&
\sphinxAtStartPar

&
\sphinxAtStartPar

&
\sphinxAtStartPar

\\
\sphinxhline
\sphinxAtStartPar
5
&
\sphinxAtStartPar

&
\sphinxAtStartPar

&
\sphinxAtStartPar

&
\sphinxAtStartPar

\\
\sphinxbottomrule
\end{tabulary}
\sphinxtableafterendhook\par
\sphinxattableend\end{savenotes}
\begin{enumerate}
\sphinxsetlistlabels{\arabic}{enumi}{enumii}{}{.}%
\setcounter{enumi}{2}
\item {} 
\sphinxAtStartPar
Since we know when the AP occurred and the distance between your electrodes, we can calculate the speed of the action potential. Record these values in the “Speed” column of the table above. \sphinxstyleemphasis{(Reminder: speed = distance/time.)}

\end{enumerate}

\sphinxAtStartPar
\sphinxstylestrong{How fast, on average, is your action potential traveling down the MGF?} \_\_\_\_\_\_\_\_\_\_\_\_\_\_\_\_\_\_\_\_\_\_


\bigskip\hrule\bigskip



\subsection{VI. Design your own experiment to test the impact of temperature on MGF conduction velocity}
\label{\detokenize{Earthworm:vi-design-your-own-experiment-to-test-the-impact-of-temperature-on-mgf-conduction-velocity}}
\sphinxAtStartPar
For this final part, you will design your own experiment to test the impact of temperature on conduction velocity. We will provide ice and a thermometer, but the rest is up to you! If you are performing this experiment on a separate worm, be sure to measure the threshold for the MGF as described above.

\sphinxAtStartPar
Work with your group to determine a protocol for running your experiment — you’ll need to include this in the Methods section of your lab report. Think about what someone needs to know in order to replicate this experiment exactly as you did.

\sphinxAtStartPar
We will test the effect of \_\_\_\_\_\_\_\_\_\_\_\_\_\_\_\_\_\_\_\_\_\_\_\_\_\_\_\_\_\_ on \_\_\_\_\_\_\_\_\_\_\_\_\_\_\_\_\_\_\_\_\_\_\_\_\_\_\_\_\_\_.

\sphinxAtStartPar
We will do this by…


\subsubsection{Conduct your experiment}
\label{\detokenize{Earthworm:conduct-your-experiment}}
\sphinxAtStartPar
Work with your group to determine a protocol for running your experiment — you’ll need to include this in the Methods section of your lab report.

\sphinxAtStartPar
Record your results in the MGF Experimental AP Latency column of Table 4. You’ll need these later for the Results section of your lab report.


\bigskip\hrule\bigskip



\subsection{Troubleshooting}
\label{\detokenize{Earthworm:troubleshooting}}

\begin{savenotes}\sphinxattablestart
\sphinxthistablewithglobalstyle
\centering
\begin{tabulary}{\linewidth}[t]{TTT}
\sphinxtoprule
\sphinxstyletheadfamily 
\sphinxAtStartPar
Observation
&\sphinxstyletheadfamily 
\sphinxAtStartPar
Likely issue
&\sphinxstyletheadfamily 
\sphinxAtStartPar
Possible solution
\\
\sphinxmidrule
\sphinxtableatstartofbodyhook
\sphinxAtStartPar
The recording trace is flat
&
\sphinxAtStartPar
Something is disconnected
&
\sphinxAtStartPar
Check all of your connections
\\
\sphinxhline
\sphinxAtStartPar
I can’t see a stimulus trace in the Stimulus channel
&
\sphinxAtStartPar
You don’t have a marker channel set up
&
\sphinxAtStartPar
Go to Setup > Stimulator to set up the Marker Channel
\\
\sphinxhline
\sphinxAtStartPar
I can’t see a stimulus artifact in my action potential channel
&
\sphinxAtStartPar
Your electrodes may not be in the worm, or stimulus amplitude is too weak
&
\sphinxAtStartPar
Ensure all electrodes are in the worm (but not too medial), and try stimulating at 0.5 V
\\
\sphinxhline
\sphinxAtStartPar
My worm is bleeding
&
\sphinxAtStartPar
The electrode is through the blood vessel
&
\sphinxAtStartPar
Move the electrode more lateral; if you don’t get an AP after this, get a new worm
\\
\sphinxhline
\sphinxAtStartPar
There is a stimulus artifact, but no visible response to stimuli between 1–3 V
&
\sphinxAtStartPar
Your electrodes may not be placed well, or your worm may be damaged
&
\sphinxAtStartPar
Try to reposition the electrodes on a different place in the worm. If this doesn’t work, get a new worm.
\\
\sphinxhline
\sphinxAtStartPar
My recording just looks like a big up and down squiggle (like a sine wave)
&
\sphinxAtStartPar
You’re picking up 60 Hz noise
&
\sphinxAtStartPar
Ground everything (Faraday cage, ground recording pin, amplifier) to the ground pin on the back of the PowerLab. Averaging multiple trials can also get rid of noise.
\\
\sphinxbottomrule
\end{tabulary}
\sphinxtableafterendhook\par
\sphinxattableend\end{savenotes}


\bigskip\hrule\bigskip



\subsection{Possible opportunities for student projects}
\label{\detokenize{Earthworm:possible-opportunities-for-student-projects}}\begin{itemize}
\item {} 
\sphinxAtStartPar
Application of chemicals, nerve blockers, etc. to test for effects on the nervous system

\item {} 
\sphinxAtStartPar
Earthworm dissection \& recordings (as in Gunther, 1976 or 1972)

\item {} 
\sphinxAtStartPar
Demonstrating facilitation of conduction rate (Bullock, 1951)

\end{itemize}

\sphinxstepscope


\chapter{\sphinxstyleemphasis{Drosophila} Behavior \& Optogenetics}
\label{\detokenize{Drosophila/Introduction:drosophila-behavior-optogenetics}}\label{\detokenize{Drosophila/Introduction::doc}}

\section{Background \& Goals}
\label{\detokenize{Drosophila/Introduction:background-goals}}

\subsection{\sphinxstyleemphasis{Drosophila} as a model in neuroscience}
\label{\detokenize{Drosophila/Introduction:drosophila-as-a-model-in-neuroscience}}
\sphinxAtStartPar
With just over 100,000 neurons but incredible genetic tractability, the fruit fly is an important model organism in neuroscience. About 65\% of human disease\sphinxhyphen{}associated genes have a homolog in \sphinxstyleemphasis{Drosophila melanogaster} (Ugur et al., 2016). Like many model organisms, fruit flies are easy to reproduce, cheap, and easy to take care of in a lab. On top of that, they’ve been studied for more than a century. \sphinxstyleemphasis{Drosophila} have been used to study learning, memory, vision, sleep, addiction, courtship, and aggression, among many other behaviors.

\sphinxAtStartPar
In these labs, we’ll aim to characterize the gustatory, ethological, and locomotor behavior of transgenic fruit flies, solving a few mysteries along the way.


\subsection{What’s coming up}
\label{\detokenize{Drosophila/Introduction:what-s-coming-up}}
\sphinxAtStartPar
This module spans \sphinxstylestrong{two days} of lab:
\begin{itemize}
\item {} 
\sphinxAtStartPar
\sphinxstylestrong{{\hyperref[\detokenize{Drosophila/Day1::doc}]{\sphinxcrossref{\DUrole{std,std-doc}{Day 1: The Case of the Mislabeled Fly Vials}}}}} — three behavioral assays (proboscis extension reflex, ethological observation, and locomotor tracking) to characterize an unknown transgenic line.

\item {} 
\sphinxAtStartPar
\sphinxstylestrong{{\hyperref[\detokenize{Drosophila/Day2::doc}]{\sphinxcrossref{\DUrole{std,std-doc}{Day 2: The Case of the Missing Methods}}}}} — design your own optogenetic experiment to figure out what the genetically modified neurons control.

\end{itemize}


\subsection{In these labs you will:}
\label{\detokenize{Drosophila/Introduction:in-these-labs-you-will}}\begin{itemize}
\item {} 
\sphinxAtStartPar
Operationalize and quantify the behavior of a fruit fly

\item {} 
\sphinxAtStartPar
Use a computer vision program to track locomotor behaviors

\item {} 
\sphinxAtStartPar
Identify different types of transgenic flies used for neuroscience research

\item {} 
\sphinxAtStartPar
Design an optogenetic experiment to interrogate neural circuit function

\end{itemize}


\subsection{Relevant references}
\label{\detokenize{Drosophila/Introduction:relevant-references}}
\sphinxAtStartPar
Al\sphinxhyphen{}Anzi, B., Tracey, W. D., \& Benzer, S. (2006). Response of \sphinxstyleemphasis{Drosophila} to wasabi is mediated by \sphinxstyleemphasis{painless}, the fly homolog of mammalian TRPA1/ANKTM1. \sphinxstyleemphasis{Current Biology} 16(10): 1034–1040.

\sphinxAtStartPar
Chen, Y.\sphinxhyphen{}C. D., \& Dahanukar, A. (2020). Recent advances in the genetic basis of taste detection in \sphinxstyleemphasis{Drosophila}. \sphinxstyleemphasis{Cellular and Molecular Life Sciences} 77(6): 1087–1101.

\sphinxAtStartPar
Devineni, A. V., \& Heberlein, U. (2009). Preferential ethanol consumption in \sphinxstyleemphasis{Drosophila} models features of addiction. \sphinxstyleemphasis{Current Biology} 19(24): 2126–2132.

\sphinxAtStartPar
Mandel, S. J., Shoaf, M. L., Bostwick, J. R., Combs, D. J., Trudeau, V. L., \& Johnson, E. C. (2018). Behavioral aggression is regulated by \sphinxstyleemphasis{Drosophila} CCAP signaling. \sphinxstyleemphasis{Frontiers in Neural Circuits} 12: 45.

\sphinxAtStartPar
Nakamura, M., Baldwin, D., Hannaford, S., Palka, J., \& Montell, C. (2002). Defective proboscis extension response (DPR), a member of the Ig superfamily required for the gustatory response to salt. \sphinxstyleemphasis{Journal of Neuroscience} 22(9): 3463–3472.

\sphinxAtStartPar
Trisal, S., VijayRaghavan, K., \& Ramaswami, M. (2023). The proboscis extension reflex assay for evaluating taste responses in \sphinxstyleemphasis{Drosophila}. \sphinxstyleemphasis{Bio\sphinxhyphen{}protocol} 13(20).

\sphinxAtStartPar
Ugur, B., Chen, K., \& Bellen, H. J. (2016). \sphinxstyleemphasis{Drosophila} tools and assays for the study of human diseases. \sphinxstyleemphasis{Disease Models \& Mechanisms} 9(3): 235–244.

\sphinxstepscope


\section{Day 1: The Case of the Mislabeled Fly Vials}
\label{\detokenize{Drosophila/Day1:day-1-the-case-of-the-mislabeled-fly-vials}}\label{\detokenize{Drosophila/Day1::doc}}
\sphinxAtStartPar
Something terrible has happened in the lab. The flies you once knew, that were once perfectly organized and labeled, are now completely jumbled. You’re pretty sure the flies are mutated, in one way or another. It’s up to you to figure out what type of flies these are, based on three behavioral assays: a \sphinxstylestrong{proboscis extension reflex (PER) assay, an ethological survey, and a locomotion test}.

\sphinxAtStartPar
Everyone will receive one batch of wildtype flies, and one batch of an experimental group. Ultimately, \sphinxstylestrong{you will present your findings to the class} (see Canvas for details) so that we can determine the phenotype of our mutant flies.


\begin{savenotes}\sphinxattablestart
\sphinxthistablewithglobalstyle
\centering
\begin{tabulary}{\linewidth}[t]{TT}
\sphinxtoprule
\sphinxstyletheadfamily 
\sphinxAtStartPar
\sphinxstylestrong{Supplies}
&\sphinxstyletheadfamily 
\sphinxAtStartPar
\sphinxstylestrong{Solutions}
\\
\sphinxmidrule
\sphinxtableatstartofbodyhook
\sphinxAtStartPar
Two vials of flies (\textasciitilde{}10 WT, \textasciitilde{}10 Mutant/Knockout)
&
\sphinxAtStartPar
DI water
\\
\sphinxhline
\sphinxAtStartPar
Dissecting microscope
&
\sphinxAtStartPar
100 mM sucrose solution
\\
\sphinxhline
\sphinxAtStartPar
Paintbrush
&
\sphinxAtStartPar
40\% EtOH in 100 mM sucrose solution
\\
\sphinxhline
\sphinxAtStartPar
Large petri dish (lid or base)
&
\sphinxAtStartPar

\\
\sphinxhline
\sphinxAtStartPar
Clear cyanoacrylate glue
&
\sphinxAtStartPar

\\
\sphinxhline
\sphinxAtStartPar
4 mL transfer pipette
&
\sphinxAtStartPar

\\
\sphinxhline
\sphinxAtStartPar
Coverslips (6 per fly group)
&
\sphinxAtStartPar

\\
\sphinxhline
\sphinxAtStartPar
Paper towels \& KimWipes
&
\sphinxAtStartPar

\\
\sphinxhline
\sphinxAtStartPar
Plastic “fly recovery” box
&
\sphinxAtStartPar

\\
\sphinxhline
\sphinxAtStartPar
Petri dish with sylgard bottom
&
\sphinxAtStartPar

\\
\sphinxhline
\sphinxAtStartPar
Magnetic stand \& camera
&
\sphinxAtStartPar

\\
\sphinxbottomrule
\end{tabulary}
\sphinxtableafterendhook\par
\sphinxattableend\end{savenotes}


\bigskip\hrule\bigskip



\subsection{Test I. Proboscis Extension Reflex (PER) Assay}
\label{\detokenize{Drosophila/Day1:test-i-proboscis-extension-reflex-per-assay}}
\sphinxAtStartPar
\sphinxincludegraphics{{per_overview}.png}

\sphinxAtStartPar
The \sphinxstylestrong{proboscis extension} is the reflexive extension of the mouthpart (proboscis) in response to taste stimuli. It is a fundamental feeding behavior that can be stimulated by exposing the antennae, labellar hairs, or tarsi (last segments of the leg) to an appetitive taste stimulus.

\sphinxAtStartPar
While a solution with a low concentration of ethanol is appetitive to fruit flies (serving as an olfactory cue that plant\sphinxhyphen{}material fermentation has begun), higher concentrations (>10\% EtOH) can be aversive and toxic to fruit flies. In this assay, you’ll compare your wildtype flies to a knockout group across three solutions: sucrose (an appetitive stimulus), DI water (a neutral control), and 40\% EtOH in sucrose (an aversive stimulus).


\subsubsection{Immobilization}
\label{\detokenize{Drosophila/Day1:immobilization}}\begin{enumerate}
\sphinxsetlistlabels{\arabic}{enumi}{enumii}{}{.}%
\item {} 
\sphinxAtStartPar
Place fly tubes on ice for 2–3 minutes until flies are immobile.

\sphinxAtStartPar
\sphinxstylestrong{Note:} Keep a close eye on your flies and do not leave them on ice too long! As soon as they’re not moving, you can transfer them.

\item {} 
\sphinxAtStartPar
Place the bottom of the petri dish upside down on ice. Wipe away condensation before placing flies on the dish.

\item {} 
\sphinxAtStartPar
Using a brush, \sphinxstyleemphasis{very gently} isolate one fly onto the petri dish and \sphinxstyleemphasis{carefully} orient the fly so it is dorsal\sphinxhyphen{}side up (Figure E below). Be gentle — the flies are very fragile.

\item {} 
\sphinxAtStartPar
Prepare a coverslip: using the transfer pipette, place a tiny drop (1 μl, <2 mm in diameter) of glue on a coverslip.

\end{enumerate}

\sphinxAtStartPar
\sphinxincludegraphics{{per_immobilization}.png}
\begin{enumerate}
\sphinxsetlistlabels{\arabic}{enumi}{enumii}{}{.}%
\setcounter{enumi}{4}
\item {} 
\sphinxAtStartPar
Gently lower the coverslip onto the dorsal side of the fly, until just submerging the back and wings, immobilizing the fly in the glue (Figures F–G above).

\sphinxAtStartPar
\sphinxstylestrong{Note:} Only the back should be submerged in the glue. The head, proboscis, and at least one leg should be exposed — if any one of those parts is submerged, try again with another fly. If you’ve only submerged the wings, try flipping the coverslip over and using forceps to gently lower the back into the glue.

\item {} 
\sphinxAtStartPar
Place one damp paper towel in the plastic container on your lab bench. Ensure the box is dry. Place the coverslip in the box and close the container to regulate humidity.

\item {} 
\sphinxAtStartPar
Repeat with the remaining flies. You should prepare \sphinxstylestrong{5 wildtype} and \sphinxstylestrong{5 knockout} flies.

\item {} 
\sphinxAtStartPar
Allow the flies to recover for at least \sphinxstylestrong{1.5 hours} before testing. Complete the other behavioral tests during the recovery period.

\sphinxAtStartPar
\sphinxstylestrong{Note:} Omit flies that do not show movement after 1.5 hours of recovery.

\end{enumerate}


\subsubsection{Running the PER assay}
\label{\detokenize{Drosophila/Day1:running-the-per-assay}}\begin{enumerate}
\sphinxsetlistlabels{\arabic}{enumi}{enumii}{}{.}%
\setcounter{enumi}{8}
\item {} 
\sphinxAtStartPar
Choose one of your responsive, immobilized flies with accessible head and legs, and place it under the dissecting microscope. Set up your lighting so you can see the legs and proboscis.

\item {} 
\sphinxAtStartPar
Dip the corner of a KimWipe in the sucrose solution and gently touch the tarsi with the KimWipe for 3 seconds.

\sphinxAtStartPar
\sphinxincludegraphics{{per_scoring}.png}

\sphinxAtStartPar
a. Score the response in \sphinxstylestrong{Data Table 1}. Assign a \sphinxstylestrong{1} for proboscis extension (right image above) and a \sphinxstylestrong{0} for no extension (left image above).

\sphinxAtStartPar
b. Wait 3 seconds or until the proboscis retracts.

\sphinxAtStartPar
c. Repeat for \sphinxstylestrong{10 trials}.

\item {} 
\sphinxAtStartPar
Repeat using \sphinxstylestrong{DI water}.

\item {} 
\sphinxAtStartPar
Repeat using \sphinxstylestrong{40\% EtOH in sucrose}.

\end{enumerate}


\subsubsection{Data Table 1: PER Responses}
\label{\detokenize{Drosophila/Day1:data-table-1-per-responses}}

\begin{savenotes}\sphinxattablestart
\sphinxthistablewithglobalstyle
\centering
\begin{tabulary}{\linewidth}[t]{TTTTT}
\sphinxtoprule
\sphinxstyletheadfamily 
\sphinxAtStartPar
Genotype
&\sphinxstyletheadfamily 
\sphinxAtStartPar
Trial
&\sphinxstyletheadfamily 
\sphinxAtStartPar
Sucrose
&\sphinxstyletheadfamily 
\sphinxAtStartPar
DI Water
&\sphinxstyletheadfamily 
\sphinxAtStartPar
40\% EtOH in Sucrose
\\
\sphinxmidrule
\sphinxtableatstartofbodyhook
\sphinxAtStartPar
Wildtype
&
\sphinxAtStartPar
1
&
\sphinxAtStartPar

&
\sphinxAtStartPar

&
\sphinxAtStartPar

\\
\sphinxhline
\sphinxAtStartPar
Wildtype
&
\sphinxAtStartPar
2
&
\sphinxAtStartPar

&
\sphinxAtStartPar

&
\sphinxAtStartPar

\\
\sphinxhline
\sphinxAtStartPar
Wildtype
&
\sphinxAtStartPar
3
&
\sphinxAtStartPar

&
\sphinxAtStartPar

&
\sphinxAtStartPar

\\
\sphinxhline
\sphinxAtStartPar
Wildtype
&
\sphinxAtStartPar
4
&
\sphinxAtStartPar

&
\sphinxAtStartPar

&
\sphinxAtStartPar

\\
\sphinxhline
\sphinxAtStartPar
Wildtype
&
\sphinxAtStartPar
5
&
\sphinxAtStartPar

&
\sphinxAtStartPar

&
\sphinxAtStartPar

\\
\sphinxhline
\sphinxAtStartPar
Wildtype
&
\sphinxAtStartPar
6
&
\sphinxAtStartPar

&
\sphinxAtStartPar

&
\sphinxAtStartPar

\\
\sphinxhline
\sphinxAtStartPar
Wildtype
&
\sphinxAtStartPar
7
&
\sphinxAtStartPar

&
\sphinxAtStartPar

&
\sphinxAtStartPar

\\
\sphinxhline
\sphinxAtStartPar
Wildtype
&
\sphinxAtStartPar
8
&
\sphinxAtStartPar

&
\sphinxAtStartPar

&
\sphinxAtStartPar

\\
\sphinxhline
\sphinxAtStartPar
Wildtype
&
\sphinxAtStartPar
9
&
\sphinxAtStartPar

&
\sphinxAtStartPar

&
\sphinxAtStartPar

\\
\sphinxhline
\sphinxAtStartPar
Wildtype
&
\sphinxAtStartPar
10
&
\sphinxAtStartPar

&
\sphinxAtStartPar

&
\sphinxAtStartPar

\\
\sphinxhline
\sphinxAtStartPar
Knockout
&
\sphinxAtStartPar
1
&
\sphinxAtStartPar

&
\sphinxAtStartPar

&
\sphinxAtStartPar

\\
\sphinxhline
\sphinxAtStartPar
Knockout
&
\sphinxAtStartPar
2
&
\sphinxAtStartPar

&
\sphinxAtStartPar

&
\sphinxAtStartPar

\\
\sphinxhline
\sphinxAtStartPar
Knockout
&
\sphinxAtStartPar
3
&
\sphinxAtStartPar

&
\sphinxAtStartPar

&
\sphinxAtStartPar

\\
\sphinxhline
\sphinxAtStartPar
Knockout
&
\sphinxAtStartPar
4
&
\sphinxAtStartPar

&
\sphinxAtStartPar

&
\sphinxAtStartPar

\\
\sphinxhline
\sphinxAtStartPar
Knockout
&
\sphinxAtStartPar
5
&
\sphinxAtStartPar

&
\sphinxAtStartPar

&
\sphinxAtStartPar

\\
\sphinxhline
\sphinxAtStartPar
Knockout
&
\sphinxAtStartPar
6
&
\sphinxAtStartPar

&
\sphinxAtStartPar

&
\sphinxAtStartPar

\\
\sphinxhline
\sphinxAtStartPar
Knockout
&
\sphinxAtStartPar
7
&
\sphinxAtStartPar

&
\sphinxAtStartPar

&
\sphinxAtStartPar

\\
\sphinxhline
\sphinxAtStartPar
Knockout
&
\sphinxAtStartPar
8
&
\sphinxAtStartPar

&
\sphinxAtStartPar

&
\sphinxAtStartPar

\\
\sphinxhline
\sphinxAtStartPar
Knockout
&
\sphinxAtStartPar
9
&
\sphinxAtStartPar

&
\sphinxAtStartPar

&
\sphinxAtStartPar

\\
\sphinxhline
\sphinxAtStartPar
Knockout
&
\sphinxAtStartPar
10
&
\sphinxAtStartPar

&
\sphinxAtStartPar

&
\sphinxAtStartPar

\\
\sphinxbottomrule
\end{tabulary}
\sphinxtableafterendhook\par
\sphinxattableend\end{savenotes}

\sphinxAtStartPar
\sphinxstylestrong{Reflection Question:} Identify the positive and negative controls in this experiment. With these controls in mind, are your findings as expected?


\bigskip\hrule\bigskip



\subsection{Test II. Ethological approach to behavior}
\label{\detokenize{Drosophila/Day1:test-ii-ethological-approach-to-behavior}}
\sphinxAtStartPar
It’s possible your flies do something different that’s \sphinxstyleemphasis{not} captured by a standard behavioral assay. The best way to figure this out is to watch the flies, choose several behaviors to quantify, and see if those behaviors differ between your experimental and wildtype flies.


\subsubsection{Choosing behaviors to quantify}
\label{\detokenize{Drosophila/Day1:choosing-behaviors-to-quantify}}
\sphinxAtStartPar
Your next goal for today will be to quantify \sphinxstylestrong{three} different behaviors that your flies perform. It is up to you to define and operationalize these behaviors.
\begin{enumerate}
\sphinxsetlistlabels{\arabic}{enumi}{enumii}{}{.}%
\item {} 
\sphinxAtStartPar
Freeze both groups of flies for 2–3 minutes until they stop moving.

\item {} 
\sphinxAtStartPar
Isolate three flies of each type in separate petri dishes.

\sphinxAtStartPar
\sphinxstylestrong{Note:} Make sure the flies are completely awake and moving before moving on to the next step. This could take about 10 minutes.

\item {} 
\sphinxAtStartPar
Take a few minutes to observe both batches of flies.

\item {} 
\sphinxAtStartPar
Choose three different behaviors that the flies do, and record them below. You also need to provide a definition for each behavior.

\end{enumerate}

\sphinxAtStartPar
\sphinxstylestrong{Table 1.}


\begin{savenotes}\sphinxattablestart
\sphinxthistablewithglobalstyle
\centering
\begin{tabulary}{\linewidth}[t]{TTTT}
\sphinxtoprule
\sphinxstyletheadfamily 
\sphinxAtStartPar
Behavior name
&\sphinxstyletheadfamily 
\sphinxAtStartPar
Definition
&\sphinxstyletheadfamily 
\sphinxAtStartPar
\# Instances in WT
&\sphinxstyletheadfamily 
\sphinxAtStartPar
\# Instances in Mutant
\\
\sphinxmidrule
\sphinxtableatstartofbodyhook
\sphinxAtStartPar

&
\sphinxAtStartPar

&
\sphinxAtStartPar

&
\sphinxAtStartPar

\\
\sphinxhline
\sphinxAtStartPar

&
\sphinxAtStartPar

&
\sphinxAtStartPar

&
\sphinxAtStartPar

\\
\sphinxhline
\sphinxAtStartPar

&
\sphinxAtStartPar

&
\sphinxAtStartPar

&
\sphinxAtStartPar

\\
\sphinxbottomrule
\end{tabulary}
\sphinxtableafterendhook\par
\sphinxattableend\end{savenotes}
\begin{enumerate}
\sphinxsetlistlabels{\arabic}{enumi}{enumii}{}{.}%
\setcounter{enumi}{4}
\item {} 
\sphinxAtStartPar
For 5 minutes, each person in your group should observe one behavior and record the number of instances those behaviors happen in the WT flies. Record the number of instances in the table above.

\item {} 
\sphinxAtStartPar
Do the same for the experimental group.

\end{enumerate}


\subsection{Test III. Recording locomotor behavior}
\label{\detokenize{Drosophila/Day1:test-iii-recording-locomotor-behavior}}
\sphinxAtStartPar
For the final test, you’ll record the movement of your two different types of flies in a petri dish and use an automated tracking program to track the velocity and heading of your flies. Your mutants might have altered locomotor behavior — this will help you test that.


\subsubsection{Set up your camera}
\label{\detokenize{Drosophila/Day1:set-up-your-camera}}
\sphinxAtStartPar
\sphinxincludegraphics{{fly_camera_setup}.png}

\sphinxAtStartPar
First, we need to arrange the camera and petri dish so that you can get everything in frame, in focus, and well lit.
\begin{enumerate}
\sphinxsetlistlabels{\arabic}{enumi}{enumii}{}{.}%
\item {} 
\sphinxAtStartPar
Ensure that your camera is connected via the USB cord to the computer. The camera should be plugged into a USB 3.0 slot (the blue ones).

\item {} 
\sphinxAtStartPar
Set up your camera and petri dish as in the figure above. You want your camera to be about 6 inches from the petri dish.

\item {} 
\sphinxAtStartPar
On the computer, open \sphinxstylestrong{Point Grey FlyCap2}.

\item {} 
\sphinxAtStartPar
Make sure that it recognizes a “Flea3” camera, and press OK.

\item {} 
\sphinxAtStartPar
Adjust the \sphinxstylestrong{focal length} knob on your camera (closest to end, ranges from infinity to 0.1) in order to get the image in focus.

\item {} 
\sphinxAtStartPar
Adjust the image on your screen using the toolbar at the top:

\sphinxAtStartPar
\sphinxincludegraphics{{flycap_toolbar}.png}

\item {} 
\sphinxAtStartPar
Isolate one fly in your petri dish.

\item {} 
\sphinxAtStartPar
Click on the button to open up the Camera settings (e.g., exposure).

\sphinxAtStartPar
a. Uncheck \sphinxstyleemphasis{auto} for brightness, exposure, and shutter.

\sphinxAtStartPar
b. Increase the exposure almost to the point where your image is overexposed. You want the flies to stand out very clearly against the white background, and the edges of the petri dish to be invisible:

\sphinxAtStartPar
\sphinxincludegraphics{{fly_arena}.png}

\item {} 
\sphinxAtStartPar
Ensure that you can easily see your fly throughout all areas of the arena.

\sphinxAtStartPar
\sphinxstylestrong{Note:} Optimizing the way your fly recording looks will make the tracking \sphinxstyleemphasis{much} easier. Make sure there are no glares from lights in the room, no dark spots on the petri dish, and clean off any fingerprints on the glass. \sphinxstylestrong{Every light and dark spot matters, and anything can throw off the tracking.}

\end{enumerate}


\subsubsection{Set up your flies \& record}
\label{\detokenize{Drosophila/Day1:set-up-your-flies-record}}
\sphinxAtStartPar
\sphinxincludegraphics{{fly_dish_setup}.png}
\begin{enumerate}
\sphinxsetlistlabels{\arabic}{enumi}{enumii}{}{.}%
\item {} 
\sphinxAtStartPar
When you’re ready to record, go to \sphinxstylestrong{File > Capture Image or Video}.

\item {} 
\sphinxAtStartPar
A dialogue will pop up. Set up a 60 second recording with the following settings:
\begin{itemize}
\item {} 
\sphinxAtStartPar
\sphinxstylestrong{Filename:} browse to create your own (make sure you label it as WT or Knockout). \sphinxstylestrong{SAVE YOUR FILE LOCALLY}, not on your username server. \sphinxstylestrong{Your video file will not save correctly if you do not do this.}

\item {} 
\sphinxAtStartPar
🔘 Capture for 60000 ms

\item {} 
\sphinxAtStartPar
🔘 Buffered

\item {} 
\sphinxAtStartPar
(Click on \sphinxstylestrong{Videos} tab)

\item {} 
\sphinxAtStartPar
\sphinxstylestrong{Video Recording Type:} M\sphinxhyphen{}JPEG

\item {} 
\sphinxAtStartPar
\sphinxstylestrong{Frame Rate:} 15.0

\item {} 
\sphinxAtStartPar
\sphinxstylestrong{AVI Split Size:} 0

\item {} 
\sphinxAtStartPar
\sphinxstylestrong{JPEG Compression Quality:} 75

\end{itemize}

\item {} 
\sphinxAtStartPar
Press \sphinxstylestrong{Start Recording}, and wait while the program records the video.

\sphinxAtStartPar
\sphinxstylestrong{Note:} This window will not close, but the message at the bottom will change.

\item {} 
\sphinxAtStartPar
When it is done, \sphinxstylestrong{check that the video recorded correctly by finding the file and opening it.}

\sphinxAtStartPar
\sphinxstylestrong{Note:} If it didn’t record correctly, close FlyCap, unplug the camera, and retry.

\item {} 
\sphinxAtStartPar
For the first fly you record, \sphinxstylestrong{move on to the next step to ensure that your video works well for tracking. If the tracking doesn’t work well, you’ll need to adjust your lighting and/or image acquisition.}

\item {} 
\sphinxAtStartPar
When you’re confident the tracking works well with your videos, obtain videos for three flies from each group (you’ve already done one).

\item {} 
\sphinxAtStartPar
When you’re done, save your videos in a clouded location for safekeeping.

\item {} 
\sphinxAtStartPar
Close FlyCap.

\end{enumerate}


\subsubsection{Quantify your locomotor behavior}
\label{\detokenize{Drosophila/Day1:quantify-your-locomotor-behavior}}
\sphinxAtStartPar
Now that you have videos of your specimen, you need to quantify the locomotor behavior. We’ll use a MATLAB program to do so.


\paragraph{Download the code we need}
\label{\detokenize{Drosophila/Day1:download-the-code-we-need}}\begin{enumerate}
\sphinxsetlistlabels{\arabic}{enumi}{enumii}{}{.}%
\item {} 
\sphinxAtStartPar
Download the code for \sphinxstylestrong{fly\_tracker} from GitHub: \sphinxurl{https://github.com/ajuavinett/fly\_tracker}

\item {} 
\sphinxAtStartPar
Extract all of the contents of the zipped folder into your \sphinxstylestrong{local} Documents folder (\sphinxstyleemphasis{not} a folder on your username).

\item {} 
\sphinxAtStartPar
Open MATLAB.

\item {} 
\sphinxAtStartPar
On the top bar of MATLAB, find the “Set Path” button:

\sphinxAtStartPar
\sphinxincludegraphics{{matlab_setpath}.png}

\item {} 
\sphinxAtStartPar
In this window, click “Add with Subfolders.”

\item {} 
\sphinxAtStartPar
Navigate to the folder where you extracted fly\_tracker. Hit \sphinxstylestrong{Save}.

\sphinxAtStartPar
\sphinxstylestrong{Note:} You might get an error about pathdef.m. Select \sphinxstylestrong{Yes} and save this file in your \sphinxstylestrong{local} Documents folder.

\end{enumerate}


\paragraph{Test the tracking for a sample video}
\label{\detokenize{Drosophila/Day1:test-the-tracking-for-a-sample-video}}\begin{enumerate}
\sphinxsetlistlabels{\arabic}{enumi}{enumii}{}{.}%
\setcounter{enumi}{6}
\item {} 
\sphinxAtStartPar
In the MATLAB window, type:

\begin{sphinxVerbatim}[commandchars=\\\{\}]
\PYG{p}{[}\PYG{n}{mean\PYGZus{}velocity}\PYG{+w}{ }\PYG{n}{SEM\PYGZus{}velocity}\PYG{p}{]}\PYG{+w}{ }\PYG{p}{=}\PYG{+w}{ }\PYG{n}{BIPN145\PYGZus{}flytrack}
\end{sphinxVerbatim}

\sphinxAtStartPar
and press enter.

\item {} 
\sphinxAtStartPar
Navigate to where your movie is and select it.

\item {} 
\sphinxAtStartPar
In the next window, draw a rectangle (click and drag) around your entire petri dish.

\item {} 
\sphinxAtStartPar
Double\sphinxhyphen{}click inside the rectangle to initiate tracking.

\item {} 
\sphinxAtStartPar
Wait for the tracking to finish. A window will pop up that shows the trajectory of your fly and its velocity across time.

\sphinxAtStartPar
\sphinxstylestrong{Note:} If you’d like to edit these figures in any way, you can go to \sphinxstylestrong{Tools > Edit Plot.}

\end{enumerate}


\paragraph{Run your analysis for a group of flies}
\label{\detokenize{Drosophila/Day1:run-your-analysis-for-a-group-of-flies}}
\sphinxAtStartPar
Once you’ve tested your tracking \& recorded all of your movies, you can run this script for \sphinxstylestrong{several movies} at a time (e.g., one group of flies). You should record videos for 3 WT flies and 3 mutant flies. Do this by re\sphinxhyphen{}doing steps 7–11 above, but \sphinxstylestrong{select all of the videos for one group of flies} in the window where you are asked to pick a movie. After you’ve done the analysis for your wildtype flies, run the \sphinxcode{\sphinxupquote{BIPN145\_flytrack}} script again with your mutant fly videos.

\sphinxAtStartPar
The script outputs a variable called \sphinxcode{\sphinxupquote{mean\_velocity}}. This will show the mean velocity (in mm/s) for \sphinxstyleemphasis{each} of your videos. \sphinxcode{\sphinxupquote{SEM\_velocity}} is the standard error of the mean of the velocity for these videos.

\sphinxAtStartPar
To generate a mean value for the video(s) you just tracked, type:

\begin{sphinxVerbatim}[commandchars=\\\{\}]
\PYG{n+nb}{mean}\PYG{p}{(}\PYG{n}{mean\PYGZus{}velocity}\PYG{p}{)}
\end{sphinxVerbatim}

\sphinxAtStartPar
and press enter.

\sphinxAtStartPar
Record these values in the table below (you’ll ultimately want to plot these).

\sphinxAtStartPar
\sphinxstylestrong{Table 2.} Comparison of mean velocities for control and mutant flies.


\begin{savenotes}\sphinxattablestart
\sphinxthistablewithglobalstyle
\centering
\begin{tabulary}{\linewidth}[t]{TTT}
\sphinxtoprule

\sphinxAtStartPar

&\sphinxstyletheadfamily 
\sphinxAtStartPar
\sphinxstylestrong{Control}
&\sphinxstyletheadfamily 
\sphinxAtStartPar
\sphinxstylestrong{Mutant}
\\
\sphinxmidrule
\sphinxtableatstartofbodyhook
\sphinxAtStartPar
Mean
&
\sphinxAtStartPar

&
\sphinxAtStartPar

\\
\sphinxhline
\sphinxAtStartPar
SEM
&
\sphinxAtStartPar

&
\sphinxAtStartPar

\\
\sphinxbottomrule
\end{tabulary}
\sphinxtableafterendhook\par
\sphinxattableend\end{savenotes}


\subsection{Troubleshooting}
\label{\detokenize{Drosophila/Day1:troubleshooting}}

\begin{savenotes}\sphinxattablestart
\sphinxthistablewithglobalstyle
\centering
\begin{tabulary}{\linewidth}[t]{TTT}
\sphinxtoprule
\sphinxstyletheadfamily 
\sphinxAtStartPar
\sphinxstylestrong{Observation}
&\sphinxstyletheadfamily 
\sphinxAtStartPar
\sphinxstylestrong{Likely issue(s)}
&\sphinxstyletheadfamily 
\sphinxAtStartPar
\sphinxstylestrong{Possible solution}
\\
\sphinxmidrule
\sphinxtableatstartofbodyhook
\sphinxAtStartPar
The camera output is black
&
\sphinxAtStartPar
Lens cap is on
&
\sphinxAtStartPar
Remove the lens cap
\\
\sphinxhline
\sphinxAtStartPar
Camera not recognized
&
\sphinxAtStartPar
Camera not plugged in
&
\sphinxAtStartPar
Plug in the camera
\\
\sphinxhline
\sphinxAtStartPar
Video file is not recognized in MATLAB or the movie is the wrong length/size
&
\sphinxAtStartPar
You wrote the video file to the server
&
\sphinxAtStartPar
You need to re\sphinxhyphen{}record your video but with the file saved \sphinxstylestrong{directly to the computer locally}
\\
\sphinxhline
\sphinxAtStartPar
Fly does not extend proboscis to sucrose
&
\sphinxAtStartPar
Fly is dehydrated or improperly mounted
&
\sphinxAtStartPar
Check humidity in recovery box; verify proboscis and at least one leg are not stuck in glue
\\
\sphinxhline
\sphinxAtStartPar
All flies look the same in tracking video
&
\sphinxAtStartPar
Lighting/exposure issues
&
\sphinxAtStartPar
Increase exposure so flies stand out as dark spots on a fully white background
\\
\sphinxbottomrule
\end{tabulary}
\sphinxtableafterendhook\par
\sphinxattableend\end{savenotes}

\sphinxstepscope


\section{Day 2: The Case of the Missing Methods}
\label{\detokenize{Drosophila/Day2:day-2-the-case-of-the-missing-methods}}\label{\detokenize{Drosophila/Day2::doc}}
\sphinxAtStartPar
Just when you thought everything was back to normal in the lab, a remarkably inconveniently placed coffee spill completely drenches several pages of the notes given to you from your mentor.

\sphinxAtStartPar
Unfortunately, those pages detailed all of the methods that you needed in order to perform your experiment. You’ve managed to save the supplies list, at least:

\sphinxAtStartPar
\sphinxincludegraphics{{coffee_supplies}.png}

\sphinxAtStartPar
Hope is not lost. You’ve got some supplies, and you suspect that this study had something to do with the role of the genetically\sphinxhyphen{}modified neurons. Certain neurons have optogenetic channels, but what do those neurons actually control?

\sphinxAtStartPar
It might take a little bit of research to refresh your memories about those transgenic flies, but that’s easy enough. And you’ve done at least some fly behavior before — you’re pretty much the fly whisperer.

\sphinxAtStartPar
You can re\sphinxhyphen{}write the protocol and get some interesting results, can’t you?


\subsection{Tips for designing your experiment}
\label{\detokenize{Drosophila/Day2:tips-for-designing-your-experiment}}
\sphinxAtStartPar
You can control the intensity of the light source through LabChart by following these steps:
\begin{enumerate}
\sphinxsetlistlabels{\arabic}{enumi}{enumii}{}{.}%
\item {} 
\sphinxAtStartPar
Plug the BNC input on the light into the \sphinxstylestrong{OUTPUT} of the PowerLab using a BNC cable.

\item {} 
\sphinxAtStartPar
Program the stimulator (\sphinxstylestrong{Setup > Stimulator}) to send different pulses with heights from 2V – 10V. You should notice that the intensity of the light changes with more voltage.

\sphinxAtStartPar
\sphinxstyleemphasis{\sphinxstylestrong{Do not}} stimulate your flies at higher than 1 Hz frequency or with stimulus pulses longer than 0.5 seconds.

\item {} 
\sphinxAtStartPar
Measure the intensity of the light using a phone app luxmeter (we recommend \sphinxstyleemphasis{Lux Light Meter Free}). You should measure the light intensity \sphinxstylestrong{from the same distance} as you are using it in the experiment.

\end{enumerate}

\sphinxAtStartPar
Don’t forget about control groups! It may help to compare these flies to each other, as well as to wildtype flies. We have wildtype flies if you need them.

\sphinxAtStartPar
You will need to put the red light \sphinxstylestrong{within several inches} of your flies in order to see an effect of the stimulation.

\sphinxstepscope


\section{Fly Tracker (Local)}
\label{\detokenize{FlyTracking/Script:fly-tracker-local}}\label{\detokenize{FlyTracking/Script::doc}}
\sphinxAtStartPar
In this lab, you’ll use a Python script to track the position and velocity of a fruit fly from video recordings. The script will detect the fly against a light background, track its movement frame\sphinxhyphen{}by\sphinxhyphen{}frame, and output position and velocity data as CSV files, along with plots.

\sphinxAtStartPar
\sphinxstylestrong{At the end of this notebook, you’ll be able to:}
\begin{itemize}
\item {} 
\sphinxAtStartPar
Run the fly tracker script on your own video recordings

\item {} 
\sphinxAtStartPar
Load and visualize tracking and velocity data in Python

\item {} 
\sphinxAtStartPar
Customize your plots for your lab report

\end{itemize}




\subsection{Part 1: Running the Fly Tracker Script}
\label{\detokenize{FlyTracking/Script:part-1-running-the-fly-tracker-script}}
\sphinxAtStartPar
The fly tracker is a Python script that you’ll run from the command line. Follow the steps below to get set up.


\subsubsection{Step 1: Open Anaconda Prompt}
\label{\detokenize{FlyTracking/Script:step-1-open-anaconda-prompt}}\begin{enumerate}
\sphinxsetlistlabels{\arabic}{enumi}{enumii}{}{.}%
\item {} 
\sphinxAtStartPar
Click the \sphinxstylestrong{Windows search bar} (bottom\sphinxhyphen{}left of your screen) and type \sphinxstylestrong{Anaconda}.

\item {} 
\sphinxAtStartPar
Click \sphinxstylestrong{Anaconda Prompt} to open it. You should see a black terminal window.

\end{enumerate}

\begin{sphinxadmonition}{note}{Note:}
\sphinxAtStartPar
This is intended to be run on the BIPN 145 Windows lab computers in York 1310. If you are running this on your own Mac computer, you can search for Python and run it that way. If you’re running it on your own Windows computer, you’ll need to download Python or Anaconda.
\end{sphinxadmonition}


\subsubsection{Step 2: Download the Scripts}
\label{\detokenize{FlyTracking/Script:step-2-download-the-scripts}}
\sphinxAtStartPar
In the Anaconda Prompt, first use the following to move into \sphinxstyleemphasis{your} folder. Replace USER with the UCSD username (not including @ucsd.edu) for the person that is logged into the computer.

\begin{sphinxadmonition}{tip}{Tip:}
\sphinxAtStartPar
\sphinxstylestrong{Copying and pasting into the Anaconda Prompt:} You can copy commands from this page by clicking the copy button (📋) at the top\sphinxhyphen{}right of each code block and using Ctrl+V to paste into the command line.
\end{sphinxadmonition}

\begin{sphinxVerbatim}[commandchars=\\\{\}]
\PYG{n}{cd} \PYG{n}{C}\PYG{p}{:}\PYGZbs{}\PYG{n}{Users}\PYGZbs{}\PYG{n}{USER}
\end{sphinxVerbatim}

\sphinxAtStartPar
Then, run the following commands to download the setup and fly tracker scripts:

\begin{sphinxVerbatim}[commandchars=\\\{\}]
\PYG{n}{curl} \PYG{o}{\PYGZhy{}}\PYG{n}{o} \PYG{n}{setup}\PYG{o}{.}\PYG{n}{py} \PYG{n}{https}\PYG{p}{:}\PYG{o}{/}\PYG{o}{/}\PYG{n}{raw}\PYG{o}{.}\PYG{n}{githubusercontent}\PYG{o}{.}\PYG{n}{com}\PYG{o}{/}\PYG{n}{ajuavinett}\PYG{o}{/}\PYG{n}{fly\PYGZus{}tracker}\PYG{o}{/}\PYG{n}{master}\PYG{o}{/}\PYG{n}{setup}\PYG{o}{.}\PYG{n}{py}
\end{sphinxVerbatim}

\begin{sphinxVerbatim}[commandchars=\\\{\}]
\PYG{n}{curl} \PYG{o}{\PYGZhy{}}\PYG{n}{o} \PYG{n}{BIPN145\PYGZus{}flytrack}\PYG{o}{.}\PYG{n}{py} \PYG{n}{https}\PYG{p}{:}\PYG{o}{/}\PYG{o}{/}\PYG{n}{raw}\PYG{o}{.}\PYG{n}{githubusercontent}\PYG{o}{.}\PYG{n}{com}\PYG{o}{/}\PYG{n}{ajuavinett}\PYG{o}{/}\PYG{n}{fly\PYGZus{}tracker}\PYG{o}{/}\PYG{n}{master}\PYG{o}{/}\PYG{n}{BIPN145\PYGZus{}flytrack}\PYG{o}{.}\PYG{n}{py}
\end{sphinxVerbatim}

\sphinxAtStartPar
If you already have older versions of these scripts, this will overwrite them with the latest versions. You can also download the scripts manually from the \sphinxhref{https://github.com/ajuavinett/fly\_tracker}{fly\_tracker GitHub page}.


\subsubsection{Step 3: Install Required Packages}
\label{\detokenize{FlyTracking/Script:step-3-install-required-packages}}
\sphinxAtStartPar
Run the setup script to install the required Python packages. You only need to do this \sphinxstylestrong{once} per computer:

\begin{sphinxVerbatim}[commandchars=\\\{\}]
\PYG{n}{python} \PYG{n}{setup}\PYG{o}{.}\PYG{n}{py}
\end{sphinxVerbatim}

\sphinxAtStartPar
You should see a message confirming that each package is installed. Once you see “Setup complete!”, you’re ready to go.


\subsubsection{Step 4: Know Where Your Video File(s) Are}
\label{\detokenize{FlyTracking/Script:step-4-know-where-your-video-file-s-are}}
\sphinxAtStartPar
Make sure you know where your fly video file(s) (\sphinxcode{\sphinxupquote{.avi}}, \sphinxcode{\sphinxupquote{.mp4}}, etc.) are saved on your computer. When you run the script, a file picker window will pop up and you can navigate to them from there.


\subsubsection{Step 5: Run the Script}
\label{\detokenize{FlyTracking/Script:step-5-run-the-script}}
\sphinxAtStartPar
In the Anaconda Prompt, type:

\begin{sphinxVerbatim}[commandchars=\\\{\}]
\PYG{n}{python} \PYG{n}{BIPN145\PYGZus{}flytrack}\PYG{o}{.}\PYG{n}{py}
\end{sphinxVerbatim}

\sphinxAtStartPar
The script will walk you through everything interactively:
\begin{enumerate}
\sphinxsetlistlabels{\arabic}{enumi}{enumii}{}{.}%
\item {} 
\sphinxAtStartPar
\sphinxstylestrong{Select your video(s)} — A file picker dialog will pop up. Navigate to your video file(s), select them, and click Open.

\item {} 
\sphinxAtStartPar
\sphinxstylestrong{Draw the ROI} — A window will pop up showing the first frame of your video. \sphinxstylestrong{Click and drag} to draw a rectangle around the dish/arena (and \sphinxstyleemphasis{only} around the dish – make it as snug as possible!), then press \sphinxstylestrong{Enter} or \sphinxstylestrong{Space} to confirm. Press \sphinxstylestrong{C} to cancel and redraw.

\item {} 
\sphinxAtStartPar
\sphinxstylestrong{Wait for tracking} — The script will process each frame. You’ll see progress updates (10\%, 20\%, etc.). This will take a few minutes to run per video.

\item {} 
\sphinxAtStartPar
\sphinxstylestrong{View results} — When done, plots will appear will save automatically along with CSV files containing the tracking results. You can either use the plots that are saved or modify them (see below).

\end{enumerate}


\subsubsection{Step 6: Find Your Output Files}
\label{\detokenize{FlyTracking/Script:step-6-find-your-output-files}}
\sphinxAtStartPar
After the script finishes, you’ll find two CSV files for each video in the same folder as your video:


\begin{savenotes}\sphinxattablestart
\sphinxthistablewithglobalstyle
\centering
\begin{tabulary}{\linewidth}[t]{TT}
\sphinxtoprule
\sphinxstyletheadfamily 
\sphinxAtStartPar
File
&\sphinxstyletheadfamily 
\sphinxAtStartPar
Contents
\\
\sphinxmidrule
\sphinxtableatstartofbodyhook
\sphinxAtStartPar
\sphinxcode{\sphinxupquote{yourvideo\_tracking.csv}}
&
\sphinxAtStartPar
Time (s), X position (cm), Y position (cm)
\\
\sphinxhline
\sphinxAtStartPar
\sphinxcode{\sphinxupquote{yourvideo\_velocity.csv}}
&
\sphinxAtStartPar
Time (s), Velocity (mm/s)
\\
\sphinxbottomrule
\end{tabulary}
\sphinxtableafterendhook\par
\sphinxattableend\end{savenotes}

\sphinxAtStartPar
If you need to edit them, you can open these in Excel, Google Sheets, or load them into Python (see below).




\subsection{Part 2: Working with Your Data in Python (Optional)}
\label{\detokenize{FlyTracking/Script:part-2-working-with-your-data-in-python-optional}}
\sphinxAtStartPar
The cells below will help you load your CSV output files and create customized plots in case you need to adjust anything after running the tracking. This is useful if you want to adjust axis labels, colors, or combine data from multiple flies for your lab report.
\begin{quote}

\sphinxAtStartPar
\sphinxstylestrong{Task}: Run the cell below to import the packages we’ll need.
\end{quote}

\begin{sphinxuseclass}{cell}\begin{sphinxVerbatimInput}

\begin{sphinxuseclass}{cell_input}
\begin{sphinxVerbatim}[commandchars=\\\{\}]
\PYG{k+kn}{import}\PYG{+w}{ }\PYG{n+nn}{numpy}\PYG{+w}{ }\PYG{k}{as}\PYG{+w}{ }\PYG{n+nn}{np}
\PYG{k+kn}{import}\PYG{+w}{ }\PYG{n+nn}{matplotlib}\PYG{n+nn}{.}\PYG{n+nn}{pyplot}\PYG{+w}{ }\PYG{k}{as}\PYG{+w}{ }\PYG{n+nn}{plt}
\PYG{k+kn}{from}\PYG{+w}{ }\PYG{n+nn}{matplotlib}\PYG{n+nn}{.}\PYG{n+nn}{collections}\PYG{+w}{ }\PYG{k+kn}{import} \PYG{n}{LineCollection}

\PYG{n+nb}{print}\PYG{p}{(}\PYG{l+s+s1}{\PYGZsq{}}\PYG{l+s+s1}{Packages imported!}\PYG{l+s+s1}{\PYGZsq{}}\PYG{p}{)}
\end{sphinxVerbatim}

\end{sphinxuseclass}\end{sphinxVerbatimInput}

\end{sphinxuseclass}

\subsubsection{Upload your CSV files}
\label{\detokenize{FlyTracking/Script:upload-your-csv-files}}
\sphinxAtStartPar
If you’re running this notebook in \sphinxstylestrong{Google Colab}, run the cell below to upload your tracking and velocity CSV files. If you’re running locally, you can skip this cell and just set the file paths directly.
\begin{quote}

\sphinxAtStartPar
\sphinxstylestrong{Task}: Run the cell below and upload your \sphinxcode{\sphinxupquote{\_tracking.csv}} and \sphinxcode{\sphinxupquote{\_velocity.csv}} files.
\end{quote}

\begin{sphinxuseclass}{cell}\begin{sphinxVerbatimInput}

\begin{sphinxuseclass}{cell_input}
\begin{sphinxVerbatim}[commandchars=\\\{\}]
\PYG{k}{try}\PYG{p}{:}
    \PYG{k+kn}{from}\PYG{+w}{ }\PYG{n+nn}{google}\PYG{n+nn}{.}\PYG{n+nn}{colab}\PYG{+w}{ }\PYG{k+kn}{import} \PYG{n}{files}
    \PYG{n}{uploaded} \PYG{o}{=} \PYG{n}{files}\PYG{o}{.}\PYG{n}{upload}\PYG{p}{(}\PYG{p}{)}
    \PYG{n}{filenames} \PYG{o}{=} \PYG{n+nb}{list}\PYG{p}{(}\PYG{n}{uploaded}\PYG{o}{.}\PYG{n}{keys}\PYG{p}{(}\PYG{p}{)}\PYG{p}{)}
    \PYG{n}{tracking\PYGZus{}file} \PYG{o}{=} \PYG{p}{[}\PYG{n}{f} \PYG{k}{for} \PYG{n}{f} \PYG{o+ow}{in} \PYG{n}{filenames} \PYG{k}{if} \PYG{l+s+s1}{\PYGZsq{}}\PYG{l+s+s1}{tracking}\PYG{l+s+s1}{\PYGZsq{}} \PYG{o+ow}{in} \PYG{n}{f}\PYG{p}{]}\PYG{p}{[}\PYG{l+m+mi}{0}\PYG{p}{]}
    \PYG{n}{velocity\PYGZus{}file} \PYG{o}{=} \PYG{p}{[}\PYG{n}{f} \PYG{k}{for} \PYG{n}{f} \PYG{o+ow}{in} \PYG{n}{filenames} \PYG{k}{if} \PYG{l+s+s1}{\PYGZsq{}}\PYG{l+s+s1}{velocity}\PYG{l+s+s1}{\PYGZsq{}} \PYG{o+ow}{in} \PYG{n}{f}\PYG{p}{]}\PYG{p}{[}\PYG{l+m+mi}{0}\PYG{p}{]}
    \PYG{n+nb}{print}\PYG{p}{(}\PYG{l+s+sa}{f}\PYG{l+s+s1}{\PYGZsq{}}\PYG{l+s+s1}{Tracking file: }\PYG{l+s+si}{\PYGZob{}}\PYG{n}{tracking\PYGZus{}file}\PYG{l+s+si}{\PYGZcb{}}\PYG{l+s+s1}{\PYGZsq{}}\PYG{p}{)}
    \PYG{n+nb}{print}\PYG{p}{(}\PYG{l+s+sa}{f}\PYG{l+s+s1}{\PYGZsq{}}\PYG{l+s+s1}{Velocity file: }\PYG{l+s+si}{\PYGZob{}}\PYG{n}{velocity\PYGZus{}file}\PYG{l+s+si}{\PYGZcb{}}\PYG{l+s+s1}{\PYGZsq{}}\PYG{p}{)}
\PYG{k}{except} \PYG{n+ne}{ImportError}\PYG{p}{:}
    \PYG{c+c1}{\PYGZsh{} Running locally — set your file paths here}
    \PYG{n}{tracking\PYGZus{}file} \PYG{o}{=} \PYG{l+s+s1}{\PYGZsq{}}\PYG{l+s+s1}{yourvideo\PYGZus{}tracking.csv}\PYG{l+s+s1}{\PYGZsq{}}  \PYG{c+c1}{\PYGZsh{} \PYGZlt{}\PYGZhy{}\PYGZhy{} change this}
    \PYG{n}{velocity\PYGZus{}file} \PYG{o}{=} \PYG{l+s+s1}{\PYGZsq{}}\PYG{l+s+s1}{yourvideo\PYGZus{}velocity.csv}\PYG{l+s+s1}{\PYGZsq{}}  \PYG{c+c1}{\PYGZsh{} \PYGZlt{}\PYGZhy{}\PYGZhy{} change this}
    \PYG{n+nb}{print}\PYG{p}{(}\PYG{l+s+sa}{f}\PYG{l+s+s1}{\PYGZsq{}}\PYG{l+s+s1}{Using local files: }\PYG{l+s+si}{\PYGZob{}}\PYG{n}{tracking\PYGZus{}file}\PYG{l+s+si}{\PYGZcb{}}\PYG{l+s+s1}{, }\PYG{l+s+si}{\PYGZob{}}\PYG{n}{velocity\PYGZus{}file}\PYG{l+s+si}{\PYGZcb{}}\PYG{l+s+s1}{\PYGZsq{}}\PYG{p}{)}
\end{sphinxVerbatim}

\end{sphinxuseclass}\end{sphinxVerbatimInput}

\end{sphinxuseclass}

\subsubsection{Load the data}
\label{\detokenize{FlyTracking/Script:load-the-data}}
\sphinxAtStartPar
Now let’s load both CSV files using \sphinxcode{\sphinxupquote{numpy}}. Each file has a header row, so we skip it with \sphinxcode{\sphinxupquote{skiprows=1}}.
\begin{quote}

\sphinxAtStartPar
\sphinxstylestrong{Task}: Run the cell below to load your data.
\end{quote}

\begin{sphinxuseclass}{cell}\begin{sphinxVerbatimInput}

\begin{sphinxuseclass}{cell_input}
\begin{sphinxVerbatim}[commandchars=\\\{\}]
\PYG{c+c1}{\PYGZsh{} Load tracking data: columns are Time (s), X (cm), Y (cm)}
\PYG{n}{tracking} \PYG{o}{=} \PYG{n}{np}\PYG{o}{.}\PYG{n}{loadtxt}\PYG{p}{(}\PYG{n}{tracking\PYGZus{}file}\PYG{p}{,} \PYG{n}{delimiter}\PYG{o}{=}\PYG{l+s+s1}{\PYGZsq{}}\PYG{l+s+s1}{,}\PYG{l+s+s1}{\PYGZsq{}}\PYG{p}{,} \PYG{n}{skiprows}\PYG{o}{=}\PYG{l+m+mi}{1}\PYG{p}{)}
\PYG{n}{time\PYGZus{}pos} \PYG{o}{=} \PYG{n}{tracking}\PYG{p}{[}\PYG{p}{:}\PYG{p}{,} \PYG{l+m+mi}{0}\PYG{p}{]}
\PYG{n}{x\PYGZus{}pos} \PYG{o}{=} \PYG{n}{tracking}\PYG{p}{[}\PYG{p}{:}\PYG{p}{,} \PYG{l+m+mi}{1}\PYG{p}{]}
\PYG{n}{y\PYGZus{}pos} \PYG{o}{=} \PYG{n}{tracking}\PYG{p}{[}\PYG{p}{:}\PYG{p}{,} \PYG{l+m+mi}{2}\PYG{p}{]}

\PYG{c+c1}{\PYGZsh{} Load velocity data: columns are Time (s), Velocity (mm/s)}
\PYG{n}{vel\PYGZus{}data} \PYG{o}{=} \PYG{n}{np}\PYG{o}{.}\PYG{n}{loadtxt}\PYG{p}{(}\PYG{n}{velocity\PYGZus{}file}\PYG{p}{,} \PYG{n}{delimiter}\PYG{o}{=}\PYG{l+s+s1}{\PYGZsq{}}\PYG{l+s+s1}{,}\PYG{l+s+s1}{\PYGZsq{}}\PYG{p}{,} \PYG{n}{skiprows}\PYG{o}{=}\PYG{l+m+mi}{1}\PYG{p}{)}
\PYG{n}{time\PYGZus{}vel} \PYG{o}{=} \PYG{n}{vel\PYGZus{}data}\PYG{p}{[}\PYG{p}{:}\PYG{p}{,} \PYG{l+m+mi}{0}\PYG{p}{]}
\PYG{n}{velocity} \PYG{o}{=} \PYG{n}{vel\PYGZus{}data}\PYG{p}{[}\PYG{p}{:}\PYG{p}{,} \PYG{l+m+mi}{1}\PYG{p}{]}

\PYG{n+nb}{print}\PYG{p}{(}\PYG{l+s+sa}{f}\PYG{l+s+s1}{\PYGZsq{}}\PYG{l+s+s1}{Tracking data: }\PYG{l+s+si}{\PYGZob{}}\PYG{n+nb}{len}\PYG{p}{(}\PYG{n}{time\PYGZus{}pos}\PYG{p}{)}\PYG{l+s+si}{\PYGZcb{}}\PYG{l+s+s1}{ frames}\PYG{l+s+s1}{\PYGZsq{}}\PYG{p}{)}
\PYG{n+nb}{print}\PYG{p}{(}\PYG{l+s+sa}{f}\PYG{l+s+s1}{\PYGZsq{}}\PYG{l+s+s1}{Velocity data: }\PYG{l+s+si}{\PYGZob{}}\PYG{n+nb}{len}\PYG{p}{(}\PYG{n}{time\PYGZus{}vel}\PYG{p}{)}\PYG{l+s+si}{\PYGZcb{}}\PYG{l+s+s1}{ time bins}\PYG{l+s+s1}{\PYGZsq{}}\PYG{p}{)}
\PYG{n+nb}{print}\PYG{p}{(}\PYG{l+s+sa}{f}\PYG{l+s+s1}{\PYGZsq{}}\PYG{l+s+s1}{Mean velocity: }\PYG{l+s+si}{\PYGZob{}}\PYG{n}{np}\PYG{o}{.}\PYG{n}{nanmean}\PYG{p}{(}\PYG{n}{velocity}\PYG{p}{)}\PYG{l+s+si}{:}\PYG{l+s+s1}{.2f}\PYG{l+s+si}{\PYGZcb{}}\PYG{l+s+s1}{ mm/s}\PYG{l+s+s1}{\PYGZsq{}}\PYG{p}{)}
\PYG{n+nb}{print}\PYG{p}{(}\PYG{l+s+sa}{f}\PYG{l+s+s1}{\PYGZsq{}}\PYG{l+s+s1}{SD velocity:   }\PYG{l+s+si}{\PYGZob{}}\PYG{n}{np}\PYG{o}{.}\PYG{n}{nanstd}\PYG{p}{(}\PYG{n}{velocity}\PYG{p}{)}\PYG{l+s+si}{:}\PYG{l+s+s1}{.2f}\PYG{l+s+si}{\PYGZcb{}}\PYG{l+s+s1}{ mm/s}\PYG{l+s+s1}{\PYGZsq{}}\PYG{p}{)}
\end{sphinxVerbatim}

\end{sphinxuseclass}\end{sphinxVerbatimInput}

\end{sphinxuseclass}

\subsubsection{Plot the fly path}
\label{\detokenize{FlyTracking/Script:plot-the-fly-path}}
\sphinxAtStartPar
The plot below shows the fly’s path through the arena, color\sphinxhyphen{}coded by time. Earlier positions are shown in purple/blue, and later positions in yellow/green.
\begin{quote}

\sphinxAtStartPar
\sphinxstylestrong{Task}: Run the cell below to plot the fly path. Edit the axis labels, title, or colormap (\sphinxcode{\sphinxupquote{cmap}}) if you’d like to customize it.
\end{quote}

\begin{sphinxuseclass}{cell}\begin{sphinxVerbatimInput}

\begin{sphinxuseclass}{cell_input}
\begin{sphinxVerbatim}[commandchars=\\\{\}]
\PYG{c+c1}{\PYGZsh{} Set the diameter of your dish (should match what you used in the script)}
\PYG{n}{diameter} \PYG{o}{=} \PYG{l+m+mi}{4}  \PYG{c+c1}{\PYGZsh{} cm}

\PYG{n}{fig}\PYG{p}{,} \PYG{n}{ax} \PYG{o}{=} \PYG{n}{plt}\PYG{o}{.}\PYG{n}{subplots}\PYG{p}{(}\PYG{n}{figsize}\PYG{o}{=}\PYG{p}{(}\PYG{l+m+mi}{6}\PYG{p}{,} \PYG{l+m+mi}{6}\PYG{p}{)}\PYG{p}{)}

\PYG{c+c1}{\PYGZsh{} Create color\PYGZhy{}coded path}
\PYG{n}{valid} \PYG{o}{=} \PYG{o}{\PYGZti{}}\PYG{n}{np}\PYG{o}{.}\PYG{n}{isnan}\PYG{p}{(}\PYG{n}{x\PYGZus{}pos}\PYG{p}{)} \PYG{o}{\PYGZam{}} \PYG{o}{\PYGZti{}}\PYG{n}{np}\PYG{o}{.}\PYG{n}{isnan}\PYG{p}{(}\PYG{n}{y\PYGZus{}pos}\PYG{p}{)}
\PYG{n}{points} \PYG{o}{=} \PYG{n}{np}\PYG{o}{.}\PYG{n}{array}\PYG{p}{(}\PYG{p}{[}\PYG{n}{x\PYGZus{}pos}\PYG{p}{[}\PYG{n}{valid}\PYG{p}{]}\PYG{p}{,} \PYG{n}{y\PYGZus{}pos}\PYG{p}{[}\PYG{n}{valid}\PYG{p}{]}\PYG{p}{]}\PYG{p}{)}\PYG{o}{.}\PYG{n}{T}\PYG{o}{.}\PYG{n}{reshape}\PYG{p}{(}\PYG{o}{\PYGZhy{}}\PYG{l+m+mi}{1}\PYG{p}{,} \PYG{l+m+mi}{1}\PYG{p}{,} \PYG{l+m+mi}{2}\PYG{p}{)}
\PYG{n}{segments} \PYG{o}{=} \PYG{n}{np}\PYG{o}{.}\PYG{n}{concatenate}\PYG{p}{(}\PYG{p}{[}\PYG{n}{points}\PYG{p}{[}\PYG{p}{:}\PYG{o}{\PYGZhy{}}\PYG{l+m+mi}{1}\PYG{p}{]}\PYG{p}{,} \PYG{n}{points}\PYG{p}{[}\PYG{l+m+mi}{1}\PYG{p}{:}\PYG{p}{]}\PYG{p}{]}\PYG{p}{,} \PYG{n}{axis}\PYG{o}{=}\PYG{l+m+mi}{1}\PYG{p}{)}

\PYG{n}{lc} \PYG{o}{=} \PYG{n}{LineCollection}\PYG{p}{(}\PYG{n}{segments}\PYG{p}{,} \PYG{n}{cmap}\PYG{o}{=}\PYG{l+s+s1}{\PYGZsq{}}\PYG{l+s+s1}{viridis}\PYG{l+s+s1}{\PYGZsq{}}\PYG{p}{,} \PYG{n}{linewidth}\PYG{o}{=}\PYG{l+m+mi}{2}\PYG{p}{)}
\PYG{n}{lc}\PYG{o}{.}\PYG{n}{set\PYGZus{}array}\PYG{p}{(}\PYG{n}{time\PYGZus{}pos}\PYG{p}{[}\PYG{n}{valid}\PYG{p}{]}\PYG{p}{[}\PYG{p}{:}\PYG{o}{\PYGZhy{}}\PYG{l+m+mi}{1}\PYG{p}{]}\PYG{p}{)}
\PYG{n}{ax}\PYG{o}{.}\PYG{n}{add\PYGZus{}collection}\PYG{p}{(}\PYG{n}{lc}\PYG{p}{)}

\PYG{n}{ax}\PYG{o}{.}\PYG{n}{set\PYGZus{}xlim}\PYG{p}{(}\PYG{l+m+mi}{0}\PYG{p}{,} \PYG{n}{diameter}\PYG{p}{)}
\PYG{n}{ax}\PYG{o}{.}\PYG{n}{set\PYGZus{}ylim}\PYG{p}{(}\PYG{n}{diameter}\PYG{p}{,} \PYG{l+m+mi}{0}\PYG{p}{)}  \PYG{c+c1}{\PYGZsh{} Inverted to match video orientation}
\PYG{n}{ax}\PYG{o}{.}\PYG{n}{set\PYGZus{}aspect}\PYG{p}{(}\PYG{l+s+s1}{\PYGZsq{}}\PYG{l+s+s1}{equal}\PYG{l+s+s1}{\PYGZsq{}}\PYG{p}{)}

\PYG{c+c1}{\PYGZsh{} \PYGZhy{}\PYGZhy{}\PYGZhy{} Customize these labels for your report \PYGZhy{}\PYGZhy{}\PYGZhy{}}
\PYG{n}{ax}\PYG{o}{.}\PYG{n}{set\PYGZus{}xlabel}\PYG{p}{(}\PYG{l+s+s1}{\PYGZsq{}}\PYG{l+s+s1}{X\PYGZhy{}coordinate (cm)}\PYG{l+s+s1}{\PYGZsq{}}\PYG{p}{,} \PYG{n}{fontsize}\PYG{o}{=}\PYG{l+m+mi}{12}\PYG{p}{)}
\PYG{n}{ax}\PYG{o}{.}\PYG{n}{set\PYGZus{}ylabel}\PYG{p}{(}\PYG{l+s+s1}{\PYGZsq{}}\PYG{l+s+s1}{Y\PYGZhy{}coordinate (cm)}\PYG{l+s+s1}{\PYGZsq{}}\PYG{p}{,} \PYG{n}{fontsize}\PYG{o}{=}\PYG{l+m+mi}{12}\PYG{p}{)}
\PYG{n}{ax}\PYG{o}{.}\PYG{n}{set\PYGZus{}title}\PYG{p}{(}\PYG{l+s+s1}{\PYGZsq{}}\PYG{l+s+s1}{Fly Path}\PYG{l+s+s1}{\PYGZsq{}}\PYG{p}{,} \PYG{n}{fontsize}\PYG{o}{=}\PYG{l+m+mi}{14}\PYG{p}{)}

\PYG{n}{cbar} \PYG{o}{=} \PYG{n}{fig}\PYG{o}{.}\PYG{n}{colorbar}\PYG{p}{(}\PYG{n}{lc}\PYG{p}{,} \PYG{n}{ax}\PYG{o}{=}\PYG{n}{ax}\PYG{p}{,} \PYG{n}{orientation}\PYG{o}{=}\PYG{l+s+s1}{\PYGZsq{}}\PYG{l+s+s1}{horizontal}\PYG{l+s+s1}{\PYGZsq{}}\PYG{p}{,} \PYG{n}{pad}\PYG{o}{=}\PYG{l+m+mf}{0.1}\PYG{p}{)}
\PYG{n}{cbar}\PYG{o}{.}\PYG{n}{set\PYGZus{}label}\PYG{p}{(}\PYG{l+s+s1}{\PYGZsq{}}\PYG{l+s+s1}{Time (s)}\PYG{l+s+s1}{\PYGZsq{}}\PYG{p}{)}

\PYG{n}{plt}\PYG{o}{.}\PYG{n}{tight\PYGZus{}layout}\PYG{p}{(}\PYG{p}{)}
\PYG{n}{plt}\PYG{o}{.}\PYG{n}{show}\PYG{p}{(}\PYG{p}{)}
\end{sphinxVerbatim}

\end{sphinxuseclass}\end{sphinxVerbatimInput}

\end{sphinxuseclass}

\subsubsection{Plot velocity over time}
\label{\detokenize{FlyTracking/Script:plot-velocity-over-time}}
\sphinxAtStartPar
The plot below shows how the fly’s velocity changes over the course of the recording.
\begin{quote}

\sphinxAtStartPar
\sphinxstylestrong{Task}: Run the cell below to plot the velocity trace. Edit the labels, colors, or axis limits as needed.
\end{quote}

\begin{sphinxuseclass}{cell}\begin{sphinxVerbatimInput}

\begin{sphinxuseclass}{cell_input}
\begin{sphinxVerbatim}[commandchars=\\\{\}]
\PYG{n}{fig}\PYG{p}{,} \PYG{n}{ax} \PYG{o}{=} \PYG{n}{plt}\PYG{o}{.}\PYG{n}{subplots}\PYG{p}{(}\PYG{n}{figsize}\PYG{o}{=}\PYG{p}{(}\PYG{l+m+mi}{10}\PYG{p}{,} \PYG{l+m+mi}{4}\PYG{p}{)}\PYG{p}{)}

\PYG{n}{ax}\PYG{o}{.}\PYG{n}{plot}\PYG{p}{(}\PYG{n}{time\PYGZus{}vel}\PYG{p}{,} \PYG{n}{velocity}\PYG{p}{,} \PYG{n}{linewidth}\PYG{o}{=}\PYG{l+m+mf}{1.5}\PYG{p}{,} \PYG{n}{color}\PYG{o}{=}\PYG{l+s+s1}{\PYGZsq{}}\PYG{l+s+s1}{steelblue}\PYG{l+s+s1}{\PYGZsq{}}\PYG{p}{)}

\PYG{c+c1}{\PYGZsh{} \PYGZhy{}\PYGZhy{}\PYGZhy{} Customize these for your report \PYGZhy{}\PYGZhy{}\PYGZhy{}}
\PYG{n}{ax}\PYG{o}{.}\PYG{n}{set\PYGZus{}xlabel}\PYG{p}{(}\PYG{l+s+s1}{\PYGZsq{}}\PYG{l+s+s1}{Time (s)}\PYG{l+s+s1}{\PYGZsq{}}\PYG{p}{,} \PYG{n}{fontsize}\PYG{o}{=}\PYG{l+m+mi}{12}\PYG{p}{)}
\PYG{n}{ax}\PYG{o}{.}\PYG{n}{set\PYGZus{}ylabel}\PYG{p}{(}\PYG{l+s+s1}{\PYGZsq{}}\PYG{l+s+s1}{Velocity (mm/s)}\PYG{l+s+s1}{\PYGZsq{}}\PYG{p}{,} \PYG{n}{fontsize}\PYG{o}{=}\PYG{l+m+mi}{12}\PYG{p}{)}
\PYG{n}{ax}\PYG{o}{.}\PYG{n}{set\PYGZus{}title}\PYG{p}{(}\PYG{l+s+s1}{\PYGZsq{}}\PYG{l+s+s1}{Fly Velocity Over Time}\PYG{l+s+s1}{\PYGZsq{}}\PYG{p}{,} \PYG{n}{fontsize}\PYG{o}{=}\PYG{l+m+mi}{14}\PYG{p}{)}

\PYG{c+c1}{\PYGZsh{} Uncomment the line below to set custom axis limits:}
\PYG{c+c1}{\PYGZsh{} ax.set\PYGZus{}xlim(0, 60)   \PYGZsh{} e.g., first 60 seconds}
\PYG{c+c1}{\PYGZsh{} ax.set\PYGZus{}ylim(0, 20)   \PYGZsh{} e.g., max 20 mm/s}

\PYG{n}{plt}\PYG{o}{.}\PYG{n}{tight\PYGZus{}layout}\PYG{p}{(}\PYG{p}{)}
\PYG{n}{plt}\PYG{o}{.}\PYG{n}{show}\PYG{p}{(}\PYG{p}{)}

\PYG{n+nb}{print}\PYG{p}{(}\PYG{l+s+sa}{f}\PYG{l+s+s1}{\PYGZsq{}}\PYG{l+s+s1}{Mean velocity: }\PYG{l+s+si}{\PYGZob{}}\PYG{n}{np}\PYG{o}{.}\PYG{n}{nanmean}\PYG{p}{(}\PYG{n}{velocity}\PYG{p}{)}\PYG{l+s+si}{:}\PYG{l+s+s1}{.2f}\PYG{l+s+si}{\PYGZcb{}}\PYG{l+s+s1}{ mm/s}\PYG{l+s+s1}{\PYGZsq{}}\PYG{p}{)}
\PYG{n+nb}{print}\PYG{p}{(}\PYG{l+s+sa}{f}\PYG{l+s+s1}{\PYGZsq{}}\PYG{l+s+s1}{SD velocity:   }\PYG{l+s+si}{\PYGZob{}}\PYG{n}{np}\PYG{o}{.}\PYG{n}{nanstd}\PYG{p}{(}\PYG{n}{velocity}\PYG{p}{)}\PYG{l+s+si}{:}\PYG{l+s+s1}{.2f}\PYG{l+s+si}{\PYGZcb{}}\PYG{l+s+s1}{ mm/s}\PYG{l+s+s1}{\PYGZsq{}}\PYG{p}{)}
\end{sphinxVerbatim}

\end{sphinxuseclass}\end{sphinxVerbatimInput}

\end{sphinxuseclass}

\subsubsection{Compare multiple flies \sphinxstyleemphasis{(optional)}}
\label{\detokenize{FlyTracking/Script:compare-multiple-flies-optional}}
\sphinxAtStartPar
If you have velocity data from multiple flies, you can plot them together and compute summary statistics.
\begin{quote}

\sphinxAtStartPar
\sphinxstylestrong{Task}: Upload additional velocity CSV files and add their filenames to the list below.
\end{quote}

\begin{sphinxuseclass}{cell}\begin{sphinxVerbatimInput}

\begin{sphinxuseclass}{cell_input}
\begin{sphinxVerbatim}[commandchars=\\\{\}]
\PYG{c+c1}{\PYGZsh{} Add your velocity CSV filenames here}
\PYG{n}{velocity\PYGZus{}files} \PYG{o}{=} \PYG{p}{[}
    \PYG{l+s+s1}{\PYGZsq{}}\PYG{l+s+s1}{fly1\PYGZus{}velocity.csv}\PYG{l+s+s1}{\PYGZsq{}}\PYG{p}{,}   \PYG{c+c1}{\PYGZsh{} \PYGZlt{}\PYGZhy{}\PYGZhy{} change these}
    \PYG{l+s+s1}{\PYGZsq{}}\PYG{l+s+s1}{fly2\PYGZus{}velocity.csv}\PYG{l+s+s1}{\PYGZsq{}}\PYG{p}{,}
    \PYG{c+c1}{\PYGZsh{} \PYGZsq{}fly3\PYGZus{}velocity.csv\PYGZsq{},}
\PYG{p}{]}

\PYG{n}{fig}\PYG{p}{,} \PYG{n}{ax} \PYG{o}{=} \PYG{n}{plt}\PYG{o}{.}\PYG{n}{subplots}\PYG{p}{(}\PYG{n}{figsize}\PYG{o}{=}\PYG{p}{(}\PYG{l+m+mi}{10}\PYG{p}{,} \PYG{l+m+mi}{5}\PYG{p}{)}\PYG{p}{)}
\PYG{n}{all\PYGZus{}means} \PYG{o}{=} \PYG{p}{[}\PYG{p}{]}

\PYG{k}{for} \PYG{n}{i}\PYG{p}{,} \PYG{n}{vf} \PYG{o+ow}{in} \PYG{n+nb}{enumerate}\PYG{p}{(}\PYG{n}{velocity\PYGZus{}files}\PYG{p}{)}\PYG{p}{:}
    \PYG{n}{data} \PYG{o}{=} \PYG{n}{np}\PYG{o}{.}\PYG{n}{loadtxt}\PYG{p}{(}\PYG{n}{vf}\PYG{p}{,} \PYG{n}{delimiter}\PYG{o}{=}\PYG{l+s+s1}{\PYGZsq{}}\PYG{l+s+s1}{,}\PYG{l+s+s1}{\PYGZsq{}}\PYG{p}{,} \PYG{n}{skiprows}\PYG{o}{=}\PYG{l+m+mi}{1}\PYG{p}{)}
    \PYG{n}{t} \PYG{o}{=} \PYG{n}{data}\PYG{p}{[}\PYG{p}{:}\PYG{p}{,} \PYG{l+m+mi}{0}\PYG{p}{]}
    \PYG{n}{v} \PYG{o}{=} \PYG{n}{data}\PYG{p}{[}\PYG{p}{:}\PYG{p}{,} \PYG{l+m+mi}{1}\PYG{p}{]}
    \PYG{n}{ax}\PYG{o}{.}\PYG{n}{plot}\PYG{p}{(}\PYG{n}{t}\PYG{p}{,} \PYG{n}{v}\PYG{p}{,} \PYG{n}{linewidth}\PYG{o}{=}\PYG{l+m+mf}{1.5}\PYG{p}{,} \PYG{n}{label}\PYG{o}{=}\PYG{l+s+sa}{f}\PYG{l+s+s1}{\PYGZsq{}}\PYG{l+s+s1}{Fly }\PYG{l+s+si}{\PYGZob{}}\PYG{n}{i}\PYG{+w}{ }\PYG{o}{+}\PYG{+w}{ }\PYG{l+m+mi}{1}\PYG{l+s+si}{\PYGZcb{}}\PYG{l+s+s1}{\PYGZsq{}}\PYG{p}{)}
    \PYG{n}{all\PYGZus{}means}\PYG{o}{.}\PYG{n}{append}\PYG{p}{(}\PYG{n}{np}\PYG{o}{.}\PYG{n}{nanmean}\PYG{p}{(}\PYG{n}{v}\PYG{p}{)}\PYG{p}{)}

\PYG{n}{ax}\PYG{o}{.}\PYG{n}{set\PYGZus{}xlabel}\PYG{p}{(}\PYG{l+s+s1}{\PYGZsq{}}\PYG{l+s+s1}{Time (s)}\PYG{l+s+s1}{\PYGZsq{}}\PYG{p}{,} \PYG{n}{fontsize}\PYG{o}{=}\PYG{l+m+mi}{12}\PYG{p}{)}
\PYG{n}{ax}\PYG{o}{.}\PYG{n}{set\PYGZus{}ylabel}\PYG{p}{(}\PYG{l+s+s1}{\PYGZsq{}}\PYG{l+s+s1}{Velocity (mm/s)}\PYG{l+s+s1}{\PYGZsq{}}\PYG{p}{,} \PYG{n}{fontsize}\PYG{o}{=}\PYG{l+m+mi}{12}\PYG{p}{)}
\PYG{n}{ax}\PYG{o}{.}\PYG{n}{set\PYGZus{}title}\PYG{p}{(}\PYG{l+s+s1}{\PYGZsq{}}\PYG{l+s+s1}{Velocity Comparison}\PYG{l+s+s1}{\PYGZsq{}}\PYG{p}{,} \PYG{n}{fontsize}\PYG{o}{=}\PYG{l+m+mi}{14}\PYG{p}{)}
\PYG{n}{ax}\PYG{o}{.}\PYG{n}{legend}\PYG{p}{(}\PYG{p}{)}
\PYG{n}{plt}\PYG{o}{.}\PYG{n}{tight\PYGZus{}layout}\PYG{p}{(}\PYG{p}{)}
\PYG{n}{plt}\PYG{o}{.}\PYG{n}{show}\PYG{p}{(}\PYG{p}{)}

\PYG{n+nb}{print}\PYG{p}{(}\PYG{l+s+sa}{f}\PYG{l+s+s1}{\PYGZsq{}}\PYG{l+s+se}{\PYGZbs{}n}\PYG{l+s+s1}{Mean velocity across flies: }\PYG{l+s+si}{\PYGZob{}}\PYG{n}{np}\PYG{o}{.}\PYG{n}{mean}\PYG{p}{(}\PYG{n}{all\PYGZus{}means}\PYG{p}{)}\PYG{l+s+si}{:}\PYG{l+s+s1}{.2f}\PYG{l+s+si}{\PYGZcb{}}\PYG{l+s+s1}{ mm/s}\PYG{l+s+s1}{\PYGZsq{}}\PYG{p}{)}
\PYG{n+nb}{print}\PYG{p}{(}\PYG{l+s+sa}{f}\PYG{l+s+s1}{\PYGZsq{}}\PYG{l+s+s1}{SD across flies: }\PYG{l+s+si}{\PYGZob{}}\PYG{n}{np}\PYG{o}{.}\PYG{n}{std}\PYG{p}{(}\PYG{n}{all\PYGZus{}means}\PYG{p}{)}\PYG{l+s+si}{:}\PYG{l+s+s1}{.2f}\PYG{l+s+si}{\PYGZcb{}}\PYG{l+s+s1}{ mm/s}\PYG{l+s+s1}{\PYGZsq{}}\PYG{p}{)}
\end{sphinxVerbatim}

\end{sphinxuseclass}\end{sphinxVerbatimInput}

\end{sphinxuseclass}



\subsection{About this Notebook}
\label{\detokenize{FlyTracking/Script:about-this-notebook}}
\sphinxAtStartPar
This notebook was created by \sphinxhref{https://github.com/ajuavinett}{Ashley Juavinett} for classes at UC San Diego. The fly tracker script is based on MATLAB code by Jeff Stafford.

\sphinxstepscope


\section{Fly Tracker (Colab)}
\label{\detokenize{FlyTracking/ColabNotebook:fly-tracker-colab}}\label{\detokenize{FlyTracking/ColabNotebook::doc}}
\sphinxAtStartPar
This notebook tracks a fruit fly against a light background in a video, calculates its position and velocity over time, and produces path and velocity plots.


\subsection{Notebook steps}
\label{\detokenize{FlyTracking/ColabNotebook:notebook-steps}}\begin{enumerate}
\sphinxsetlistlabels{\arabic}{enumi}{enumii}{}{.}%
\item {} 
\sphinxAtStartPar
Run the \sphinxstylestrong{Setup} cell to import packages.

\item {} 
\sphinxAtStartPar
\sphinxstylestrong{Mount your Google Drive} and copy the path to your video file.

\item {} 
\sphinxAtStartPar
Create the \sphinxstylestrong{Helper Functions}.

\item {} 
\sphinxAtStartPar
Run the ROI cell and \sphinxstylestrong{draw your ROI} on the first frame, then confirm it.

\item {} 
\sphinxAtStartPar
Run the \sphinxstylestrong{Track Fly} cell and wait for it to finish.

\item {} 
\sphinxAtStartPar
Run the \sphinxstylestrong{Plot} and \sphinxstylestrong{Velocity} cells to view results.

\item {} 
\sphinxAtStartPar
Re\sphinxhyphen{}run steps 3\sphinxhyphen{}7 if you would like to track an additional video.

\end{enumerate}


\subsection{Setup}
\label{\detokenize{FlyTracking/ColabNotebook:setup}}
\begin{sphinxuseclass}{cell}\begin{sphinxVerbatimInput}

\begin{sphinxuseclass}{cell_input}
\begin{sphinxVerbatim}[commandchars=\\\{\}]
\PYG{c+c1}{\PYGZsh{}@title Install \PYGZam{} import packages \PYGZob{} display\PYGZhy{}mode: \PYGZdq{}form\PYGZdq{} \PYGZcb{}}
\PYG{k+kn}{import}\PYG{+w}{ }\PYG{n+nn}{numpy}\PYG{+w}{ }\PYG{k}{as}\PYG{+w}{ }\PYG{n+nn}{np}
\PYG{k+kn}{import}\PYG{+w}{ }\PYG{n+nn}{cv2}
\PYG{k+kn}{import}\PYG{+w}{ }\PYG{n+nn}{matplotlib}\PYG{n+nn}{.}\PYG{n+nn}{pyplot}\PYG{+w}{ }\PYG{k}{as}\PYG{+w}{ }\PYG{n+nn}{plt}
\PYG{k+kn}{import}\PYG{+w}{ }\PYG{n+nn}{matplotlib}\PYG{n+nn}{.}\PYG{n+nn}{colors}\PYG{+w}{ }\PYG{k}{as}\PYG{+w}{ }\PYG{n+nn}{mcolors}
\PYG{k+kn}{from}\PYG{+w}{ }\PYG{n+nn}{matplotlib}\PYG{n+nn}{.}\PYG{n+nn}{collections}\PYG{+w}{ }\PYG{k+kn}{import} \PYG{n}{LineCollection}
\PYG{k+kn}{from}\PYG{+w}{ }\PYG{n+nn}{scipy}\PYG{n+nn}{.}\PYG{n+nn}{spatial}\PYG{n+nn}{.}\PYG{n+nn}{distance}\PYG{+w}{ }\PYG{k+kn}{import} \PYG{n}{euclidean}
\PYG{k+kn}{from}\PYG{+w}{ }\PYG{n+nn}{google}\PYG{n+nn}{.}\PYG{n+nn}{colab}\PYG{+w}{ }\PYG{k+kn}{import} \PYG{n}{files}\PYG{p}{,} \PYG{n}{output}
\PYG{k+kn}{from}\PYG{+w}{ }\PYG{n+nn}{IPython}\PYG{n+nn}{.}\PYG{n+nn}{display}\PYG{+w}{ }\PYG{k+kn}{import} \PYG{n}{display}\PYG{p}{,} \PYG{n}{clear\PYGZus{}output}\PYG{p}{,} \PYG{n}{HTML}\PYG{p}{,} \PYG{n}{Image} \PYG{k}{as} \PYG{n}{IPImage}
\PYG{k+kn}{import}\PYG{+w}{ }\PYG{n+nn}{ipywidgets}\PYG{+w}{ }\PYG{k}{as}\PYG{+w}{ }\PYG{n+nn}{widgets}
\PYG{k+kn}{import}\PYG{+w}{ }\PYG{n+nn}{os}
\PYG{k+kn}{import}\PYG{+w}{ }\PYG{n+nn}{base64}

\PYG{c+c1}{\PYGZsh{} Initialize result containers — each is created independently so adding}
\PYG{c+c1}{\PYGZsh{} new lists in future won\PYGZsq{}t require a full session restart}
\PYG{c+c1}{\PYGZsh{} You can uncomment the following to reset all results if needed:}
\PYG{c+c1}{\PYGZsh{}all\PYGZus{}corrected  = []}
\PYG{c+c1}{\PYGZsh{}all\PYGZus{}velocity   = []}
\PYG{c+c1}{\PYGZsh{}all\PYGZus{}fps        = []}
\PYG{c+c1}{\PYGZsh{}all\PYGZus{}video\PYGZus{}paths = []}
\PYG{c+c1}{\PYGZsh{}print(\PYGZsq{}Results cleared — ready to process new videos.\PYGZsq{})}

\PYG{k}{if} \PYG{l+s+s1}{\PYGZsq{}}\PYG{l+s+s1}{all\PYGZus{}corrected}\PYG{l+s+s1}{\PYGZsq{}} \PYG{o+ow}{not} \PYG{o+ow}{in} \PYG{n+nb}{dir}\PYG{p}{(}\PYG{p}{)}\PYG{p}{:} \PYG{n}{all\PYGZus{}corrected} \PYG{o}{=} \PYG{p}{[}\PYG{p}{]}
\PYG{k}{if} \PYG{l+s+s1}{\PYGZsq{}}\PYG{l+s+s1}{all\PYGZus{}velocity}\PYG{l+s+s1}{\PYGZsq{}}  \PYG{o+ow}{not} \PYG{o+ow}{in} \PYG{n+nb}{dir}\PYG{p}{(}\PYG{p}{)}\PYG{p}{:} \PYG{n}{all\PYGZus{}velocity}  \PYG{o}{=} \PYG{p}{[}\PYG{p}{]}
\PYG{k}{if} \PYG{l+s+s1}{\PYGZsq{}}\PYG{l+s+s1}{all\PYGZus{}fps}\PYG{l+s+s1}{\PYGZsq{}}       \PYG{o+ow}{not} \PYG{o+ow}{in} \PYG{n+nb}{dir}\PYG{p}{(}\PYG{p}{)}\PYG{p}{:} \PYG{n}{all\PYGZus{}fps}       \PYG{o}{=} \PYG{p}{[}\PYG{p}{]}
\PYG{k}{if} \PYG{l+s+s1}{\PYGZsq{}}\PYG{l+s+s1}{all\PYGZus{}video\PYGZus{}paths}\PYG{l+s+s1}{\PYGZsq{}} \PYG{o+ow}{not} \PYG{o+ow}{in} \PYG{n+nb}{dir}\PYG{p}{(}\PYG{p}{)}\PYG{p}{:} \PYG{n}{all\PYGZus{}video\PYGZus{}paths} \PYG{o}{=} \PYG{p}{[}\PYG{p}{]}
\PYG{n+nb}{print}\PYG{p}{(}\PYG{l+s+sa}{f}\PYG{l+s+s1}{\PYGZsq{}}\PYG{l+s+s1}{Ready. }\PYG{l+s+si}{\PYGZob{}}\PYG{n+nb}{len}\PYG{p}{(}\PYG{n}{all\PYGZus{}corrected}\PYG{p}{)}\PYG{l+s+si}{\PYGZcb{}}\PYG{l+s+s1}{ video(s) processed so far.}\PYG{l+s+s1}{\PYGZsq{}}\PYG{p}{)}

\PYG{n+nb}{print}\PYG{p}{(}\PYG{l+s+s1}{\PYGZsq{}}\PYG{l+s+s1}{All packages loaded successfully!}\PYG{l+s+s1}{\PYGZsq{}}\PYG{p}{)}
\end{sphinxVerbatim}

\end{sphinxuseclass}\end{sphinxVerbatimInput}

\end{sphinxuseclass}

\subsection{Mount your drive}
\label{\detokenize{FlyTracking/ColabNotebook:mount-your-drive}}
\sphinxAtStartPar
Run the cell below to mount your drive, which should contain your video files.

\sphinxAtStartPar
If you’d prefer to upload your videos, you can do that instead (and do not need to run the cell below. See instructions here.)

\begin{sphinxuseclass}{cell}\begin{sphinxVerbatimInput}

\begin{sphinxuseclass}{cell_input}
\begin{sphinxVerbatim}[commandchars=\\\{\}]
\PYG{k+kn}{from}\PYG{+w}{ }\PYG{n+nn}{google}\PYG{n+nn}{.}\PYG{n+nn}{colab}\PYG{+w}{ }\PYG{k+kn}{import} \PYG{n}{drive}
\PYG{n}{drive}\PYG{o}{.}\PYG{n}{mount}\PYG{p}{(}\PYG{l+s+s1}{\PYGZsq{}}\PYG{l+s+s1}{/content/drive}\PYG{l+s+s1}{\PYGZsq{}}\PYG{p}{)}
\end{sphinxVerbatim}

\end{sphinxuseclass}\end{sphinxVerbatimInput}

\end{sphinxuseclass}

\subsection{Select Video File}
\label{\detokenize{FlyTracking/ColabNotebook:select-video-file}}
\sphinxAtStartPar
Update the path below to your video file. You’ll process one video at a time, and results will accumulate in \sphinxcode{\sphinxupquote{all\_corrected}} and \sphinxcode{\sphinxupquote{all\_velocity}}.
\begin{quote}

\sphinxAtStartPar
\sphinxstylestrong{Tip}: In the window at left, you can use the three vertical dots to copy and paste the path to your video file.
\end{quote}

\begin{sphinxuseclass}{cell}\begin{sphinxVerbatimInput}

\begin{sphinxuseclass}{cell_input}
\begin{sphinxVerbatim}[commandchars=\\\{\}]
\PYG{c+c1}{\PYGZsh{} Set the current video file to process}
\PYG{n}{current\PYGZus{}video} \PYG{o}{=} \PYG{l+s+s1}{\PYGZsq{}}\PYG{l+s+s1}{/content/drive/MyDrive/video1.avi}\PYG{l+s+s1}{\PYGZsq{}}

\PYG{k}{if} \PYG{n}{os}\PYG{o}{.}\PYG{n}{path}\PYG{o}{.}\PYG{n}{exists}\PYG{p}{(}\PYG{n}{current\PYGZus{}video}\PYG{p}{)}\PYG{p}{:}
    \PYG{n}{size\PYGZus{}mb} \PYG{o}{=} \PYG{n}{os}\PYG{o}{.}\PYG{n}{path}\PYG{o}{.}\PYG{n}{getsize}\PYG{p}{(}\PYG{n}{current\PYGZus{}video}\PYG{p}{)} \PYG{o}{/} \PYG{l+m+mf}{1e6}
    \PYG{n+nb}{print}\PYG{p}{(}\PYG{l+s+sa}{f}\PYG{l+s+s1}{\PYGZsq{}}\PYG{l+s+s1}{File found: }\PYG{l+s+si}{\PYGZob{}}\PYG{n}{os}\PYG{o}{.}\PYG{n}{path}\PYG{o}{.}\PYG{n}{basename}\PYG{p}{(}\PYG{n}{current\PYGZus{}video}\PYG{p}{)}\PYG{l+s+si}{\PYGZcb{}}\PYG{l+s+s1}{ (}\PYG{l+s+si}{\PYGZob{}}\PYG{n}{size\PYGZus{}mb}\PYG{l+s+si}{:}\PYG{l+s+s1}{.1f}\PYG{l+s+si}{\PYGZcb{}}\PYG{l+s+s1}{ MB)}\PYG{l+s+s1}{\PYGZsq{}}\PYG{p}{)}
\PYG{k}{else}\PYG{p}{:}
    \PYG{n+nb}{print}\PYG{p}{(}\PYG{l+s+sa}{f}\PYG{l+s+s1}{\PYGZsq{}}\PYG{l+s+s1}{File not found: }\PYG{l+s+si}{\PYGZob{}}\PYG{n}{current\PYGZus{}video}\PYG{l+s+si}{\PYGZcb{}}\PYG{l+s+s1}{\PYGZsq{}}\PYG{p}{)}
    \PYG{n+nb}{print}\PYG{p}{(}\PYG{p}{)}
    \PYG{n+nb}{print}\PYG{p}{(}\PYG{l+s+s1}{\PYGZsq{}}\PYG{l+s+s1}{Troubleshooting tips:}\PYG{l+s+s1}{\PYGZsq{}}\PYG{p}{)}
    \PYG{n+nb}{print}\PYG{p}{(}\PYG{l+s+s1}{\PYGZsq{}}\PYG{l+s+s1}{  1. Make sure you have run the }\PYG{l+s+s1}{\PYGZdq{}}\PYG{l+s+s1}{Mount your drive}\PYG{l+s+s1}{\PYGZdq{}}\PYG{l+s+s1}{ cell above.}\PYG{l+s+s1}{\PYGZsq{}}\PYG{p}{)}
    \PYG{n+nb}{print}\PYG{p}{(}\PYG{l+s+s1}{\PYGZsq{}}\PYG{l+s+s1}{  2. In the left sidebar, click the folder icon to open the file browser.}\PYG{l+s+s1}{\PYGZsq{}}\PYG{p}{)}
    \PYG{n+nb}{print}\PYG{p}{(}\PYG{l+s+s1}{\PYGZsq{}}\PYG{l+s+s1}{  3. Navigate to your video file, right\PYGZhy{}click it, and select }\PYG{l+s+s1}{\PYGZdq{}}\PYG{l+s+s1}{Copy path}\PYG{l+s+s1}{\PYGZdq{}}\PYG{l+s+s1}{.}\PYG{l+s+s1}{\PYGZsq{}}\PYG{p}{)}
    \PYG{n+nb}{print}\PYG{p}{(}\PYG{l+s+s1}{\PYGZsq{}}\PYG{l+s+s1}{  4. Paste the copied path above to replace the current\PYGZus{}video value.}\PYG{l+s+s1}{\PYGZsq{}}\PYG{p}{)}
    \PYG{n+nb}{print}\PYG{p}{(}\PYG{p}{)}
    \PYG{c+c1}{\PYGZsh{} Show video files in the same folder if it exists, to help spot typos}
    \PYG{n}{folder} \PYG{o}{=} \PYG{n}{os}\PYG{o}{.}\PYG{n}{path}\PYG{o}{.}\PYG{n}{dirname}\PYG{p}{(}\PYG{n}{current\PYGZus{}video}\PYG{p}{)}
    \PYG{k}{if} \PYG{n}{os}\PYG{o}{.}\PYG{n}{path}\PYG{o}{.}\PYG{n}{exists}\PYG{p}{(}\PYG{n}{folder}\PYG{p}{)}\PYG{p}{:}
        \PYG{n}{videos} \PYG{o}{=} \PYG{p}{[}\PYG{n}{f} \PYG{k}{for} \PYG{n}{f} \PYG{o+ow}{in} \PYG{n}{os}\PYG{o}{.}\PYG{n}{listdir}\PYG{p}{(}\PYG{n}{folder}\PYG{p}{)}
                  \PYG{k}{if} \PYG{n}{f}\PYG{o}{.}\PYG{n}{lower}\PYG{p}{(}\PYG{p}{)}\PYG{o}{.}\PYG{n}{endswith}\PYG{p}{(}\PYG{p}{(}\PYG{l+s+s1}{\PYGZsq{}}\PYG{l+s+s1}{.avi}\PYG{l+s+s1}{\PYGZsq{}}\PYG{p}{,} \PYG{l+s+s1}{\PYGZsq{}}\PYG{l+s+s1}{.mp4}\PYG{l+s+s1}{\PYGZsq{}}\PYG{p}{,} \PYG{l+s+s1}{\PYGZsq{}}\PYG{l+s+s1}{.mov}\PYG{l+s+s1}{\PYGZsq{}}\PYG{p}{,} \PYG{l+s+s1}{\PYGZsq{}}\PYG{l+s+s1}{.mkv}\PYG{l+s+s1}{\PYGZsq{}}\PYG{p}{,} \PYG{l+s+s1}{\PYGZsq{}}\PYG{l+s+s1}{.mj2}\PYG{l+s+s1}{\PYGZsq{}}\PYG{p}{,} \PYG{l+s+s1}{\PYGZsq{}}\PYG{l+s+s1}{.mpg}\PYG{l+s+s1}{\PYGZsq{}}\PYG{p}{,} \PYG{l+s+s1}{\PYGZsq{}}\PYG{l+s+s1}{.m4v}\PYG{l+s+s1}{\PYGZsq{}}\PYG{p}{)}\PYG{p}{)}\PYG{p}{]}
        \PYG{k}{if} \PYG{n}{videos}\PYG{p}{:}
            \PYG{n+nb}{print}\PYG{p}{(}\PYG{l+s+sa}{f}\PYG{l+s+s1}{\PYGZsq{}}\PYG{l+s+s1}{Video files found in }\PYG{l+s+si}{\PYGZob{}}\PYG{n}{folder}\PYG{l+s+si}{\PYGZcb{}}\PYG{l+s+s1}{:}\PYG{l+s+s1}{\PYGZsq{}}\PYG{p}{)}
            \PYG{k}{for} \PYG{n}{v} \PYG{o+ow}{in} \PYG{n+nb}{sorted}\PYG{p}{(}\PYG{n}{videos}\PYG{p}{)}\PYG{p}{:}
                \PYG{n+nb}{print}\PYG{p}{(}\PYG{l+s+sa}{f}\PYG{l+s+s1}{\PYGZsq{}}\PYG{l+s+s1}{  }\PYG{l+s+si}{\PYGZob{}}\PYG{n}{os}\PYG{o}{.}\PYG{n}{path}\PYG{o}{.}\PYG{n}{join}\PYG{p}{(}\PYG{n}{folder}\PYG{p}{,}\PYG{+w}{ }\PYG{n}{v}\PYG{p}{)}\PYG{l+s+si}{\PYGZcb{}}\PYG{l+s+s1}{\PYGZsq{}}\PYG{p}{)}
        \PYG{k}{else}\PYG{p}{:}
            \PYG{n+nb}{print}\PYG{p}{(}\PYG{l+s+sa}{f}\PYG{l+s+s1}{\PYGZsq{}}\PYG{l+s+s1}{No video files found in }\PYG{l+s+si}{\PYGZob{}}\PYG{n}{folder}\PYG{l+s+si}{\PYGZcb{}}\PYG{l+s+s1}{ — check that you are looking in the right folder.}\PYG{l+s+s1}{\PYGZsq{}}\PYG{p}{)}
    \PYG{k}{else}\PYG{p}{:}
        \PYG{n+nb}{print}\PYG{p}{(}\PYG{l+s+sa}{f}\PYG{l+s+s1}{\PYGZsq{}}\PYG{l+s+s1}{Folder not found: }\PYG{l+s+si}{\PYGZob{}}\PYG{n}{folder}\PYG{l+s+si}{\PYGZcb{}}\PYG{l+s+s1}{\PYGZsq{}}\PYG{p}{)}
        \PYG{n+nb}{print}\PYG{p}{(}\PYG{l+s+s1}{\PYGZsq{}}\PYG{l+s+s1}{  Make sure your Drive is mounted and the path is correct.}\PYG{l+s+s1}{\PYGZsq{}}\PYG{p}{)}
\end{sphinxVerbatim}

\end{sphinxuseclass}\end{sphinxVerbatimInput}

\end{sphinxuseclass}
\begin{sphinxuseclass}{cell}\begin{sphinxVerbatimInput}

\begin{sphinxuseclass}{cell_input}
\begin{sphinxVerbatim}[commandchars=\\\{\}]
\PYG{c+c1}{\PYGZsh{}@title Define helper functions (fly\PYGZus{}finder, dist\PYGZus{}filter, interpolate\PYGZus{}pos) \PYGZob{} display\PYGZhy{}mode: \PYGZdq{}form\PYGZdq{} \PYGZcb{}}

\PYG{c+c1}{\PYGZsh{} \PYGZhy{}\PYGZhy{}\PYGZhy{} Parameters \PYGZhy{}\PYGZhy{}\PYGZhy{}}
\PYG{n}{frame\PYGZus{}rate} \PYG{o}{=} \PYG{k+kc}{None}  \PYG{c+c1}{\PYGZsh{} auto\PYGZhy{}detect from video (set e.g. frame\PYGZus{}rate = 30 to override)}
\PYG{n}{diameter} \PYG{o}{=} \PYG{l+m+mi}{4}       \PYG{c+c1}{\PYGZsh{} dish diameter in centimeters}
\PYG{n}{search\PYGZus{}size} \PYG{o}{=} \PYG{l+m+mi}{20}   \PYG{c+c1}{\PYGZsh{} fly search window size in pixels}
\PYG{n}{per\PYGZus{}pixel\PYGZus{}threshold} \PYG{o}{=} \PYG{l+m+mf}{1.5}
\PYG{n}{bin\PYGZus{}size} \PYG{o}{=} \PYG{l+m+mi}{1}       \PYG{c+c1}{\PYGZsh{} velocity bin size in seconds}
\PYG{c+c1}{\PYGZsh{} \PYGZhy{}\PYGZhy{}\PYGZhy{}\PYGZhy{}\PYGZhy{}\PYGZhy{}\PYGZhy{}\PYGZhy{}\PYGZhy{}\PYGZhy{}\PYGZhy{}\PYGZhy{}\PYGZhy{}\PYGZhy{}\PYGZhy{}\PYGZhy{}\PYGZhy{}\PYGZhy{}\PYGZhy{}}


\PYG{k}{def}\PYG{+w}{ }\PYG{n+nf}{fly\PYGZus{}finder}\PYG{p}{(}\PYG{n}{roi\PYGZus{}image}\PYG{p}{,} \PYG{n}{half\PYGZus{}search}\PYG{p}{,} \PYG{n}{threshold}\PYG{p}{,} \PYG{n}{flip}\PYG{o}{=}\PYG{k+kc}{True}\PYG{p}{)}\PYG{p}{:}
\PYG{+w}{    }\PYG{l+s+sd}{\PYGZdq{}\PYGZdq{}\PYGZdq{}}
\PYG{l+s+sd}{    Find a fly (dark region) in a grayscale image.}
\PYG{l+s+sd}{    Locates the darkest pixel, retrieves a search area around it,}
\PYG{l+s+sd}{    and finds the center of pixel intensity.}

\PYG{l+s+sd}{    Returns (x, y) position or (NaN, NaN) if not found.}
\PYG{l+s+sd}{    \PYGZdq{}\PYGZdq{}\PYGZdq{}}
    \PYG{k}{if} \PYG{n}{flip}\PYG{p}{:}
        \PYG{n}{val} \PYG{o}{=} \PYG{n}{np}\PYG{o}{.}\PYG{n}{nanmin}\PYG{p}{(}\PYG{n}{roi\PYGZus{}image}\PYG{p}{)}
    \PYG{k}{else}\PYG{p}{:}
        \PYG{n}{val} \PYG{o}{=} \PYG{n}{np}\PYG{o}{.}\PYG{n}{nanmax}\PYG{p}{(}\PYG{n}{roi\PYGZus{}image}\PYG{p}{)}

    \PYG{n}{ys}\PYG{p}{,} \PYG{n}{xs} \PYG{o}{=} \PYG{n}{np}\PYG{o}{.}\PYG{n}{where}\PYG{p}{(}\PYG{n}{roi\PYGZus{}image} \PYG{o}{==} \PYG{n}{val}\PYG{p}{)}
    \PYG{n}{xpos} \PYG{o}{=} \PYG{n}{np}\PYG{o}{.}\PYG{n}{mean}\PYG{p}{(}\PYG{n}{xs}\PYG{p}{)}
    \PYG{n}{ypos} \PYG{o}{=} \PYG{n}{np}\PYG{o}{.}\PYG{n}{mean}\PYG{p}{(}\PYG{n}{ys}\PYG{p}{)}

    \PYG{n}{h}\PYG{p}{,} \PYG{n}{w} \PYG{o}{=} \PYG{n}{roi\PYGZus{}image}\PYG{o}{.}\PYG{n}{shape}
    \PYG{n}{left} \PYG{o}{=} \PYG{n+nb}{max}\PYG{p}{(}\PYG{n+nb}{int}\PYG{p}{(}\PYG{n+nb}{round}\PYG{p}{(}\PYG{n}{xpos}\PYG{p}{)} \PYG{o}{\PYGZhy{}} \PYG{n}{half\PYGZus{}search}\PYG{p}{)}\PYG{p}{,} \PYG{l+m+mi}{0}\PYG{p}{)}
    \PYG{n}{right} \PYG{o}{=} \PYG{n+nb}{min}\PYG{p}{(}\PYG{n+nb}{int}\PYG{p}{(}\PYG{n+nb}{round}\PYG{p}{(}\PYG{n}{xpos}\PYG{p}{)} \PYG{o}{+} \PYG{n}{half\PYGZus{}search}\PYG{p}{)}\PYG{p}{,} \PYG{n}{w} \PYG{o}{\PYGZhy{}} \PYG{l+m+mi}{1}\PYG{p}{)}
    \PYG{n}{top} \PYG{o}{=} \PYG{n+nb}{max}\PYG{p}{(}\PYG{n+nb}{int}\PYG{p}{(}\PYG{n+nb}{round}\PYG{p}{(}\PYG{n}{ypos}\PYG{p}{)} \PYG{o}{\PYGZhy{}} \PYG{n}{half\PYGZus{}search}\PYG{p}{)}\PYG{p}{,} \PYG{l+m+mi}{0}\PYG{p}{)}
    \PYG{n}{bottom} \PYG{o}{=} \PYG{n+nb}{min}\PYG{p}{(}\PYG{n+nb}{int}\PYG{p}{(}\PYG{n+nb}{round}\PYG{p}{(}\PYG{n}{ypos}\PYG{p}{)} \PYG{o}{+} \PYG{n}{half\PYGZus{}search}\PYG{p}{)}\PYG{p}{,} \PYG{n}{h} \PYG{o}{\PYGZhy{}} \PYG{l+m+mi}{1}\PYG{p}{)}

    \PYG{n}{search\PYGZus{}area} \PYG{o}{=} \PYG{n}{roi\PYGZus{}image}\PYG{p}{[}\PYG{n}{top}\PYG{p}{:}\PYG{n}{bottom}\PYG{o}{+}\PYG{l+m+mi}{1}\PYG{p}{,} \PYG{n}{left}\PYG{p}{:}\PYG{n}{right}\PYG{o}{+}\PYG{l+m+mi}{1}\PYG{p}{]}\PYG{o}{.}\PYG{n}{astype}\PYG{p}{(}\PYG{n}{np}\PYG{o}{.}\PYG{n}{float64}\PYG{p}{)}

    \PYG{k}{if} \PYG{n}{flip}\PYG{p}{:}
        \PYG{n}{search\PYGZus{}area} \PYG{o}{=} \PYG{l+m+mf}{255.0} \PYG{o}{\PYGZhy{}} \PYG{n}{search\PYGZus{}area}

    \PYG{n}{total} \PYG{o}{=} \PYG{n}{np}\PYG{o}{.}\PYG{n}{sum}\PYG{p}{(}\PYG{n}{search\PYGZus{}area}\PYG{p}{)}

    \PYG{k}{if} \PYG{n}{total} \PYG{o}{\PYGZgt{}}\PYG{o}{=} \PYG{n}{threshold}\PYG{p}{:}
        \PYG{n}{x\PYGZus{}indices} \PYG{o}{=} \PYG{n}{np}\PYG{o}{.}\PYG{n}{arange}\PYG{p}{(}\PYG{n}{search\PYGZus{}area}\PYG{o}{.}\PYG{n}{shape}\PYG{p}{[}\PYG{l+m+mi}{1}\PYG{p}{]}\PYG{p}{)}
        \PYG{n}{y\PYGZus{}indices} \PYG{o}{=} \PYG{n}{np}\PYG{o}{.}\PYG{n}{arange}\PYG{p}{(}\PYG{n}{search\PYGZus{}area}\PYG{o}{.}\PYG{n}{shape}\PYG{p}{[}\PYG{l+m+mi}{0}\PYG{p}{]}\PYG{p}{)}
        \PYG{n}{x} \PYG{o}{=} \PYG{n}{np}\PYG{o}{.}\PYG{n}{sum}\PYG{p}{(}\PYG{n}{search\PYGZus{}area} \PYG{o}{@} \PYG{n}{x\PYGZus{}indices}\PYG{p}{)} \PYG{o}{/} \PYG{n}{total} \PYG{o}{+} \PYG{n}{left}
        \PYG{n}{y} \PYG{o}{=} \PYG{n}{np}\PYG{o}{.}\PYG{n}{sum}\PYG{p}{(}\PYG{n}{search\PYGZus{}area}\PYG{o}{.}\PYG{n}{T} \PYG{o}{@} \PYG{n}{y\PYGZus{}indices}\PYG{p}{)} \PYG{o}{/} \PYG{n}{total} \PYG{o}{+} \PYG{n}{top}
        \PYG{k}{return} \PYG{n}{x}\PYG{p}{,} \PYG{n}{y}
    \PYG{k}{else}\PYG{p}{:}
        \PYG{k}{return} \PYG{n}{np}\PYG{o}{.}\PYG{n}{nan}\PYG{p}{,} \PYG{n}{np}\PYG{o}{.}\PYG{n}{nan}


\PYG{k}{def}\PYG{+w}{ }\PYG{n+nf}{dist\PYGZus{}filter}\PYG{p}{(}\PYG{n}{array}\PYG{p}{,} \PYG{n}{tele\PYGZus{}dist\PYGZus{}threshold}\PYG{p}{,} \PYG{n}{num\PYGZus{}avg}\PYG{o}{=}\PYG{l+m+mi}{5}\PYG{p}{)}\PYG{p}{:}
\PYG{+w}{    }\PYG{l+s+sd}{\PYGZdq{}\PYGZdq{}\PYGZdq{}}
\PYG{l+s+sd}{    Teleport filter: removes spurious points where fly position}
\PYG{l+s+sd}{    jumps far from the mean of surrounding frames.}

\PYG{l+s+sd}{    array: Nx3 array [time, x, y]}
\PYG{l+s+sd}{    \PYGZdq{}\PYGZdq{}\PYGZdq{}}
    \PYG{n}{filtered} \PYG{o}{=} \PYG{n}{array}\PYG{o}{.}\PYG{n}{copy}\PYG{p}{(}\PYG{p}{)}
    \PYG{n}{n} \PYG{o}{=} \PYG{n+nb}{len}\PYG{p}{(}\PYG{n}{filtered}\PYG{p}{)}
    \PYG{n}{x} \PYG{o}{=} \PYG{n}{filtered}\PYG{p}{[}\PYG{p}{:}\PYG{p}{,} \PYG{l+m+mi}{1}\PYG{p}{]}
    \PYG{n}{y} \PYG{o}{=} \PYG{n}{filtered}\PYG{p}{[}\PYG{p}{:}\PYG{p}{,} \PYG{l+m+mi}{2}\PYG{p}{]}

    \PYG{k}{def}\PYG{+w}{ }\PYG{n+nf}{rolling\PYGZus{}nanmean}\PYG{p}{(}\PYG{n}{arr}\PYG{p}{,} \PYG{n}{window}\PYG{p}{,} \PYG{n}{forward}\PYG{o}{=}\PYG{k+kc}{False}\PYG{p}{)}\PYG{p}{:}
\PYG{+w}{        }\PYG{l+s+sd}{\PYGZdq{}\PYGZdq{}\PYGZdq{}Rolling mean over \PYGZsq{}window\PYGZsq{} preceding (or following) valid elements.\PYGZdq{}\PYGZdq{}\PYGZdq{}}
        \PYG{n}{valid} \PYG{o}{=} \PYG{p}{(}\PYG{o}{\PYGZti{}}\PYG{n}{np}\PYG{o}{.}\PYG{n}{isnan}\PYG{p}{(}\PYG{n}{arr}\PYG{p}{)}\PYG{p}{)}\PYG{o}{.}\PYG{n}{astype}\PYG{p}{(}\PYG{n}{np}\PYG{o}{.}\PYG{n}{float64}\PYG{p}{)}
        \PYG{n}{filled} \PYG{o}{=} \PYG{n}{np}\PYG{o}{.}\PYG{n}{where}\PYG{p}{(}\PYG{n}{np}\PYG{o}{.}\PYG{n}{isnan}\PYG{p}{(}\PYG{n}{arr}\PYG{p}{)}\PYG{p}{,} \PYG{l+m+mf}{0.0}\PYG{p}{,} \PYG{n}{arr}\PYG{p}{)}
        \PYG{n}{cs\PYGZus{}val} \PYG{o}{=} \PYG{n}{np}\PYG{o}{.}\PYG{n}{concatenate}\PYG{p}{(}\PYG{p}{[}\PYG{p}{[}\PYG{l+m+mf}{0.0}\PYG{p}{]}\PYG{p}{,} \PYG{n}{np}\PYG{o}{.}\PYG{n}{cumsum}\PYG{p}{(}\PYG{n}{filled}\PYG{p}{)}\PYG{p}{]}\PYG{p}{)}
        \PYG{n}{cs\PYGZus{}cnt} \PYG{o}{=} \PYG{n}{np}\PYG{o}{.}\PYG{n}{concatenate}\PYG{p}{(}\PYG{p}{[}\PYG{p}{[}\PYG{l+m+mf}{0.0}\PYG{p}{]}\PYG{p}{,} \PYG{n}{np}\PYG{o}{.}\PYG{n}{cumsum}\PYG{p}{(}\PYG{n}{valid}\PYG{p}{)}\PYG{p}{]}\PYG{p}{)}
        \PYG{n}{idx} \PYG{o}{=} \PYG{n}{np}\PYG{o}{.}\PYG{n}{arange}\PYG{p}{(}\PYG{n}{n}\PYG{p}{)}
        \PYG{k}{if} \PYG{n}{forward}\PYG{p}{:}
            \PYG{n}{starts} \PYG{o}{=} \PYG{n}{idx} \PYG{o}{+} \PYG{l+m+mi}{1}
            \PYG{n}{ends} \PYG{o}{=} \PYG{n}{np}\PYG{o}{.}\PYG{n}{minimum}\PYG{p}{(}\PYG{n}{n}\PYG{p}{,} \PYG{n}{idx} \PYG{o}{+} \PYG{l+m+mi}{1} \PYG{o}{+} \PYG{n}{window}\PYG{p}{)}
        \PYG{k}{else}\PYG{p}{:}
            \PYG{n}{starts} \PYG{o}{=} \PYG{n}{np}\PYG{o}{.}\PYG{n}{maximum}\PYG{p}{(}\PYG{l+m+mi}{0}\PYG{p}{,} \PYG{n}{idx} \PYG{o}{\PYGZhy{}} \PYG{n}{window}\PYG{p}{)}
            \PYG{n}{ends} \PYG{o}{=} \PYG{n}{idx}
        \PYG{n}{sum\PYGZus{}v} \PYG{o}{=} \PYG{n}{cs\PYGZus{}val}\PYG{p}{[}\PYG{n}{ends}\PYG{p}{]} \PYG{o}{\PYGZhy{}} \PYG{n}{cs\PYGZus{}val}\PYG{p}{[}\PYG{n}{starts}\PYG{p}{]}
        \PYG{n}{cnt\PYGZus{}v} \PYG{o}{=} \PYG{n}{cs\PYGZus{}cnt}\PYG{p}{[}\PYG{n}{ends}\PYG{p}{]} \PYG{o}{\PYGZhy{}} \PYG{n}{cs\PYGZus{}cnt}\PYG{p}{[}\PYG{n}{starts}\PYG{p}{]}
        \PYG{k}{with} \PYG{n}{np}\PYG{o}{.}\PYG{n}{errstate}\PYG{p}{(}\PYG{n}{invalid}\PYG{o}{=}\PYG{l+s+s1}{\PYGZsq{}}\PYG{l+s+s1}{ignore}\PYG{l+s+s1}{\PYGZsq{}}\PYG{p}{,} \PYG{n}{divide}\PYG{o}{=}\PYG{l+s+s1}{\PYGZsq{}}\PYG{l+s+s1}{ignore}\PYG{l+s+s1}{\PYGZsq{}}\PYG{p}{)}\PYG{p}{:}
            \PYG{k}{return} \PYG{n}{np}\PYG{o}{.}\PYG{n}{where}\PYG{p}{(}\PYG{n}{cnt\PYGZus{}v} \PYG{o}{\PYGZgt{}} \PYG{l+m+mi}{0}\PYG{p}{,} \PYG{n}{sum\PYGZus{}v} \PYG{o}{/} \PYG{n}{cnt\PYGZus{}v}\PYG{p}{,} \PYG{n}{np}\PYG{o}{.}\PYG{n}{nan}\PYG{p}{)}

    \PYG{n}{last\PYGZus{}x} \PYG{o}{=} \PYG{n}{rolling\PYGZus{}nanmean}\PYG{p}{(}\PYG{n}{x}\PYG{p}{,} \PYG{n}{num\PYGZus{}avg}\PYG{p}{)}
    \PYG{n}{last\PYGZus{}y} \PYG{o}{=} \PYG{n}{rolling\PYGZus{}nanmean}\PYG{p}{(}\PYG{n}{y}\PYG{p}{,} \PYG{n}{num\PYGZus{}avg}\PYG{p}{)}
    \PYG{n}{next\PYGZus{}x} \PYG{o}{=} \PYG{n}{rolling\PYGZus{}nanmean}\PYG{p}{(}\PYG{n}{x}\PYG{p}{,} \PYG{n}{num\PYGZus{}avg}\PYG{p}{,} \PYG{n}{forward}\PYG{o}{=}\PYG{k+kc}{True}\PYG{p}{)}
    \PYG{n}{next\PYGZus{}y} \PYG{o}{=} \PYG{n}{rolling\PYGZus{}nanmean}\PYG{p}{(}\PYG{n}{y}\PYG{p}{,} \PYG{n}{num\PYGZus{}avg}\PYG{p}{,} \PYG{n}{forward}\PYG{o}{=}\PYG{k+kc}{True}\PYG{p}{)}

    \PYG{n}{dist\PYGZus{}last} \PYG{o}{=} \PYG{n}{np}\PYG{o}{.}\PYG{n}{sqrt}\PYG{p}{(}\PYG{p}{(}\PYG{n}{x} \PYG{o}{\PYGZhy{}} \PYG{n}{last\PYGZus{}x}\PYG{p}{)} \PYG{o}{*}\PYG{o}{*} \PYG{l+m+mi}{2} \PYG{o}{+} \PYG{p}{(}\PYG{n}{y} \PYG{o}{\PYGZhy{}} \PYG{n}{last\PYGZus{}y}\PYG{p}{)} \PYG{o}{*}\PYG{o}{*} \PYG{l+m+mi}{2}\PYG{p}{)}
    \PYG{n}{dist\PYGZus{}next} \PYG{o}{=} \PYG{n}{np}\PYG{o}{.}\PYG{n}{sqrt}\PYG{p}{(}\PYG{p}{(}\PYG{n}{x} \PYG{o}{\PYGZhy{}} \PYG{n}{next\PYGZus{}x}\PYG{p}{)} \PYG{o}{*}\PYG{o}{*} \PYG{l+m+mi}{2} \PYG{o}{+} \PYG{p}{(}\PYG{n}{y} \PYG{o}{\PYGZhy{}} \PYG{n}{next\PYGZus{}y}\PYG{p}{)} \PYG{o}{*}\PYG{o}{*} \PYG{l+m+mi}{2}\PYG{p}{)}

    \PYG{n}{valid\PYGZus{}pts} \PYG{o}{=} \PYG{o}{\PYGZti{}}\PYG{n}{np}\PYG{o}{.}\PYG{n}{isnan}\PYG{p}{(}\PYG{n}{x}\PYG{p}{)}
    \PYG{n}{tele\PYGZus{}mask} \PYG{o}{=} \PYG{n}{valid\PYGZus{}pts} \PYG{o}{\PYGZam{}} \PYG{p}{(}
        \PYG{p}{(}\PYG{n}{dist\PYGZus{}last} \PYG{o}{\PYGZgt{}} \PYG{n}{tele\PYGZus{}dist\PYGZus{}threshold}\PYG{p}{)} \PYG{o}{|} \PYG{p}{(}\PYG{n}{dist\PYGZus{}next} \PYG{o}{\PYGZgt{}} \PYG{n}{tele\PYGZus{}dist\PYGZus{}threshold}\PYG{p}{)}
    \PYG{p}{)}
    \PYG{n}{tele\PYGZus{}mask}\PYG{p}{[}\PYG{p}{:}\PYG{n}{num\PYGZus{}avg}\PYG{p}{]} \PYG{o}{=} \PYG{k+kc}{False}
    \PYG{n}{tele\PYGZus{}mask}\PYG{p}{[}\PYG{n}{n} \PYG{o}{\PYGZhy{}} \PYG{n}{num\PYGZus{}avg}\PYG{p}{:}\PYG{p}{]} \PYG{o}{=} \PYG{k+kc}{False}

    \PYG{n}{filtered}\PYG{p}{[}\PYG{n}{tele\PYGZus{}mask}\PYG{p}{,} \PYG{l+m+mi}{1}\PYG{p}{:}\PYG{l+m+mi}{3}\PYG{p}{]} \PYG{o}{=} \PYG{n}{np}\PYG{o}{.}\PYG{n}{nan}
    \PYG{n}{tele\PYGZus{}count} \PYG{o}{=} \PYG{n+nb}{int}\PYG{p}{(}\PYG{n}{tele\PYGZus{}mask}\PYG{o}{.}\PYG{n}{sum}\PYG{p}{(}\PYG{p}{)}\PYG{p}{)}

    \PYG{n}{edge\PYGZus{}indices} \PYG{o}{=} \PYG{n+nb}{list}\PYG{p}{(}\PYG{n+nb}{range}\PYG{p}{(}\PYG{l+m+mi}{0}\PYG{p}{,} \PYG{n+nb}{min}\PYG{p}{(}\PYG{l+m+mi}{5}\PYG{p}{,} \PYG{n}{n} \PYG{o}{\PYGZhy{}} \PYG{l+m+mi}{1}\PYG{p}{)}\PYG{p}{)}\PYG{p}{)} \PYG{o}{+} \PYG{n+nb}{list}\PYG{p}{(}\PYG{n+nb}{range}\PYG{p}{(}\PYG{n+nb}{max}\PYG{p}{(}\PYG{l+m+mi}{0}\PYG{p}{,} \PYG{n}{n} \PYG{o}{\PYGZhy{}} \PYG{l+m+mi}{6}\PYG{p}{)}\PYG{p}{,} \PYG{n}{n} \PYG{o}{\PYGZhy{}} \PYG{l+m+mi}{1}\PYG{p}{)}\PYG{p}{)}
    \PYG{k}{for} \PYG{n}{idx} \PYG{o+ow}{in} \PYG{n}{edge\PYGZus{}indices}\PYG{p}{:}
        \PYG{k}{if} \PYG{n}{np}\PYG{o}{.}\PYG{n}{any}\PYG{p}{(}\PYG{n}{np}\PYG{o}{.}\PYG{n}{isnan}\PYG{p}{(}\PYG{n}{filtered}\PYG{p}{[}\PYG{n}{idx}\PYG{p}{,} \PYG{l+m+mi}{1}\PYG{p}{:}\PYG{l+m+mi}{3}\PYG{p}{]}\PYG{p}{)}\PYG{p}{)} \PYG{o+ow}{or} \PYG{n}{np}\PYG{o}{.}\PYG{n}{any}\PYG{p}{(}\PYG{n}{np}\PYG{o}{.}\PYG{n}{isnan}\PYG{p}{(}\PYG{n}{filtered}\PYG{p}{[}\PYG{n}{idx} \PYG{o}{+} \PYG{l+m+mi}{1}\PYG{p}{,} \PYG{l+m+mi}{1}\PYG{p}{:}\PYG{l+m+mi}{3}\PYG{p}{]}\PYG{p}{)}\PYG{p}{)}\PYG{p}{:}
            \PYG{k}{continue}
        \PYG{k}{if} \PYG{n}{euclidean}\PYG{p}{(}\PYG{n}{filtered}\PYG{p}{[}\PYG{n}{idx}\PYG{p}{,} \PYG{l+m+mi}{1}\PYG{p}{:}\PYG{l+m+mi}{3}\PYG{p}{]}\PYG{p}{,} \PYG{n}{filtered}\PYG{p}{[}\PYG{n}{idx} \PYG{o}{+} \PYG{l+m+mi}{1}\PYG{p}{,} \PYG{l+m+mi}{1}\PYG{p}{:}\PYG{l+m+mi}{3}\PYG{p}{]}\PYG{p}{)} \PYG{o}{\PYGZgt{}} \PYG{n}{tele\PYGZus{}dist\PYGZus{}threshold} \PYG{o}{/} \PYG{l+m+mi}{2}\PYG{p}{:}
            \PYG{n}{filtered}\PYG{p}{[}\PYG{n}{idx}\PYG{p}{,} \PYG{l+m+mi}{1}\PYG{p}{:}\PYG{l+m+mi}{3}\PYG{p}{]} \PYG{o}{=} \PYG{n}{np}\PYG{o}{.}\PYG{n}{nan}
            \PYG{n}{tele\PYGZus{}count} \PYG{o}{+}\PYG{o}{=} \PYG{l+m+mi}{1}

    \PYG{n+nb}{print}\PYG{p}{(}\PYG{l+s+sa}{f}\PYG{l+s+s1}{\PYGZsq{}}\PYG{l+s+si}{\PYGZob{}}\PYG{n}{tele\PYGZus{}count}\PYG{l+s+si}{\PYGZcb{}}\PYG{l+s+s1}{ points removed by the teleportation filter.}\PYG{l+s+s1}{\PYGZsq{}}\PYG{p}{)}
    \PYG{k}{return} \PYG{n}{filtered}


\PYG{k}{def}\PYG{+w}{ }\PYG{n+nf}{interpolate\PYGZus{}pos}\PYG{p}{(}\PYG{n}{array}\PYG{p}{,} \PYG{n}{inter\PYGZus{}dist\PYGZus{}threshold}\PYG{p}{)}\PYG{p}{:}
\PYG{+w}{    }\PYG{l+s+sd}{\PYGZdq{}\PYGZdq{}\PYGZdq{}}
\PYG{l+s+sd}{    Linearly interpolate fly position between NaN gaps,}
\PYG{l+s+sd}{    as long as the gap endpoints are within inter\PYGZus{}dist\PYGZus{}threshold.}

\PYG{l+s+sd}{    array: Nx3 array [time, x, y]}
\PYG{l+s+sd}{    \PYGZdq{}\PYGZdq{}\PYGZdq{}}
    \PYG{n}{result} \PYG{o}{=} \PYG{n}{array}\PYG{o}{.}\PYG{n}{copy}\PYG{p}{(}\PYG{p}{)}
    \PYG{n}{interp\PYGZus{}count} \PYG{o}{=} \PYG{l+m+mi}{0}

    \PYG{n}{x} \PYG{o}{=} \PYG{n}{result}\PYG{p}{[}\PYG{p}{:}\PYG{p}{,} \PYG{l+m+mi}{1}\PYG{p}{]}
    \PYG{n}{y} \PYG{o}{=} \PYG{n}{result}\PYG{p}{[}\PYG{p}{:}\PYG{p}{,} \PYG{l+m+mi}{2}\PYG{p}{]}
    \PYG{n}{nan\PYGZus{}mask} \PYG{o}{=} \PYG{n}{np}\PYG{o}{.}\PYG{n}{isnan}\PYG{p}{(}\PYG{n}{x}\PYG{p}{)}

    \PYG{k}{if} \PYG{o+ow}{not} \PYG{n}{nan\PYGZus{}mask}\PYG{o}{.}\PYG{n}{any}\PYG{p}{(}\PYG{p}{)}\PYG{p}{:}
        \PYG{n+nb}{print}\PYG{p}{(}\PYG{l+s+s1}{\PYGZsq{}}\PYG{l+s+s1}{0 points recovered through interpolation.}\PYG{l+s+s1}{\PYGZsq{}}\PYG{p}{)}
        \PYG{k}{return} \PYG{n}{result}

    \PYG{n}{padded} \PYG{o}{=} \PYG{n}{np}\PYG{o}{.}\PYG{n}{concatenate}\PYG{p}{(}\PYG{p}{[}\PYG{p}{[}\PYG{k+kc}{False}\PYG{p}{]}\PYG{p}{,} \PYG{n}{nan\PYGZus{}mask}\PYG{p}{,} \PYG{p}{[}\PYG{k+kc}{False}\PYG{p}{]}\PYG{p}{]}\PYG{p}{)}
    \PYG{n}{diff} \PYG{o}{=} \PYG{n}{np}\PYG{o}{.}\PYG{n}{diff}\PYG{p}{(}\PYG{n}{padded}\PYG{o}{.}\PYG{n}{astype}\PYG{p}{(}\PYG{n}{np}\PYG{o}{.}\PYG{n}{int8}\PYG{p}{)}\PYG{p}{)}
    \PYG{n}{starts} \PYG{o}{=} \PYG{n}{np}\PYG{o}{.}\PYG{n}{where}\PYG{p}{(}\PYG{n}{diff} \PYG{o}{==} \PYG{l+m+mi}{1}\PYG{p}{)}\PYG{p}{[}\PYG{l+m+mi}{0}\PYG{p}{]}
    \PYG{n}{ends} \PYG{o}{=} \PYG{n}{np}\PYG{o}{.}\PYG{n}{where}\PYG{p}{(}\PYG{n}{diff} \PYG{o}{==} \PYG{o}{\PYGZhy{}}\PYG{l+m+mi}{1}\PYG{p}{)}\PYG{p}{[}\PYG{l+m+mi}{0}\PYG{p}{]} \PYG{o}{\PYGZhy{}} \PYG{l+m+mi}{1}

    \PYG{k}{for} \PYG{n}{s}\PYG{p}{,} \PYG{n}{e} \PYG{o+ow}{in} \PYG{n+nb}{zip}\PYG{p}{(}\PYG{n}{starts}\PYG{p}{,} \PYG{n}{ends}\PYG{p}{)}\PYG{p}{:}
        \PYG{n}{last\PYGZus{}idx} \PYG{o}{=} \PYG{n}{s} \PYG{o}{\PYGZhy{}} \PYG{l+m+mi}{1}
        \PYG{n}{next\PYGZus{}idx} \PYG{o}{=} \PYG{n}{e} \PYG{o}{+} \PYG{l+m+mi}{1}
        \PYG{k}{if} \PYG{n}{last\PYGZus{}idx} \PYG{o}{\PYGZlt{}} \PYG{l+m+mi}{0} \PYG{o+ow}{or} \PYG{n}{next\PYGZus{}idx} \PYG{o}{\PYGZgt{}}\PYG{o}{=} \PYG{n+nb}{len}\PYG{p}{(}\PYG{n}{result}\PYG{p}{)}\PYG{p}{:}
            \PYG{k}{continue}
        \PYG{n}{last\PYGZus{}pt} \PYG{o}{=} \PYG{n}{result}\PYG{p}{[}\PYG{n}{last\PYGZus{}idx}\PYG{p}{,} \PYG{l+m+mi}{1}\PYG{p}{:}\PYG{l+m+mi}{3}\PYG{p}{]}
        \PYG{n}{next\PYGZus{}pt} \PYG{o}{=} \PYG{n}{result}\PYG{p}{[}\PYG{n}{next\PYGZus{}idx}\PYG{p}{,} \PYG{l+m+mi}{1}\PYG{p}{:}\PYG{l+m+mi}{3}\PYG{p}{]}
        \PYG{k}{if} \PYG{n}{np}\PYG{o}{.}\PYG{n}{any}\PYG{p}{(}\PYG{n}{np}\PYG{o}{.}\PYG{n}{isnan}\PYG{p}{(}\PYG{n}{last\PYGZus{}pt}\PYG{p}{)}\PYG{p}{)} \PYG{o+ow}{or} \PYG{n}{np}\PYG{o}{.}\PYG{n}{any}\PYG{p}{(}\PYG{n}{np}\PYG{o}{.}\PYG{n}{isnan}\PYG{p}{(}\PYG{n}{next\PYGZus{}pt}\PYG{p}{)}\PYG{p}{)}\PYG{p}{:}
            \PYG{k}{continue}
        \PYG{k}{if} \PYG{n}{euclidean}\PYG{p}{(}\PYG{n}{last\PYGZus{}pt}\PYG{p}{,} \PYG{n}{next\PYGZus{}pt}\PYG{p}{)} \PYG{o}{\PYGZlt{}}\PYG{o}{=} \PYG{n}{inter\PYGZus{}dist\PYGZus{}threshold}\PYG{p}{:}
            \PYG{n}{gap} \PYG{o}{=} \PYG{n}{next\PYGZus{}idx} \PYG{o}{\PYGZhy{}} \PYG{n}{last\PYGZus{}idx}
            \PYG{n}{fracs} \PYG{o}{=} \PYG{n}{np}\PYG{o}{.}\PYG{n}{arange}\PYG{p}{(}\PYG{l+m+mi}{1}\PYG{p}{,} \PYG{n}{gap}\PYG{p}{)} \PYG{o}{/} \PYG{n}{gap}
            \PYG{n}{result}\PYG{p}{[}\PYG{n}{s}\PYG{p}{:}\PYG{n}{e} \PYG{o}{+} \PYG{l+m+mi}{1}\PYG{p}{,} \PYG{l+m+mi}{1}\PYG{p}{]} \PYG{o}{=} \PYG{n}{last\PYGZus{}pt}\PYG{p}{[}\PYG{l+m+mi}{0}\PYG{p}{]} \PYG{o}{+} \PYG{n}{fracs} \PYG{o}{*} \PYG{p}{(}\PYG{n}{next\PYGZus{}pt}\PYG{p}{[}\PYG{l+m+mi}{0}\PYG{p}{]} \PYG{o}{\PYGZhy{}} \PYG{n}{last\PYGZus{}pt}\PYG{p}{[}\PYG{l+m+mi}{0}\PYG{p}{]}\PYG{p}{)}
            \PYG{n}{result}\PYG{p}{[}\PYG{n}{s}\PYG{p}{:}\PYG{n}{e} \PYG{o}{+} \PYG{l+m+mi}{1}\PYG{p}{,} \PYG{l+m+mi}{2}\PYG{p}{]} \PYG{o}{=} \PYG{n}{last\PYGZus{}pt}\PYG{p}{[}\PYG{l+m+mi}{1}\PYG{p}{]} \PYG{o}{+} \PYG{n}{fracs} \PYG{o}{*} \PYG{p}{(}\PYG{n}{next\PYGZus{}pt}\PYG{p}{[}\PYG{l+m+mi}{1}\PYG{p}{]} \PYG{o}{\PYGZhy{}} \PYG{n}{last\PYGZus{}pt}\PYG{p}{[}\PYG{l+m+mi}{1}\PYG{p}{]}\PYG{p}{)}
            \PYG{n}{interp\PYGZus{}count} \PYG{o}{+}\PYG{o}{=} \PYG{n}{e} \PYG{o}{\PYGZhy{}} \PYG{n}{s} \PYG{o}{+} \PYG{l+m+mi}{1}

    \PYG{n+nb}{print}\PYG{p}{(}\PYG{l+s+sa}{f}\PYG{l+s+s1}{\PYGZsq{}}\PYG{l+s+si}{\PYGZob{}}\PYG{n}{interp\PYGZus{}count}\PYG{l+s+si}{\PYGZcb{}}\PYG{l+s+s1}{ points recovered through interpolation.}\PYG{l+s+s1}{\PYGZsq{}}\PYG{p}{)}
    \PYG{k}{return} \PYG{n}{result}


\PYG{k}{def}\PYG{+w}{ }\PYG{n+nf}{boundary\PYGZus{}filter}\PYG{p}{(}\PYG{n}{array}\PYG{p}{,} \PYG{n}{diameter}\PYG{p}{)}\PYG{p}{:}
\PYG{+w}{    }\PYG{l+s+sd}{\PYGZdq{}\PYGZdq{}\PYGZdq{}}
\PYG{l+s+sd}{    Remove points that fall outside the circular dish.}
\PYG{l+s+sd}{    \PYGZdq{}\PYGZdq{}\PYGZdq{}}
    \PYG{n}{filtered} \PYG{o}{=} \PYG{n}{array}\PYG{o}{.}\PYG{n}{copy}\PYG{p}{(}\PYG{p}{)}
    \PYG{n}{cx}\PYG{p}{,} \PYG{n}{cy} \PYG{o}{=} \PYG{n}{diameter} \PYG{o}{/} \PYG{l+m+mi}{2}\PYG{p}{,} \PYG{n}{diameter} \PYG{o}{/} \PYG{l+m+mi}{2}
    \PYG{n}{radius} \PYG{o}{=} \PYG{n}{diameter} \PYG{o}{/} \PYG{l+m+mi}{2}
    \PYG{n}{x} \PYG{o}{=} \PYG{n}{filtered}\PYG{p}{[}\PYG{p}{:}\PYG{p}{,} \PYG{l+m+mi}{1}\PYG{p}{]}
    \PYG{n}{y} \PYG{o}{=} \PYG{n}{filtered}\PYG{p}{[}\PYG{p}{:}\PYG{p}{,} \PYG{l+m+mi}{2}\PYG{p}{]}
    \PYG{n}{dist\PYGZus{}from\PYGZus{}center} \PYG{o}{=} \PYG{n}{np}\PYG{o}{.}\PYG{n}{sqrt}\PYG{p}{(}\PYG{p}{(}\PYG{n}{x} \PYG{o}{\PYGZhy{}} \PYG{n}{cx}\PYG{p}{)} \PYG{o}{*}\PYG{o}{*} \PYG{l+m+mi}{2} \PYG{o}{+} \PYG{p}{(}\PYG{n}{y} \PYG{o}{\PYGZhy{}} \PYG{n}{cy}\PYG{p}{)} \PYG{o}{*}\PYG{o}{*} \PYG{l+m+mi}{2}\PYG{p}{)}
    \PYG{n}{outside} \PYG{o}{=} \PYG{n}{dist\PYGZus{}from\PYGZus{}center} \PYG{o}{\PYGZgt{}} \PYG{n}{radius}
    \PYG{n}{outside} \PYG{o}{=} \PYG{n}{outside} \PYG{o}{\PYGZam{}} \PYG{o}{\PYGZti{}}\PYG{n}{np}\PYG{o}{.}\PYG{n}{isnan}\PYG{p}{(}\PYG{n}{x}\PYG{p}{)}
    \PYG{n}{filtered}\PYG{p}{[}\PYG{n}{outside}\PYG{p}{,} \PYG{l+m+mi}{1}\PYG{p}{:}\PYG{l+m+mi}{3}\PYG{p}{]} \PYG{o}{=} \PYG{n}{np}\PYG{o}{.}\PYG{n}{nan}
    \PYG{n}{count} \PYG{o}{=} \PYG{n}{np}\PYG{o}{.}\PYG{n}{sum}\PYG{p}{(}\PYG{n}{outside}\PYG{p}{)}
    \PYG{k}{if} \PYG{n}{count} \PYG{o}{\PYGZgt{}} \PYG{l+m+mi}{0}\PYG{p}{:}
        \PYG{n+nb}{print}\PYG{p}{(}\PYG{l+s+sa}{f}\PYG{l+s+s1}{\PYGZsq{}}\PYG{l+s+s1}{  }\PYG{l+s+si}{\PYGZob{}}\PYG{n}{count}\PYG{l+s+si}{\PYGZcb{}}\PYG{l+s+s1}{ points removed by the boundary filter (outside dish).}\PYG{l+s+s1}{\PYGZsq{}}\PYG{p}{)}
    \PYG{k}{return} \PYG{n}{filtered}


\PYG{k}{def}\PYG{+w}{ }\PYG{n+nf}{displacement\PYGZus{}filter}\PYG{p}{(}\PYG{n}{array}\PYG{p}{,} \PYG{n}{multiplier}\PYG{o}{=}\PYG{l+m+mi}{5}\PYG{p}{)}\PYG{p}{:}
\PYG{+w}{    }\PYG{l+s+sd}{\PYGZdq{}\PYGZdq{}\PYGZdq{}}
\PYG{l+s+sd}{    Remove points where the frame\PYGZhy{}to\PYGZhy{}frame displacement is much larger}
\PYG{l+s+sd}{    than the median displacement. This catches sudden jumps that the}
\PYG{l+s+sd}{    teleport filter might miss (e.g., runs of bad frames).}
\PYG{l+s+sd}{    \PYGZdq{}\PYGZdq{}\PYGZdq{}}
    \PYG{n}{filtered} \PYG{o}{=} \PYG{n}{array}\PYG{o}{.}\PYG{n}{copy}\PYG{p}{(}\PYG{p}{)}
    \PYG{n}{x} \PYG{o}{=} \PYG{n}{filtered}\PYG{p}{[}\PYG{p}{:}\PYG{p}{,} \PYG{l+m+mi}{1}\PYG{p}{]}
    \PYG{n}{y} \PYG{o}{=} \PYG{n}{filtered}\PYG{p}{[}\PYG{p}{:}\PYG{p}{,} \PYG{l+m+mi}{2}\PYG{p}{]}

    \PYG{n}{dx} \PYG{o}{=} \PYG{n}{np}\PYG{o}{.}\PYG{n}{diff}\PYG{p}{(}\PYG{n}{x}\PYG{p}{)}
    \PYG{n}{dy} \PYG{o}{=} \PYG{n}{np}\PYG{o}{.}\PYG{n}{diff}\PYG{p}{(}\PYG{n}{y}\PYG{p}{)}
    \PYG{n}{displacements} \PYG{o}{=} \PYG{n}{np}\PYG{o}{.}\PYG{n}{sqrt}\PYG{p}{(}\PYG{n}{dx} \PYG{o}{*}\PYG{o}{*} \PYG{l+m+mi}{2} \PYG{o}{+} \PYG{n}{dy} \PYG{o}{*}\PYG{o}{*} \PYG{l+m+mi}{2}\PYG{p}{)}

    \PYG{n}{valid\PYGZus{}disp} \PYG{o}{=} \PYG{n}{displacements}\PYG{p}{[}\PYG{o}{\PYGZti{}}\PYG{n}{np}\PYG{o}{.}\PYG{n}{isnan}\PYG{p}{(}\PYG{n}{displacements}\PYG{p}{)}\PYG{p}{]}
    \PYG{k}{if} \PYG{n+nb}{len}\PYG{p}{(}\PYG{n}{valid\PYGZus{}disp}\PYG{p}{)} \PYG{o}{==} \PYG{l+m+mi}{0}\PYG{p}{:}
        \PYG{k}{return} \PYG{n}{filtered}
    \PYG{n}{median\PYGZus{}disp} \PYG{o}{=} \PYG{n}{np}\PYG{o}{.}\PYG{n}{nanmedian}\PYG{p}{(}\PYG{n}{valid\PYGZus{}disp}\PYG{p}{)}

    \PYG{k}{if} \PYG{n}{median\PYGZus{}disp} \PYG{o}{==} \PYG{l+m+mi}{0}\PYG{p}{:}
        \PYG{k}{return} \PYG{n}{filtered}

    \PYG{n}{threshold} \PYG{o}{=} \PYG{n}{median\PYGZus{}disp} \PYG{o}{*} \PYG{n}{multiplier}
    \PYG{n}{count} \PYG{o}{=} \PYG{l+m+mi}{0}

    \PYG{k}{for} \PYG{n}{i} \PYG{o+ow}{in} \PYG{n+nb}{range}\PYG{p}{(}\PYG{n+nb}{len}\PYG{p}{(}\PYG{n}{displacements}\PYG{p}{)}\PYG{p}{)}\PYG{p}{:}
        \PYG{k}{if} \PYG{n}{np}\PYG{o}{.}\PYG{n}{isnan}\PYG{p}{(}\PYG{n}{displacements}\PYG{p}{[}\PYG{n}{i}\PYG{p}{]}\PYG{p}{)}\PYG{p}{:}
            \PYG{k}{continue}
        \PYG{k}{if} \PYG{n}{displacements}\PYG{p}{[}\PYG{n}{i}\PYG{p}{]} \PYG{o}{\PYGZgt{}} \PYG{n}{threshold}\PYG{p}{:}
            \PYG{n}{filtered}\PYG{p}{[}\PYG{n}{i} \PYG{o}{+} \PYG{l+m+mi}{1}\PYG{p}{,} \PYG{l+m+mi}{1}\PYG{p}{:}\PYG{l+m+mi}{3}\PYG{p}{]} \PYG{o}{=} \PYG{n}{np}\PYG{o}{.}\PYG{n}{nan}
            \PYG{n}{count} \PYG{o}{+}\PYG{o}{=} \PYG{l+m+mi}{1}

    \PYG{k}{if} \PYG{n}{count} \PYG{o}{\PYGZgt{}} \PYG{l+m+mi}{0}\PYG{p}{:}
        \PYG{n+nb}{print}\PYG{p}{(}\PYG{l+s+sa}{f}\PYG{l+s+s1}{\PYGZsq{}}\PYG{l+s+s1}{  }\PYG{l+s+si}{\PYGZob{}}\PYG{n}{count}\PYG{l+s+si}{\PYGZcb{}}\PYG{l+s+s1}{ points removed by the displacement filter (\PYGZgt{}}\PYG{l+s+si}{\PYGZob{}}\PYG{n}{multiplier}\PYG{l+s+si}{\PYGZcb{}}\PYG{l+s+s1}{x median jump).}\PYG{l+s+s1}{\PYGZsq{}}\PYG{p}{)}
    \PYG{k}{return} \PYG{n}{filtered}


\PYG{k}{def}\PYG{+w}{ }\PYG{n+nf}{process\PYGZus{}video}\PYG{p}{(}\PYG{n}{video\PYGZus{}path}\PYG{p}{,} \PYG{n}{roi}\PYG{p}{,} \PYG{n}{diameter}\PYG{p}{,} \PYG{n}{frame\PYGZus{}rate}\PYG{p}{,} \PYG{n}{search\PYGZus{}size}\PYG{p}{,} \PYG{n}{per\PYGZus{}pixel\PYGZus{}threshold}\PYG{p}{)}\PYG{p}{:}
\PYG{+w}{    }\PYG{l+s+sd}{\PYGZdq{}\PYGZdq{}\PYGZdq{}}
\PYG{l+s+sd}{    Process a single fly video and return the corrected position array.}
\PYG{l+s+sd}{    Returns (Nx3 array [time\PYGZus{}s, x\PYGZus{}cm, y\PYGZus{}cm], frame\PYGZus{}rate).}
\PYG{l+s+sd}{    \PYGZdq{}\PYGZdq{}\PYGZdq{}}
    \PYG{n}{height} \PYG{o}{=} \PYG{n}{diameter}
    \PYG{n}{width} \PYG{o}{=} \PYG{n}{diameter}
    \PYG{n}{roi\PYGZus{}x}\PYG{p}{,} \PYG{n}{roi\PYGZus{}y}\PYG{p}{,} \PYG{n}{roi\PYGZus{}w}\PYG{p}{,} \PYG{n}{roi\PYGZus{}h} \PYG{o}{=} \PYG{n}{roi}

    \PYG{n}{cap} \PYG{o}{=} \PYG{n}{cv2}\PYG{o}{.}\PYG{n}{VideoCapture}\PYG{p}{(}\PYG{n}{video\PYGZus{}path}\PYG{p}{)}
    \PYG{n}{fps} \PYG{o}{=} \PYG{n}{cap}\PYG{o}{.}\PYG{n}{get}\PYG{p}{(}\PYG{n}{cv2}\PYG{o}{.}\PYG{n}{CAP\PYGZus{}PROP\PYGZus{}FPS}\PYG{p}{)}
    \PYG{n}{total\PYGZus{}frames} \PYG{o}{=} \PYG{n+nb}{int}\PYG{p}{(}\PYG{n}{cap}\PYG{o}{.}\PYG{n}{get}\PYG{p}{(}\PYG{n}{cv2}\PYG{o}{.}\PYG{n}{CAP\PYGZus{}PROP\PYGZus{}FRAME\PYGZus{}COUNT}\PYG{p}{)}\PYG{p}{)}
    \PYG{n}{nfrm} \PYG{o}{=} \PYG{n}{total\PYGZus{}frames} \PYG{o}{\PYGZhy{}} \PYG{l+m+mi}{1}

    \PYG{c+c1}{\PYGZsh{} Auto\PYGZhy{}detect frame rate from video if not specified}
    \PYG{k}{if} \PYG{n}{frame\PYGZus{}rate} \PYG{o+ow}{is} \PYG{k+kc}{None}\PYG{p}{:}
        \PYG{k}{if} \PYG{n}{fps} \PYG{o}{\PYGZgt{}} \PYG{l+m+mi}{0}\PYG{p}{:}
            \PYG{n}{frame\PYGZus{}rate} \PYG{o}{=} \PYG{n+nb}{round}\PYG{p}{(}\PYG{n}{fps}\PYG{p}{)}
            \PYG{n+nb}{print}\PYG{p}{(}\PYG{l+s+sa}{f}\PYG{l+s+s2}{\PYGZdq{}}\PYG{l+s+se}{\PYGZbs{}n}\PYG{l+s+s2}{Processing: }\PYG{l+s+si}{\PYGZob{}}\PYG{n}{os}\PYG{o}{.}\PYG{n}{path}\PYG{o}{.}\PYG{n}{basename}\PYG{p}{(}\PYG{n}{video\PYGZus{}path}\PYG{p}{)}\PYG{l+s+si}{\PYGZcb{}}\PYG{l+s+s2}{\PYGZdq{}}\PYG{p}{)}
            \PYG{n+nb}{print}\PYG{p}{(}\PYG{l+s+sa}{f}\PYG{l+s+s2}{\PYGZdq{}}\PYG{l+s+s2}{  Auto\PYGZhy{}detected frame rate: }\PYG{l+s+si}{\PYGZob{}}\PYG{n}{frame\PYGZus{}rate}\PYG{l+s+si}{\PYGZcb{}}\PYG{l+s+s2}{ fps (}\PYG{l+s+si}{\PYGZob{}}\PYG{n}{total\PYGZus{}frames}\PYG{l+s+si}{\PYGZcb{}}\PYG{l+s+s2}{ frames)}\PYG{l+s+s2}{\PYGZdq{}}\PYG{p}{)}
        \PYG{k}{else}\PYG{p}{:}
            \PYG{n}{frame\PYGZus{}rate} \PYG{o}{=} \PYG{l+m+mi}{30}
            \PYG{n+nb}{print}\PYG{p}{(}\PYG{l+s+sa}{f}\PYG{l+s+s2}{\PYGZdq{}}\PYG{l+s+se}{\PYGZbs{}n}\PYG{l+s+s2}{Processing: }\PYG{l+s+si}{\PYGZob{}}\PYG{n}{os}\PYG{o}{.}\PYG{n}{path}\PYG{o}{.}\PYG{n}{basename}\PYG{p}{(}\PYG{n}{video\PYGZus{}path}\PYG{p}{)}\PYG{l+s+si}{\PYGZcb{}}\PYG{l+s+s2}{\PYGZdq{}}\PYG{p}{)}
            \PYG{n+nb}{print}\PYG{p}{(}\PYG{l+s+sa}{f}\PYG{l+s+s2}{\PYGZdq{}}\PYG{l+s+s2}{  WARNING: Could not detect frame rate, using default: }\PYG{l+s+si}{\PYGZob{}}\PYG{n}{frame\PYGZus{}rate}\PYG{l+s+si}{\PYGZcb{}}\PYG{l+s+s2}{ fps}\PYG{l+s+s2}{\PYGZdq{}}\PYG{p}{)}
            \PYG{n+nb}{print}\PYG{p}{(}\PYG{l+s+sa}{f}\PYG{l+s+s2}{\PYGZdq{}}\PYG{l+s+s2}{    Override by setting frame\PYGZus{}rate = 30 (or your value) in the Set Parameters cell.}\PYG{l+s+s2}{\PYGZdq{}}\PYG{p}{)}
    \PYG{k}{else}\PYG{p}{:}
        \PYG{n+nb}{print}\PYG{p}{(}\PYG{l+s+sa}{f}\PYG{l+s+s2}{\PYGZdq{}}\PYG{l+s+se}{\PYGZbs{}n}\PYG{l+s+s2}{Processing: }\PYG{l+s+si}{\PYGZob{}}\PYG{n}{os}\PYG{o}{.}\PYG{n}{path}\PYG{o}{.}\PYG{n}{basename}\PYG{p}{(}\PYG{n}{video\PYGZus{}path}\PYG{p}{)}\PYG{l+s+si}{\PYGZcb{}}\PYG{l+s+s2}{\PYGZdq{}}\PYG{p}{)}
        \PYG{n+nb}{print}\PYG{p}{(}\PYG{l+s+sa}{f}\PYG{l+s+s2}{\PYGZdq{}}\PYG{l+s+s2}{  Video FPS: }\PYG{l+s+si}{\PYGZob{}}\PYG{n}{fps}\PYG{l+s+si}{\PYGZcb{}}\PYG{l+s+s2}{, using override: }\PYG{l+s+si}{\PYGZob{}}\PYG{n}{frame\PYGZus{}rate}\PYG{l+s+si}{\PYGZcb{}}\PYG{l+s+s2}{ fps (}\PYG{l+s+si}{\PYGZob{}}\PYG{n}{total\PYGZus{}frames}\PYG{l+s+si}{\PYGZcb{}}\PYG{l+s+s2}{ frames)}\PYG{l+s+s2}{\PYGZdq{}}\PYG{p}{)}

    \PYG{c+c1}{\PYGZsh{} Sanity check: warn if video duration seems unusual}
    \PYG{k}{if} \PYG{n}{frame\PYGZus{}rate} \PYG{o}{\PYGZgt{}} \PYG{l+m+mi}{0}\PYG{p}{:}
        \PYG{n}{duration} \PYG{o}{=} \PYG{n}{total\PYGZus{}frames} \PYG{o}{/} \PYG{n}{frame\PYGZus{}rate}
        \PYG{n+nb}{print}\PYG{p}{(}\PYG{l+s+sa}{f}\PYG{l+s+s2}{\PYGZdq{}}\PYG{l+s+s2}{  Estimated duration: }\PYG{l+s+si}{\PYGZob{}}\PYG{n}{duration}\PYG{l+s+si}{:}\PYG{l+s+s2}{.1f}\PYG{l+s+si}{\PYGZcb{}}\PYG{l+s+s2}{ seconds}\PYG{l+s+s2}{\PYGZdq{}}\PYG{p}{)}
        \PYG{k}{if} \PYG{n}{duration} \PYG{o}{\PYGZlt{}} \PYG{l+m+mi}{50} \PYG{o+ow}{or} \PYG{n}{duration} \PYG{o}{\PYGZgt{}} \PYG{l+m+mi}{70}\PYG{p}{:}
            \PYG{n+nb}{print}\PYG{p}{(}\PYG{l+s+sa}{f}\PYG{l+s+s2}{\PYGZdq{}}\PYG{l+s+s2}{  WARNING: Expected \PYGZti{}60s video but got }\PYG{l+s+si}{\PYGZob{}}\PYG{n}{duration}\PYG{l+s+si}{:}\PYG{l+s+s2}{.1f}\PYG{l+s+si}{\PYGZcb{}}\PYG{l+s+s2}{s.}\PYG{l+s+s2}{\PYGZdq{}}\PYG{p}{)}
            \PYG{n+nb}{print}\PYG{p}{(}\PYG{l+s+sa}{f}\PYG{l+s+s2}{\PYGZdq{}}\PYG{l+s+s2}{    Override by setting frame\PYGZus{}rate = 30 (or your value) in the Set Parameters cell.}\PYG{l+s+s2}{\PYGZdq{}}\PYG{p}{)}

    \PYG{c+c1}{\PYGZsh{} \PYGZhy{}\PYGZhy{}\PYGZhy{} Create background from 100 random frames \PYGZhy{}\PYGZhy{}\PYGZhy{}}
    \PYG{n+nb}{print}\PYG{p}{(}\PYG{l+s+s2}{\PYGZdq{}}\PYG{l+s+s2}{  Calculating background...}\PYG{l+s+s2}{\PYGZdq{}}\PYG{p}{)}
    \PYG{n}{bg\PYGZus{}number} \PYG{o}{=} \PYG{n+nb}{min}\PYG{p}{(}\PYG{l+m+mi}{100}\PYG{p}{,} \PYG{n}{nfrm}\PYG{p}{)}
    \PYG{n}{bg\PYGZus{}indices} \PYG{o}{=} \PYG{n+nb}{sorted}\PYG{p}{(}\PYG{n}{np}\PYG{o}{.}\PYG{n}{random}\PYG{o}{.}\PYG{n}{choice}\PYG{p}{(}\PYG{n}{nfrm}\PYG{p}{,} \PYG{n}{bg\PYGZus{}number}\PYG{p}{,} \PYG{n}{replace}\PYG{o}{=}\PYG{k+kc}{False}\PYG{p}{)}\PYG{p}{)}

    \PYG{n}{cap}\PYG{o}{.}\PYG{n}{set}\PYG{p}{(}\PYG{n}{cv2}\PYG{o}{.}\PYG{n}{CAP\PYGZus{}PROP\PYGZus{}POS\PYGZus{}FRAMES}\PYG{p}{,} \PYG{l+m+mi}{0}\PYG{p}{)}
    \PYG{n}{ret}\PYG{p}{,} \PYG{n}{sample} \PYG{o}{=} \PYG{n}{cap}\PYG{o}{.}\PYG{n}{read}\PYG{p}{(}\PYG{p}{)}
    \PYG{n}{gray\PYGZus{}sample} \PYG{o}{=} \PYG{n}{cv2}\PYG{o}{.}\PYG{n}{cvtColor}\PYG{p}{(}\PYG{n}{sample}\PYG{p}{,} \PYG{n}{cv2}\PYG{o}{.}\PYG{n}{COLOR\PYGZus{}BGR2GRAY}\PYG{p}{)}
    \PYG{n}{bg\PYGZus{}array} \PYG{o}{=} \PYG{n}{np}\PYG{o}{.}\PYG{n}{zeros}\PYG{p}{(}\PYG{p}{(}\PYG{o}{*}\PYG{n}{gray\PYGZus{}sample}\PYG{o}{.}\PYG{n}{shape}\PYG{p}{,} \PYG{n}{bg\PYGZus{}number}\PYG{p}{)}\PYG{p}{,} \PYG{n}{dtype}\PYG{o}{=}\PYG{n}{np}\PYG{o}{.}\PYG{n}{uint8}\PYG{p}{)}

    \PYG{k}{for} \PYG{n}{idx}\PYG{p}{,} \PYG{n}{frame\PYGZus{}num} \PYG{o+ow}{in} \PYG{n+nb}{enumerate}\PYG{p}{(}\PYG{n}{bg\PYGZus{}indices}\PYG{p}{)}\PYG{p}{:}
        \PYG{n}{cap}\PYG{o}{.}\PYG{n}{set}\PYG{p}{(}\PYG{n}{cv2}\PYG{o}{.}\PYG{n}{CAP\PYGZus{}PROP\PYGZus{}POS\PYGZus{}FRAMES}\PYG{p}{,} \PYG{n}{frame\PYGZus{}num}\PYG{p}{)}
        \PYG{n}{ret}\PYG{p}{,} \PYG{n}{frame} \PYG{o}{=} \PYG{n}{cap}\PYG{o}{.}\PYG{n}{read}\PYG{p}{(}\PYG{p}{)}
        \PYG{k}{if} \PYG{n}{ret}\PYG{p}{:}
            \PYG{n}{bg\PYGZus{}array}\PYG{p}{[}\PYG{p}{:}\PYG{p}{,} \PYG{p}{:}\PYG{p}{,} \PYG{n}{idx}\PYG{p}{]} \PYG{o}{=} \PYG{n}{cv2}\PYG{o}{.}\PYG{n}{cvtColor}\PYG{p}{(}\PYG{n}{frame}\PYG{p}{,} \PYG{n}{cv2}\PYG{o}{.}\PYG{n}{COLOR\PYGZus{}BGR2GRAY}\PYG{p}{)}

    \PYG{n}{background} \PYG{o}{=} \PYG{n}{np}\PYG{o}{.}\PYG{n}{mean}\PYG{p}{(}\PYG{n}{bg\PYGZus{}array}\PYG{p}{,} \PYG{n}{axis}\PYG{o}{=}\PYG{l+m+mi}{2}\PYG{p}{)}\PYG{o}{.}\PYG{n}{astype}\PYG{p}{(}\PYG{n}{np}\PYG{o}{.}\PYG{n}{uint8}\PYG{p}{)}
    \PYG{n}{bg\PYGZus{}scale} \PYG{o}{=} \PYG{l+m+mf}{256.0} \PYG{o}{/} \PYG{p}{(}\PYG{n}{background}\PYG{o}{.}\PYG{n}{astype}\PYG{p}{(}\PYG{n}{np}\PYG{o}{.}\PYG{n}{float64}\PYG{p}{)} \PYG{o}{+} \PYG{l+m+mi}{1}\PYG{p}{)}
    \PYG{c+c1}{\PYGZsh{} Pre\PYGZhy{}crop bg\PYGZus{}scale to ROI so per\PYGZhy{}frame math runs on ROI pixels only}
    \PYG{n}{bg\PYGZus{}scale\PYGZus{}roi} \PYG{o}{=} \PYG{n}{bg\PYGZus{}scale}\PYG{p}{[}\PYG{n}{roi\PYGZus{}y}\PYG{p}{:}\PYG{n}{roi\PYGZus{}y} \PYG{o}{+} \PYG{n}{roi\PYGZus{}h}\PYG{p}{,} \PYG{n}{roi\PYGZus{}x}\PYG{p}{:}\PYG{n}{roi\PYGZus{}x} \PYG{o}{+} \PYG{n}{roi\PYGZus{}w}\PYG{p}{]}

    \PYG{c+c1}{\PYGZsh{} \PYGZhy{}\PYGZhy{}\PYGZhy{} Process each frame \PYGZhy{}\PYGZhy{}\PYGZhy{}}
    \PYG{n+nb}{print}\PYG{p}{(}\PYG{l+s+s2}{\PYGZdq{}}\PYG{l+s+s2}{  Tracking fly positions...}\PYG{l+s+s2}{\PYGZdq{}}\PYG{p}{)}
    \PYG{n}{threshold} \PYG{o}{=} \PYG{p}{(}\PYG{n}{search\PYGZus{}size} \PYG{o}{*}\PYG{o}{*} \PYG{l+m+mi}{2}\PYG{p}{)} \PYG{o}{*} \PYG{n}{per\PYGZus{}pixel\PYGZus{}threshold}
    \PYG{n}{half\PYGZus{}search} \PYG{o}{=} \PYG{n+nb}{round}\PYG{p}{(}\PYG{n}{search\PYGZus{}size} \PYG{o}{/} \PYG{l+m+mi}{2}\PYG{p}{)}

    \PYG{n}{pos\PYGZus{}array} \PYG{o}{=} \PYG{n}{np}\PYG{o}{.}\PYG{n}{zeros}\PYG{p}{(}\PYG{p}{(}\PYG{n}{nfrm}\PYG{p}{,} \PYG{l+m+mi}{3}\PYG{p}{)}\PYG{p}{)}

    \PYG{c+c1}{\PYGZsh{} Seek once; sequential cap.read() is far faster than per\PYGZhy{}frame seeking}
    \PYG{c+c1}{\PYGZsh{} (compressed formats must decode from the last keyframe on every seek)}
    \PYG{n}{cap}\PYG{o}{.}\PYG{n}{set}\PYG{p}{(}\PYG{n}{cv2}\PYG{o}{.}\PYG{n}{CAP\PYGZus{}PROP\PYGZus{}POS\PYGZus{}FRAMES}\PYG{p}{,} \PYG{l+m+mi}{0}\PYG{p}{)}
    \PYG{k}{for} \PYG{n}{nofr} \PYG{o+ow}{in} \PYG{n+nb}{range}\PYG{p}{(}\PYG{n}{nfrm}\PYG{p}{)}\PYG{p}{:}
        \PYG{n}{ret}\PYG{p}{,} \PYG{n}{frame} \PYG{o}{=} \PYG{n}{cap}\PYG{o}{.}\PYG{n}{read}\PYG{p}{(}\PYG{p}{)}
        \PYG{k}{if} \PYG{o+ow}{not} \PYG{n}{ret}\PYG{p}{:}
            \PYG{n}{pos\PYGZus{}array}\PYG{p}{[}\PYG{n}{nofr}\PYG{p}{]} \PYG{o}{=} \PYG{p}{[}\PYG{n}{nofr}\PYG{p}{,} \PYG{n}{np}\PYG{o}{.}\PYG{n}{nan}\PYG{p}{,} \PYG{n}{np}\PYG{o}{.}\PYG{n}{nan}\PYG{p}{]}
            \PYG{k}{continue}

        \PYG{c+c1}{\PYGZsh{} Crop to ROI before float conversion so background subtraction}
        \PYG{c+c1}{\PYGZsh{} operates on ROI pixels only, not the full frame}
        \PYG{n}{frame\PYGZus{}roi} \PYG{o}{=} \PYG{n}{cv2}\PYG{o}{.}\PYG{n}{cvtColor}\PYG{p}{(}\PYG{n}{frame}\PYG{p}{,} \PYG{n}{cv2}\PYG{o}{.}\PYG{n}{COLOR\PYGZus{}BGR2GRAY}\PYG{p}{)}\PYG{p}{[}
            \PYG{n}{roi\PYGZus{}y}\PYG{p}{:}\PYG{n}{roi\PYGZus{}y} \PYG{o}{+} \PYG{n}{roi\PYGZus{}h}\PYG{p}{,} \PYG{n}{roi\PYGZus{}x}\PYG{p}{:}\PYG{n}{roi\PYGZus{}x} \PYG{o}{+} \PYG{n}{roi\PYGZus{}w}
        \PYG{p}{]}\PYG{o}{.}\PYG{n}{astype}\PYG{p}{(}\PYG{n}{np}\PYG{o}{.}\PYG{n}{float64}\PYG{p}{)}
        \PYG{n}{frame\PYGZus{}div} \PYG{o}{=} \PYG{n}{np}\PYG{o}{.}\PYG{n}{clip}\PYG{p}{(}\PYG{n}{frame\PYGZus{}roi} \PYG{o}{*} \PYG{n}{bg\PYGZus{}scale\PYGZus{}roi}\PYG{p}{,} \PYG{l+m+mi}{0}\PYG{p}{,} \PYG{l+m+mi}{255}\PYG{p}{)}\PYG{o}{.}\PYG{n}{astype}\PYG{p}{(}\PYG{n}{np}\PYG{o}{.}\PYG{n}{uint8}\PYG{p}{)}

        \PYG{c+c1}{\PYGZsh{} Find fly}
        \PYG{n}{fx}\PYG{p}{,} \PYG{n}{fy} \PYG{o}{=} \PYG{n}{fly\PYGZus{}finder}\PYG{p}{(}\PYG{n}{frame\PYGZus{}div}\PYG{p}{,} \PYG{n}{half\PYGZus{}search}\PYG{p}{,} \PYG{n}{threshold}\PYG{p}{,} \PYG{n}{flip}\PYG{o}{=}\PYG{k+kc}{True}\PYG{p}{)}
        \PYG{n}{pos\PYGZus{}array}\PYG{p}{[}\PYG{n}{nofr}\PYG{p}{]} \PYG{o}{=} \PYG{p}{[}\PYG{n}{nofr}\PYG{p}{,} \PYG{n}{fx}\PYG{p}{,} \PYG{n}{fy}\PYG{p}{]}

        \PYG{c+c1}{\PYGZsh{} Progress update every 10\PYGZpc{}}
        \PYG{k}{if} \PYG{p}{(}\PYG{n}{nofr} \PYG{o}{+} \PYG{l+m+mi}{1}\PYG{p}{)} \PYG{o}{\PYGZpc{}} \PYG{n+nb}{max}\PYG{p}{(}\PYG{l+m+mi}{1}\PYG{p}{,} \PYG{n}{nfrm} \PYG{o}{/}\PYG{o}{/} \PYG{l+m+mi}{10}\PYG{p}{)} \PYG{o}{==} \PYG{l+m+mi}{0}\PYG{p}{:}
            \PYG{n}{pct} \PYG{o}{=} \PYG{p}{(}\PYG{n}{nofr} \PYG{o}{+} \PYG{l+m+mi}{1}\PYG{p}{)} \PYG{o}{/} \PYG{n}{nfrm} \PYG{o}{*} \PYG{l+m+mi}{100}
            \PYG{n+nb}{print}\PYG{p}{(}\PYG{l+s+sa}{f}\PYG{l+s+s2}{\PYGZdq{}}\PYG{l+s+s2}{    }\PYG{l+s+si}{\PYGZob{}}\PYG{n}{pct}\PYG{l+s+si}{:}\PYG{l+s+s2}{.0f}\PYG{l+s+si}{\PYGZcb{}}\PYG{l+s+s2}{\PYGZpc{} complete (}\PYG{l+s+si}{\PYGZob{}}\PYG{n}{nofr}\PYG{+w}{ }\PYG{o}{+}\PYG{+w}{ }\PYG{l+m+mi}{1}\PYG{l+s+si}{\PYGZcb{}}\PYG{l+s+s2}{/}\PYG{l+s+si}{\PYGZob{}}\PYG{n}{nfrm}\PYG{l+s+si}{\PYGZcb{}}\PYG{l+s+s2}{ frames)}\PYG{l+s+s2}{\PYGZdq{}}\PYG{p}{)}

    \PYG{n}{cap}\PYG{o}{.}\PYG{n}{release}\PYG{p}{(}\PYG{p}{)}

    \PYG{c+c1}{\PYGZsh{} \PYGZhy{}\PYGZhy{}\PYGZhy{} Convert to real coordinates \PYGZhy{}\PYGZhy{}\PYGZhy{}}
    \PYG{n}{xscale} \PYG{o}{=} \PYG{n}{width} \PYG{o}{/} \PYG{n}{roi\PYGZus{}w}
    \PYG{n}{yscale} \PYG{o}{=} \PYG{n}{height} \PYG{o}{/} \PYG{n}{roi\PYGZus{}h}

    \PYG{n}{corrected\PYGZus{}array} \PYG{o}{=} \PYG{n}{np}\PYG{o}{.}\PYG{n}{column\PYGZus{}stack}\PYG{p}{(}\PYG{p}{[}
        \PYG{n}{pos\PYGZus{}array}\PYG{p}{[}\PYG{p}{:}\PYG{p}{,} \PYG{l+m+mi}{0}\PYG{p}{]} \PYG{o}{/} \PYG{n}{frame\PYGZus{}rate}\PYG{p}{,}
        \PYG{n}{pos\PYGZus{}array}\PYG{p}{[}\PYG{p}{:}\PYG{p}{,} \PYG{l+m+mi}{1}\PYG{p}{]} \PYG{o}{*} \PYG{n}{xscale}\PYG{p}{,}
        \PYG{n}{pos\PYGZus{}array}\PYG{p}{[}\PYG{p}{:}\PYG{p}{,} \PYG{l+m+mi}{2}\PYG{p}{]} \PYG{o}{*} \PYG{n}{yscale}\PYG{p}{,}
    \PYG{p}{]}\PYG{p}{)}

    \PYG{n}{skipped} \PYG{o}{=} \PYG{n}{np}\PYG{o}{.}\PYG{n}{sum}\PYG{p}{(}\PYG{n}{np}\PYG{o}{.}\PYG{n}{isnan}\PYG{p}{(}\PYG{n}{corrected\PYGZus{}array}\PYG{p}{[}\PYG{p}{:}\PYG{p}{,} \PYG{l+m+mi}{1}\PYG{p}{]}\PYG{p}{)}\PYG{p}{)}
    \PYG{n+nb}{print}\PYG{p}{(}\PYG{l+s+sa}{f}\PYG{l+s+s2}{\PYGZdq{}}\PYG{l+s+s2}{  }\PYG{l+s+si}{\PYGZob{}}\PYG{n}{skipped}\PYG{l+s+si}{\PYGZcb{}}\PYG{l+s+s2}{ points skipped out of }\PYG{l+s+si}{\PYGZob{}}\PYG{n}{nfrm}\PYG{l+s+si}{\PYGZcb{}}\PYG{l+s+s2}{.}\PYG{l+s+s2}{\PYGZdq{}}\PYG{p}{)}

    \PYG{c+c1}{\PYGZsh{} Apply filters and interpolation}
    \PYG{n}{corrected\PYGZus{}array} \PYG{o}{=} \PYG{n}{boundary\PYGZus{}filter}\PYG{p}{(}\PYG{n}{corrected\PYGZus{}array}\PYG{p}{,} \PYG{n}{diameter}\PYG{p}{)}
    \PYG{n}{corrected\PYGZus{}array} \PYG{o}{=} \PYG{n}{displacement\PYGZus{}filter}\PYG{p}{(}\PYG{n}{corrected\PYGZus{}array}\PYG{p}{)}
    \PYG{n}{corrected\PYGZus{}array} \PYG{o}{=} \PYG{n}{dist\PYGZus{}filter}\PYG{p}{(}\PYG{n}{corrected\PYGZus{}array}\PYG{p}{,} \PYG{l+m+mi}{2}\PYG{p}{)}
    \PYG{n}{corrected\PYGZus{}array} \PYG{o}{=} \PYG{n}{interpolate\PYGZus{}pos}\PYG{p}{(}\PYG{n}{corrected\PYGZus{}array}\PYG{p}{,} \PYG{l+m+mi}{2}\PYG{p}{)}

    \PYG{k}{return} \PYG{n}{corrected\PYGZus{}array}\PYG{p}{,} \PYG{n}{frame\PYGZus{}rate}

\PYG{n+nb}{print}\PYG{p}{(}\PYG{l+s+s1}{\PYGZsq{}}\PYG{l+s+s1}{Helper functions defined.}\PYG{l+s+s1}{\PYGZsq{}}\PYG{p}{)}
\end{sphinxVerbatim}

\end{sphinxuseclass}\end{sphinxVerbatimInput}

\end{sphinxuseclass}
\sphinxAtStartPar
The cell below is going to prompt you to select an ROI using a square. You want to select \sphinxstylestrong{the entire petri dish} (you may need to guess where the edges are if they are hard to see).

\begin{sphinxuseclass}{cell}\begin{sphinxVerbatimInput}

\begin{sphinxuseclass}{cell_input}
\begin{sphinxVerbatim}[commandchars=\\\{\}]
\PYG{c+c1}{\PYGZsh{}@title Interactive ROI selector — click and drag on the image \PYGZob{} display\PYGZhy{}mode: \PYGZdq{}form\PYGZdq{} \PYGZcb{}}

\PYG{c+c1}{\PYGZsh{} Read first frame and encode as base64 for display}
\PYG{n}{cap} \PYG{o}{=} \PYG{n}{cv2}\PYG{o}{.}\PYG{n}{VideoCapture}\PYG{p}{(}\PYG{n}{current\PYGZus{}video}\PYG{p}{)}
\PYG{n}{ret}\PYG{p}{,} \PYG{n}{frame} \PYG{o}{=} \PYG{n}{cap}\PYG{o}{.}\PYG{n}{read}\PYG{p}{(}\PYG{p}{)}
\PYG{n}{cap}\PYG{o}{.}\PYG{n}{release}\PYG{p}{(}\PYG{p}{)}
\PYG{k}{if} \PYG{o+ow}{not} \PYG{n}{ret}\PYG{p}{:}
    \PYG{k}{raise} \PYG{n+ne}{ValueError}\PYG{p}{(}\PYG{l+s+sa}{f}\PYG{l+s+s1}{\PYGZsq{}}\PYG{l+s+s1}{Could not read first frame of }\PYG{l+s+si}{\PYGZob{}}\PYG{n}{current\PYGZus{}video}\PYG{l+s+si}{\PYGZcb{}}\PYG{l+s+s1}{\PYGZsq{}}\PYG{p}{)}

\PYG{n}{frame\PYGZus{}rgb} \PYG{o}{=} \PYG{n}{cv2}\PYG{o}{.}\PYG{n}{cvtColor}\PYG{p}{(}\PYG{n}{frame}\PYG{p}{,} \PYG{n}{cv2}\PYG{o}{.}\PYG{n}{COLOR\PYGZus{}BGR2RGB}\PYG{p}{)}
\PYG{n}{img\PYGZus{}h}\PYG{p}{,} \PYG{n}{img\PYGZus{}w} \PYG{o}{=} \PYG{n}{frame\PYGZus{}rgb}\PYG{o}{.}\PYG{n}{shape}\PYG{p}{[}\PYG{p}{:}\PYG{l+m+mi}{2}\PYG{p}{]}

\PYG{c+c1}{\PYGZsh{} Encode frame as PNG base64}
\PYG{n}{\PYGZus{}}\PYG{p}{,} \PYG{n}{buf} \PYG{o}{=} \PYG{n}{cv2}\PYG{o}{.}\PYG{n}{imencode}\PYG{p}{(}\PYG{l+s+s1}{\PYGZsq{}}\PYG{l+s+s1}{.png}\PYG{l+s+s1}{\PYGZsq{}}\PYG{p}{,} \PYG{n}{cv2}\PYG{o}{.}\PYG{n}{cvtColor}\PYG{p}{(}\PYG{n}{frame\PYGZus{}rgb}\PYG{p}{,} \PYG{n}{cv2}\PYG{o}{.}\PYG{n}{COLOR\PYGZus{}RGB2BGR}\PYG{p}{)}\PYG{p}{)}
\PYG{n}{img\PYGZus{}b64} \PYG{o}{=} \PYG{n}{base64}\PYG{o}{.}\PYG{n}{b64encode}\PYG{p}{(}\PYG{n}{buf}\PYG{p}{)}\PYG{o}{.}\PYG{n}{decode}\PYG{p}{(}\PYG{l+s+s1}{\PYGZsq{}}\PYG{l+s+s1}{utf\PYGZhy{}8}\PYG{l+s+s1}{\PYGZsq{}}\PYG{p}{)}

\PYG{n+nb}{print}\PYG{p}{(}\PYG{l+s+sa}{f}\PYG{l+s+s1}{\PYGZsq{}}\PYG{l+s+s1}{Image size: }\PYG{l+s+si}{\PYGZob{}}\PYG{n}{img\PYGZus{}w}\PYG{l+s+si}{\PYGZcb{}}\PYG{l+s+s1}{ x }\PYG{l+s+si}{\PYGZob{}}\PYG{n}{img\PYGZus{}h}\PYG{l+s+si}{\PYGZcb{}}\PYG{l+s+s1}{ pixels}\PYG{l+s+s1}{\PYGZsq{}}\PYG{p}{)}
\PYG{n+nb}{print}\PYG{p}{(}\PYG{l+s+s1}{\PYGZsq{}}\PYG{l+s+s1}{Click and drag on the image below to select the ROI, then click }\PYG{l+s+s1}{\PYGZdq{}}\PYG{l+s+s1}{Confirm ROI}\PYG{l+s+s1}{\PYGZdq{}}\PYG{l+s+s1}{.}\PYG{l+s+se}{\PYGZbs{}n}\PYG{l+s+s1}{\PYGZsq{}}\PYG{p}{)}

\PYG{c+c1}{\PYGZsh{} JavaScript/HTML interactive ROI selector}
\PYG{n}{html\PYGZus{}code} \PYG{o}{=} \PYG{l+s+sa}{f}\PYG{l+s+s2}{\PYGZdq{}\PYGZdq{}\PYGZdq{}}
\PYG{l+s+s2}{\PYGZlt{}div id=}\PYG{l+s+s2}{\PYGZdq{}}\PYG{l+s+s2}{roi\PYGZhy{}container}\PYG{l+s+s2}{\PYGZdq{}}\PYG{l+s+s2}{ style=}\PYG{l+s+s2}{\PYGZdq{}}\PYG{l+s+s2}{position:relative; display:inline\PYGZhy{}block; cursor:crosshair;}\PYG{l+s+s2}{\PYGZdq{}}\PYG{l+s+s2}{\PYGZgt{}}
\PYG{l+s+s2}{  \PYGZlt{}canvas id=}\PYG{l+s+s2}{\PYGZdq{}}\PYG{l+s+s2}{roi\PYGZhy{}canvas}\PYG{l+s+s2}{\PYGZdq{}}\PYG{l+s+s2}{ width=}\PYG{l+s+s2}{\PYGZdq{}}\PYG{l+s+si}{\PYGZob{}}\PYG{n}{img\PYGZus{}w}\PYG{l+s+si}{\PYGZcb{}}\PYG{l+s+s2}{\PYGZdq{}}\PYG{l+s+s2}{ height=}\PYG{l+s+s2}{\PYGZdq{}}\PYG{l+s+si}{\PYGZob{}}\PYG{n}{img\PYGZus{}h}\PYG{l+s+si}{\PYGZcb{}}\PYG{l+s+s2}{\PYGZdq{}}
\PYG{l+s+s2}{    style=}\PYG{l+s+s2}{\PYGZdq{}}\PYG{l+s+s2}{max\PYGZhy{}width:67\PYGZpc{}; height:auto; border:1px solid \PYGZsh{}ccc;}\PYG{l+s+s2}{\PYGZdq{}}\PYG{l+s+s2}{\PYGZgt{}\PYGZlt{}/canvas\PYGZgt{}}
\PYG{l+s+s2}{\PYGZlt{}/div\PYGZgt{}}
\PYG{l+s+s2}{\PYGZlt{}br\PYGZgt{}}
\PYG{l+s+s2}{\PYGZlt{}button id=}\PYG{l+s+s2}{\PYGZdq{}}\PYG{l+s+s2}{roi\PYGZhy{}confirm\PYGZhy{}btn}\PYG{l+s+s2}{\PYGZdq{}}\PYG{l+s+s2}{ style=}\PYG{l+s+s2}{\PYGZdq{}}\PYG{l+s+s2}{font\PYGZhy{}size:16px; padding:8px 20px; background:\PYGZsh{}4CAF50;}
\PYG{l+s+s2}{  color:white; border:none; border\PYGZhy{}radius:4px; cursor:pointer; margin\PYGZhy{}top:8px;}\PYG{l+s+s2}{\PYGZdq{}}\PYG{l+s+s2}{\PYGZgt{}}
\PYG{l+s+s2}{  Confirm ROI}
\PYG{l+s+s2}{\PYGZlt{}/button\PYGZgt{}}
\PYG{l+s+s2}{\PYGZlt{}p id=}\PYG{l+s+s2}{\PYGZdq{}}\PYG{l+s+s2}{roi\PYGZhy{}info}\PYG{l+s+s2}{\PYGZdq{}}\PYG{l+s+s2}{ style=}\PYG{l+s+s2}{\PYGZdq{}}\PYG{l+s+s2}{font\PYGZhy{}size:14px; color:\PYGZsh{}333;}\PYG{l+s+s2}{\PYGZdq{}}\PYG{l+s+s2}{\PYGZgt{}\PYGZlt{}/p\PYGZgt{}}

\PYG{l+s+s2}{\PYGZlt{}script\PYGZgt{}}
\PYG{l+s+s2}{(function() }\PYG{l+s+se}{\PYGZob{}\PYGZob{}}
\PYG{l+s+s2}{  var canvas = document.getElementById(}\PYG{l+s+s2}{\PYGZsq{}}\PYG{l+s+s2}{roi\PYGZhy{}canvas}\PYG{l+s+s2}{\PYGZsq{}}\PYG{l+s+s2}{);}
\PYG{l+s+s2}{  var ctx = canvas.getContext(}\PYG{l+s+s2}{\PYGZsq{}}\PYG{l+s+s2}{2d}\PYG{l+s+s2}{\PYGZsq{}}\PYG{l+s+s2}{);}
\PYG{l+s+s2}{  var img = new window.Image();}
\PYG{l+s+s2}{  img.src = }\PYG{l+s+s2}{\PYGZsq{}}\PYG{l+s+s2}{data:image/png;base64,}\PYG{l+s+si}{\PYGZob{}}\PYG{n}{img\PYGZus{}b64}\PYG{l+s+si}{\PYGZcb{}}\PYG{l+s+s2}{\PYGZsq{}}\PYG{l+s+s2}{;}

\PYG{l+s+s2}{  var drawing = false;}
\PYG{l+s+s2}{  var x0 = 0, y0 = 0, x1 = 0, y1 = 0;}
\PYG{l+s+s2}{  var hasROI = false;}

\PYG{l+s+s2}{  // Scale factor: canvas CSS size vs actual pixel size}
\PYG{l+s+s2}{  function getScale() }\PYG{l+s+se}{\PYGZob{}\PYGZob{}}
\PYG{l+s+s2}{    var rect = canvas.getBoundingClientRect();}
\PYG{l+s+s2}{    return }\PYG{l+s+se}{\PYGZob{}\PYGZob{}}\PYG{l+s+s2}{ sx: canvas.width / rect.width, sy: canvas.height / rect.height }\PYG{l+s+se}{\PYGZcb{}\PYGZcb{}}\PYG{l+s+s2}{;}
\PYG{l+s+s2}{  }\PYG{l+s+se}{\PYGZcb{}\PYGZcb{}}

\PYG{l+s+s2}{  img.onload = function() }\PYG{l+s+se}{\PYGZob{}\PYGZob{}}
\PYG{l+s+s2}{    ctx.drawImage(img, 0, 0);}
\PYG{l+s+s2}{  }\PYG{l+s+se}{\PYGZcb{}\PYGZcb{}}\PYG{l+s+s2}{;}

\PYG{l+s+s2}{  function redraw() }\PYG{l+s+se}{\PYGZob{}\PYGZob{}}
\PYG{l+s+s2}{    ctx.clearRect(0, 0, canvas.width, canvas.height);}
\PYG{l+s+s2}{    ctx.drawImage(img, 0, 0);}
\PYG{l+s+s2}{    if (hasROI) }\PYG{l+s+se}{\PYGZob{}\PYGZob{}}
\PYG{l+s+s2}{      var rx = Math.min(x0, x1);}
\PYG{l+s+s2}{      var ry = Math.min(y0, y1);}
\PYG{l+s+s2}{      var rw = Math.abs(x1 \PYGZhy{} x0);}
\PYG{l+s+s2}{      var rh = Math.abs(y1 \PYGZhy{} y0);}
\PYG{l+s+s2}{      ctx.strokeStyle = }\PYG{l+s+s2}{\PYGZsq{}}\PYG{l+s+s2}{red}\PYG{l+s+s2}{\PYGZsq{}}\PYG{l+s+s2}{;}
\PYG{l+s+s2}{      ctx.lineWidth = 3;}
\PYG{l+s+s2}{      ctx.strokeRect(rx, ry, rw, rh);}
\PYG{l+s+s2}{      document.getElementById(}\PYG{l+s+s2}{\PYGZsq{}}\PYG{l+s+s2}{roi\PYGZhy{}info}\PYG{l+s+s2}{\PYGZsq{}}\PYG{l+s+s2}{).innerText =}
\PYG{l+s+s2}{        }\PYG{l+s+s2}{\PYGZsq{}}\PYG{l+s+s2}{ROI: x=}\PYG{l+s+s2}{\PYGZsq{}}\PYG{l+s+s2}{ + Math.round(rx) + }\PYG{l+s+s2}{\PYGZsq{}}\PYG{l+s+s2}{, y=}\PYG{l+s+s2}{\PYGZsq{}}\PYG{l+s+s2}{ + Math.round(ry) +}
\PYG{l+s+s2}{        }\PYG{l+s+s2}{\PYGZsq{}}\PYG{l+s+s2}{, w=}\PYG{l+s+s2}{\PYGZsq{}}\PYG{l+s+s2}{ + Math.round(rw) + }\PYG{l+s+s2}{\PYGZsq{}}\PYG{l+s+s2}{, h=}\PYG{l+s+s2}{\PYGZsq{}}\PYG{l+s+s2}{ + Math.round(rh);}
\PYG{l+s+s2}{    }\PYG{l+s+se}{\PYGZcb{}\PYGZcb{}}
\PYG{l+s+s2}{  }\PYG{l+s+se}{\PYGZcb{}\PYGZcb{}}

\PYG{l+s+s2}{  canvas.addEventListener(}\PYG{l+s+s2}{\PYGZsq{}}\PYG{l+s+s2}{mousedown}\PYG{l+s+s2}{\PYGZsq{}}\PYG{l+s+s2}{, function(e) }\PYG{l+s+se}{\PYGZob{}\PYGZob{}}
\PYG{l+s+s2}{    var rect = canvas.getBoundingClientRect();}
\PYG{l+s+s2}{    var scale = getScale();}
\PYG{l+s+s2}{    x0 = (e.clientX \PYGZhy{} rect.left) * scale.sx;}
\PYG{l+s+s2}{    y0 = (e.clientY \PYGZhy{} rect.top) * scale.sy;}
\PYG{l+s+s2}{    drawing = true;}
\PYG{l+s+s2}{    hasROI = false;}
\PYG{l+s+s2}{  }\PYG{l+s+se}{\PYGZcb{}\PYGZcb{}}\PYG{l+s+s2}{);}

\PYG{l+s+s2}{  canvas.addEventListener(}\PYG{l+s+s2}{\PYGZsq{}}\PYG{l+s+s2}{mousemove}\PYG{l+s+s2}{\PYGZsq{}}\PYG{l+s+s2}{, function(e) }\PYG{l+s+se}{\PYGZob{}\PYGZob{}}
\PYG{l+s+s2}{    if (!drawing) return;}
\PYG{l+s+s2}{    var rect = canvas.getBoundingClientRect();}
\PYG{l+s+s2}{    var scale = getScale();}
\PYG{l+s+s2}{    x1 = (e.clientX \PYGZhy{} rect.left) * scale.sx;}
\PYG{l+s+s2}{    y1 = (e.clientY \PYGZhy{} rect.top) * scale.sy;}
\PYG{l+s+s2}{    hasROI = true;}
\PYG{l+s+s2}{    redraw();}
\PYG{l+s+s2}{  }\PYG{l+s+se}{\PYGZcb{}\PYGZcb{}}\PYG{l+s+s2}{);}

\PYG{l+s+s2}{  canvas.addEventListener(}\PYG{l+s+s2}{\PYGZsq{}}\PYG{l+s+s2}{mouseup}\PYG{l+s+s2}{\PYGZsq{}}\PYG{l+s+s2}{, function(e) }\PYG{l+s+se}{\PYGZob{}\PYGZob{}}
\PYG{l+s+s2}{    drawing = false;}
\PYG{l+s+s2}{  }\PYG{l+s+se}{\PYGZcb{}\PYGZcb{}}\PYG{l+s+s2}{);}

\PYG{l+s+s2}{  document.getElementById(}\PYG{l+s+s2}{\PYGZsq{}}\PYG{l+s+s2}{roi\PYGZhy{}confirm\PYGZhy{}btn}\PYG{l+s+s2}{\PYGZsq{}}\PYG{l+s+s2}{).addEventListener(}\PYG{l+s+s2}{\PYGZsq{}}\PYG{l+s+s2}{click}\PYG{l+s+s2}{\PYGZsq{}}\PYG{l+s+s2}{, function() }\PYG{l+s+se}{\PYGZob{}\PYGZob{}}
\PYG{l+s+s2}{    if (!hasROI) }\PYG{l+s+se}{\PYGZob{}\PYGZob{}}
\PYG{l+s+s2}{      document.getElementById(}\PYG{l+s+s2}{\PYGZsq{}}\PYG{l+s+s2}{roi\PYGZhy{}info}\PYG{l+s+s2}{\PYGZsq{}}\PYG{l+s+s2}{).innerText =}
\PYG{l+s+s2}{        }\PYG{l+s+s2}{\PYGZsq{}}\PYG{l+s+s2}{No ROI drawn yet! Click and drag on the image first.}\PYG{l+s+s2}{\PYGZsq{}}\PYG{l+s+s2}{;}
\PYG{l+s+s2}{      return;}
\PYG{l+s+s2}{    }\PYG{l+s+se}{\PYGZcb{}\PYGZcb{}}
\PYG{l+s+s2}{    var rx = Math.round(Math.min(x0, x1));}
\PYG{l+s+s2}{    var ry = Math.round(Math.min(y0, y1));}
\PYG{l+s+s2}{    var rw = Math.round(Math.abs(x1 \PYGZhy{} x0));}
\PYG{l+s+s2}{    var rh = Math.round(Math.abs(y1 \PYGZhy{} y0));}

\PYG{l+s+s2}{    // Send ROI back to Python via Colab}\PYG{l+s+s2}{\PYGZsq{}}\PYG{l+s+s2}{s kernel eval}
\PYG{l+s+s2}{    google.colab.kernel.invokeFunction(}\PYG{l+s+s2}{\PYGZsq{}}\PYG{l+s+s2}{set\PYGZus{}roi}\PYG{l+s+s2}{\PYGZsq{}}\PYG{l+s+s2}{, [rx, ry, rw, rh], }\PYG{l+s+se}{\PYGZob{}\PYGZob{}}\PYG{l+s+se}{\PYGZcb{}\PYGZcb{}}\PYG{l+s+s2}{);}
\PYG{l+s+s2}{    document.getElementById(}\PYG{l+s+s2}{\PYGZsq{}}\PYG{l+s+s2}{roi\PYGZhy{}info}\PYG{l+s+s2}{\PYGZsq{}}\PYG{l+s+s2}{).innerText =}
\PYG{l+s+s2}{      }\PYG{l+s+s2}{\PYGZsq{}}\PYG{l+s+s2}{ROI CONFIRMED: x=}\PYG{l+s+s2}{\PYGZsq{}}\PYG{l+s+s2}{ + rx + }\PYG{l+s+s2}{\PYGZsq{}}\PYG{l+s+s2}{, y=}\PYG{l+s+s2}{\PYGZsq{}}\PYG{l+s+s2}{ + ry + }\PYG{l+s+s2}{\PYGZsq{}}\PYG{l+s+s2}{, w=}\PYG{l+s+s2}{\PYGZsq{}}\PYG{l+s+s2}{ + rw + }\PYG{l+s+s2}{\PYGZsq{}}\PYG{l+s+s2}{, h=}\PYG{l+s+s2}{\PYGZsq{}}\PYG{l+s+s2}{ + rh;}
\PYG{l+s+s2}{    document.getElementById(}\PYG{l+s+s2}{\PYGZsq{}}\PYG{l+s+s2}{roi\PYGZhy{}confirm\PYGZhy{}btn}\PYG{l+s+s2}{\PYGZsq{}}\PYG{l+s+s2}{).style.background = }\PYG{l+s+s2}{\PYGZsq{}}\PYG{l+s+s2}{\PYGZsh{}888}\PYG{l+s+s2}{\PYGZsq{}}\PYG{l+s+s2}{;}
\PYG{l+s+s2}{    document.getElementById(}\PYG{l+s+s2}{\PYGZsq{}}\PYG{l+s+s2}{roi\PYGZhy{}confirm\PYGZhy{}btn}\PYG{l+s+s2}{\PYGZsq{}}\PYG{l+s+s2}{).innerText = }\PYG{l+s+s2}{\PYGZsq{}}\PYG{l+s+s2}{ROI Confirmed!}\PYG{l+s+s2}{\PYGZsq{}}\PYG{l+s+s2}{;}
\PYG{l+s+s2}{  }\PYG{l+s+se}{\PYGZcb{}\PYGZcb{}}\PYG{l+s+s2}{);}
\PYG{l+s+se}{\PYGZcb{}\PYGZcb{}}\PYG{l+s+s2}{)();}
\PYG{l+s+s2}{\PYGZlt{}/script\PYGZgt{}}
\PYG{l+s+s2}{\PYGZdq{}\PYGZdq{}\PYGZdq{}}

\PYG{c+c1}{\PYGZsh{} Register a callback so JavaScript can send the ROI to Python}
\PYG{n}{roi\PYGZus{}result} \PYG{o}{=} \PYG{p}{\PYGZob{}}\PYG{l+s+s1}{\PYGZsq{}}\PYG{l+s+s1}{value}\PYG{l+s+s1}{\PYGZsq{}}\PYG{p}{:} \PYG{k+kc}{None}\PYG{p}{\PYGZcb{}}

\PYG{k}{def}\PYG{+w}{ }\PYG{n+nf}{set\PYGZus{}roi}\PYG{p}{(}\PYG{n}{x}\PYG{p}{,} \PYG{n}{y}\PYG{p}{,} \PYG{n}{w}\PYG{p}{,} \PYG{n}{h}\PYG{p}{)}\PYG{p}{:}
    \PYG{n}{roi\PYGZus{}result}\PYG{p}{[}\PYG{l+s+s1}{\PYGZsq{}}\PYG{l+s+s1}{value}\PYG{l+s+s1}{\PYGZsq{}}\PYG{p}{]} \PYG{o}{=} \PYG{p}{(}\PYG{n+nb}{int}\PYG{p}{(}\PYG{n}{x}\PYG{p}{)}\PYG{p}{,} \PYG{n+nb}{int}\PYG{p}{(}\PYG{n}{y}\PYG{p}{)}\PYG{p}{,} \PYG{n+nb}{int}\PYG{p}{(}\PYG{n}{w}\PYG{p}{)}\PYG{p}{,} \PYG{n+nb}{int}\PYG{p}{(}\PYG{n}{h}\PYG{p}{)}\PYG{p}{)}
    \PYG{n+nb}{print}\PYG{p}{(}\PYG{l+s+sa}{f}\PYG{l+s+s1}{\PYGZsq{}}\PYG{l+s+s1}{ROI received: x=}\PYG{l+s+si}{\PYGZob{}}\PYG{n+nb}{int}\PYG{p}{(}\PYG{n}{x}\PYG{p}{)}\PYG{l+s+si}{\PYGZcb{}}\PYG{l+s+s1}{, y=}\PYG{l+s+si}{\PYGZob{}}\PYG{n+nb}{int}\PYG{p}{(}\PYG{n}{y}\PYG{p}{)}\PYG{l+s+si}{\PYGZcb{}}\PYG{l+s+s1}{, w=}\PYG{l+s+si}{\PYGZob{}}\PYG{n+nb}{int}\PYG{p}{(}\PYG{n}{w}\PYG{p}{)}\PYG{l+s+si}{\PYGZcb{}}\PYG{l+s+s1}{, h=}\PYG{l+s+si}{\PYGZob{}}\PYG{n+nb}{int}\PYG{p}{(}\PYG{n}{h}\PYG{p}{)}\PYG{l+s+si}{\PYGZcb{}}\PYG{l+s+s1}{\PYGZsq{}}\PYG{p}{)}

\PYG{n}{output}\PYG{o}{.}\PYG{n}{register\PYGZus{}callback}\PYG{p}{(}\PYG{l+s+s1}{\PYGZsq{}}\PYG{l+s+s1}{set\PYGZus{}roi}\PYG{l+s+s1}{\PYGZsq{}}\PYG{p}{,} \PYG{n}{set\PYGZus{}roi}\PYG{p}{)}

\PYG{n}{display}\PYG{p}{(}\PYG{n}{HTML}\PYG{p}{(}\PYG{n}{html\PYGZus{}code}\PYG{p}{)}\PYG{p}{)}
\end{sphinxVerbatim}

\end{sphinxuseclass}\end{sphinxVerbatimInput}

\end{sphinxuseclass}
\begin{sphinxuseclass}{cell}\begin{sphinxVerbatimInput}

\begin{sphinxuseclass}{cell_input}
\begin{sphinxVerbatim}[commandchars=\\\{\}]
\PYG{c+c1}{\PYGZsh{} Run this cell AFTER you have confirmed your ROI above}
\PYG{n}{roi} \PYG{o}{=} \PYG{n}{roi\PYGZus{}result}\PYG{p}{[}\PYG{l+s+s1}{\PYGZsq{}}\PYG{l+s+s1}{value}\PYG{l+s+s1}{\PYGZsq{}}\PYG{p}{]}

\PYG{k}{if} \PYG{n}{roi} \PYG{o+ow}{is} \PYG{k+kc}{None}\PYG{p}{:}
    \PYG{k}{raise} \PYG{n+ne}{ValueError}\PYG{p}{(}\PYG{l+s+s2}{\PYGZdq{}}\PYG{l+s+s2}{No ROI selected! Go back and draw a rectangle on the image, then click }\PYG{l+s+s2}{\PYGZsq{}}\PYG{l+s+s2}{Confirm ROI}\PYG{l+s+s2}{\PYGZsq{}}\PYG{l+s+s2}{.}\PYG{l+s+s2}{\PYGZdq{}}\PYG{p}{)}

\PYG{n+nb}{print}\PYG{p}{(}\PYG{l+s+sa}{f}\PYG{l+s+s1}{\PYGZsq{}}\PYG{l+s+s1}{Using ROI: x=}\PYG{l+s+si}{\PYGZob{}}\PYG{n}{roi}\PYG{p}{[}\PYG{l+m+mi}{0}\PYG{p}{]}\PYG{l+s+si}{\PYGZcb{}}\PYG{l+s+s1}{, y=}\PYG{l+s+si}{\PYGZob{}}\PYG{n}{roi}\PYG{p}{[}\PYG{l+m+mi}{1}\PYG{p}{]}\PYG{l+s+si}{\PYGZcb{}}\PYG{l+s+s1}{, width=}\PYG{l+s+si}{\PYGZob{}}\PYG{n}{roi}\PYG{p}{[}\PYG{l+m+mi}{2}\PYG{p}{]}\PYG{l+s+si}{\PYGZcb{}}\PYG{l+s+s1}{, height=}\PYG{l+s+si}{\PYGZob{}}\PYG{n}{roi}\PYG{p}{[}\PYG{l+m+mi}{3}\PYG{p}{]}\PYG{l+s+si}{\PYGZcb{}}\PYG{l+s+s1}{\PYGZsq{}}\PYG{p}{)}
\end{sphinxVerbatim}

\end{sphinxuseclass}\end{sphinxVerbatimInput}

\end{sphinxuseclass}

\subsection{Track Fly in Video(s)}
\label{\detokenize{FlyTracking/ColabNotebook:track-fly-in-video-s}}
\sphinxAtStartPar
This processes each video frame\sphinxhyphen{}by\sphinxhyphen{}frame:
\begin{enumerate}
\sphinxsetlistlabels{\arabic}{enumi}{enumii}{}{.}%
\item {} 
\sphinxAtStartPar
Creates a background image from 100 random frames

\item {} 
\sphinxAtStartPar
Subtracts the background from each frame

\item {} 
\sphinxAtStartPar
Finds the fly position using center\sphinxhyphen{}of\sphinxhyphen{}mass of dark pixels

\item {} 
\sphinxAtStartPar
Applies teleportation filter and interpolation

\item {} 
\sphinxAtStartPar
Converts pixel positions to centimeters

\end{enumerate}

\begin{sphinxuseclass}{cell}\begin{sphinxVerbatimInput}

\begin{sphinxuseclass}{cell_input}
\begin{sphinxVerbatim}[commandchars=\\\{\}]
\PYG{c+c1}{\PYGZsh{} \PYGZhy{}\PYGZhy{}\PYGZhy{} Process the current video \PYGZhy{}\PYGZhy{}\PYGZhy{}}
\PYG{n}{corrected}\PYG{p}{,} \PYG{n}{detected\PYGZus{}fps} \PYG{o}{=} \PYG{n}{process\PYGZus{}video}\PYG{p}{(}\PYG{n}{current\PYGZus{}video}\PYG{p}{,} \PYG{n}{roi}\PYG{p}{,} \PYG{n}{diameter}\PYG{p}{,} \PYG{n}{frame\PYGZus{}rate}\PYG{p}{,}
                                        \PYG{n}{search\PYGZus{}size}\PYG{p}{,} \PYG{n}{per\PYGZus{}pixel\PYGZus{}threshold}\PYG{p}{)}
\PYG{n}{all\PYGZus{}corrected}\PYG{o}{.}\PYG{n}{append}\PYG{p}{(}\PYG{n}{corrected}\PYG{p}{)}
\PYG{n}{all\PYGZus{}fps}\PYG{o}{.}\PYG{n}{append}\PYG{p}{(}\PYG{n}{detected\PYGZus{}fps}\PYG{p}{)}
\PYG{n}{all\PYGZus{}video\PYGZus{}paths}\PYG{o}{.}\PYG{n}{append}\PYG{p}{(}\PYG{n}{current\PYGZus{}video}\PYG{p}{)}

\PYG{n+nb}{print}\PYG{p}{(}\PYG{l+s+sa}{f}\PYG{l+s+s1}{\PYGZsq{}}\PYG{l+s+se}{\PYGZbs{}n}\PYG{l+s+s1}{Done! Total videos processed: }\PYG{l+s+si}{\PYGZob{}}\PYG{n+nb}{len}\PYG{p}{(}\PYG{n}{all\PYGZus{}corrected}\PYG{p}{)}\PYG{l+s+si}{\PYGZcb{}}\PYG{l+s+s1}{\PYGZsq{}}\PYG{p}{)}
\PYG{n+nb}{print}\PYG{p}{(}\PYG{l+s+s1}{\PYGZsq{}}\PYG{l+s+s1}{To process another video, go back to }\PYG{l+s+s1}{\PYGZdq{}}\PYG{l+s+s1}{Select Video File}\PYG{l+s+s1}{\PYGZdq{}}\PYG{l+s+s1}{ cell and change `current\PYGZus{}video`, then run this cell again.}\PYG{l+s+s1}{\PYGZsq{}}\PYG{p}{)}
\end{sphinxVerbatim}

\end{sphinxuseclass}\end{sphinxVerbatimInput}

\end{sphinxuseclass}

\subsection{Plot Fly Path(s)}
\label{\detokenize{FlyTracking/ColabNotebook:plot-fly-path-s}}
\sphinxAtStartPar
Color\sphinxhyphen{}coded by time (light blue = start, dark purple = end).

\begin{sphinxuseclass}{cell}\begin{sphinxVerbatimInput}

\begin{sphinxuseclass}{cell_input}
\begin{sphinxVerbatim}[commandchars=\\\{\}]
\PYG{k}{for} \PYG{n}{idx}\PYG{p}{,} \PYG{n}{corrected} \PYG{o+ow}{in} \PYG{n+nb}{enumerate}\PYG{p}{(}\PYG{n}{all\PYGZus{}corrected}\PYG{p}{)}\PYG{p}{:}
    \PYG{n}{x} \PYG{o}{=} \PYG{n}{corrected}\PYG{p}{[}\PYG{p}{:}\PYG{p}{,} \PYG{l+m+mi}{1}\PYG{p}{]}
    \PYG{n}{y} \PYG{o}{=} \PYG{n}{corrected}\PYG{p}{[}\PYG{p}{:}\PYG{p}{,} \PYG{l+m+mi}{2}\PYG{p}{]}
    \PYG{n}{t} \PYG{o}{=} \PYG{n}{corrected}\PYG{p}{[}\PYG{p}{:}\PYG{p}{,} \PYG{l+m+mi}{0}\PYG{p}{]}

    \PYG{n}{fig}\PYG{p}{,} \PYG{n}{ax} \PYG{o}{=} \PYG{n}{plt}\PYG{o}{.}\PYG{n}{subplots}\PYG{p}{(}\PYG{n}{figsize}\PYG{o}{=}\PYG{p}{(}\PYG{l+m+mi}{6}\PYG{p}{,} \PYG{l+m+mi}{6}\PYG{p}{)}\PYG{p}{)}

    \PYG{c+c1}{\PYGZsh{} Plot path colored by time}
    \PYG{n}{valid} \PYG{o}{=} \PYG{o}{\PYGZti{}}\PYG{n}{np}\PYG{o}{.}\PYG{n}{isnan}\PYG{p}{(}\PYG{n}{x}\PYG{p}{)} \PYG{o}{\PYGZam{}} \PYG{o}{\PYGZti{}}\PYG{n}{np}\PYG{o}{.}\PYG{n}{isnan}\PYG{p}{(}\PYG{n}{y}\PYG{p}{)}
    \PYG{n}{points} \PYG{o}{=} \PYG{n}{np}\PYG{o}{.}\PYG{n}{array}\PYG{p}{(}\PYG{p}{[}\PYG{n}{x}\PYG{p}{[}\PYG{n}{valid}\PYG{p}{]}\PYG{p}{,} \PYG{n}{y}\PYG{p}{[}\PYG{n}{valid}\PYG{p}{]}\PYG{p}{]}\PYG{p}{)}\PYG{o}{.}\PYG{n}{T}\PYG{o}{.}\PYG{n}{reshape}\PYG{p}{(}\PYG{o}{\PYGZhy{}}\PYG{l+m+mi}{1}\PYG{p}{,} \PYG{l+m+mi}{1}\PYG{p}{,} \PYG{l+m+mi}{2}\PYG{p}{)}
    \PYG{n}{segments} \PYG{o}{=} \PYG{n}{np}\PYG{o}{.}\PYG{n}{concatenate}\PYG{p}{(}\PYG{p}{[}\PYG{n}{points}\PYG{p}{[}\PYG{p}{:}\PYG{o}{\PYGZhy{}}\PYG{l+m+mi}{1}\PYG{p}{]}\PYG{p}{,} \PYG{n}{points}\PYG{p}{[}\PYG{l+m+mi}{1}\PYG{p}{:}\PYG{p}{]}\PYG{p}{]}\PYG{p}{,} \PYG{n}{axis}\PYG{o}{=}\PYG{l+m+mi}{1}\PYG{p}{)}

    \PYG{n}{lc} \PYG{o}{=} \PYG{n}{LineCollection}\PYG{p}{(}\PYG{n}{segments}\PYG{p}{,} \PYG{n}{cmap}\PYG{o}{=}\PYG{l+s+s1}{\PYGZsq{}}\PYG{l+s+s1}{BuPu}\PYG{l+s+s1}{\PYGZsq{}}\PYG{p}{,} \PYG{n}{linewidth}\PYG{o}{=}\PYG{l+m+mi}{2}\PYG{p}{)}
    \PYG{n}{lc}\PYG{o}{.}\PYG{n}{set\PYGZus{}array}\PYG{p}{(}\PYG{n}{t}\PYG{p}{[}\PYG{n}{valid}\PYG{p}{]}\PYG{p}{[}\PYG{p}{:}\PYG{o}{\PYGZhy{}}\PYG{l+m+mi}{1}\PYG{p}{]}\PYG{p}{)}
    \PYG{n}{ax}\PYG{o}{.}\PYG{n}{add\PYGZus{}collection}\PYG{p}{(}\PYG{n}{lc}\PYG{p}{)}

    \PYG{n}{ax}\PYG{o}{.}\PYG{n}{set\PYGZus{}xlim}\PYG{p}{(}\PYG{l+m+mi}{0}\PYG{p}{,} \PYG{n}{diameter}\PYG{p}{)}
    \PYG{n}{ax}\PYG{o}{.}\PYG{n}{set\PYGZus{}ylim}\PYG{p}{(}\PYG{n}{diameter}\PYG{p}{,} \PYG{l+m+mi}{0}\PYG{p}{)}  \PYG{c+c1}{\PYGZsh{} Invert y to match video}
    \PYG{n}{ax}\PYG{o}{.}\PYG{n}{set\PYGZus{}aspect}\PYG{p}{(}\PYG{l+s+s1}{\PYGZsq{}}\PYG{l+s+s1}{equal}\PYG{l+s+s1}{\PYGZsq{}}\PYG{p}{)}
    \PYG{n}{ax}\PYG{o}{.}\PYG{n}{set\PYGZus{}xlabel}\PYG{p}{(}\PYG{l+s+s1}{\PYGZsq{}}\PYG{l+s+s1}{X\PYGZhy{}coordinate (cm)}\PYG{l+s+s1}{\PYGZsq{}}\PYG{p}{,} \PYG{n}{fontsize}\PYG{o}{=}\PYG{l+m+mi}{11}\PYG{p}{)}
    \PYG{n}{ax}\PYG{o}{.}\PYG{n}{set\PYGZus{}ylabel}\PYG{p}{(}\PYG{l+s+s1}{\PYGZsq{}}\PYG{l+s+s1}{Y\PYGZhy{}coordinate (cm)}\PYG{l+s+s1}{\PYGZsq{}}\PYG{p}{,} \PYG{n}{fontsize}\PYG{o}{=}\PYG{l+m+mi}{11}\PYG{p}{)}
    \PYG{n}{ax}\PYG{o}{.}\PYG{n}{set\PYGZus{}title}\PYG{p}{(}\PYG{l+s+sa}{f}\PYG{l+s+s1}{\PYGZsq{}}\PYG{l+s+s1}{Fly Path — Video }\PYG{l+s+si}{\PYGZob{}}\PYG{n}{idx}\PYG{+w}{ }\PYG{o}{+}\PYG{+w}{ }\PYG{l+m+mi}{1}\PYG{l+s+si}{\PYGZcb{}}\PYG{l+s+s1}{\PYGZsq{}}\PYG{p}{)}

    \PYG{n}{cbar} \PYG{o}{=} \PYG{n}{fig}\PYG{o}{.}\PYG{n}{colorbar}\PYG{p}{(}\PYG{n}{lc}\PYG{p}{,} \PYG{n}{ax}\PYG{o}{=}\PYG{n}{ax}\PYG{p}{,} \PYG{n}{orientation}\PYG{o}{=}\PYG{l+s+s1}{\PYGZsq{}}\PYG{l+s+s1}{horizontal}\PYG{l+s+s1}{\PYGZsq{}}\PYG{p}{,} \PYG{n}{pad}\PYG{o}{=}\PYG{l+m+mf}{0.1}\PYG{p}{)}
    \PYG{n}{cbar}\PYG{o}{.}\PYG{n}{set\PYGZus{}label}\PYG{p}{(}\PYG{l+s+s1}{\PYGZsq{}}\PYG{l+s+s1}{Time (s)}\PYG{l+s+s1}{\PYGZsq{}}\PYG{p}{)}

    \PYG{n}{plt}\PYG{o}{.}\PYG{n}{tight\PYGZus{}layout}\PYG{p}{(}\PYG{p}{)}
    \PYG{n}{plt}\PYG{o}{.}\PYG{n}{show}\PYG{p}{(}\PYG{p}{)}
\end{sphinxVerbatim}

\end{sphinxuseclass}\end{sphinxVerbatimInput}

\end{sphinxuseclass}

\subsection{Calculate \& Plot Velocity}
\label{\detokenize{FlyTracking/ColabNotebook:calculate-plot-velocity}}
\begin{sphinxuseclass}{cell}\begin{sphinxVerbatimInput}

\begin{sphinxuseclass}{cell_input}
\begin{sphinxVerbatim}[commandchars=\\\{\}]
\PYG{c+c1}{\PYGZsh{} Calculate velocity for videos that haven\PYGZsq{}t been processed yet}
\PYG{n}{start\PYGZus{}idx} \PYG{o}{=} \PYG{n+nb}{len}\PYG{p}{(}\PYG{n}{all\PYGZus{}velocity}\PYG{p}{)}  \PYG{c+c1}{\PYGZsh{} Only calculate for newly added videos}

\PYG{k}{for} \PYG{n}{idx} \PYG{o+ow}{in} \PYG{n+nb}{range}\PYG{p}{(}\PYG{n}{start\PYGZus{}idx}\PYG{p}{,} \PYG{n+nb}{len}\PYG{p}{(}\PYG{n}{all\PYGZus{}corrected}\PYG{p}{)}\PYG{p}{)}\PYG{p}{:}
    \PYG{n}{corrected} \PYG{o}{=} \PYG{n}{all\PYGZus{}corrected}\PYG{p}{[}\PYG{n}{idx}\PYG{p}{]}
    \PYG{n}{x} \PYG{o}{=} \PYG{n}{corrected}\PYG{p}{[}\PYG{p}{:}\PYG{p}{,} \PYG{l+m+mi}{1}\PYG{p}{]}
    \PYG{n}{y} \PYG{o}{=} \PYG{n}{corrected}\PYG{p}{[}\PYG{p}{:}\PYG{p}{,} \PYG{l+m+mi}{2}\PYG{p}{]}

    \PYG{n}{fps} \PYG{o}{=} \PYG{n+nb}{round}\PYG{p}{(}\PYG{l+m+mf}{1.0} \PYG{o}{/} \PYG{n}{corrected}\PYG{p}{[}\PYG{l+m+mi}{1}\PYG{p}{,} \PYG{l+m+mi}{0}\PYG{p}{]}\PYG{p}{)} \PYG{k}{if} \PYG{n}{corrected}\PYG{p}{[}\PYG{l+m+mi}{1}\PYG{p}{,} \PYG{l+m+mi}{0}\PYG{p}{]} \PYG{o}{\PYGZgt{}} \PYG{l+m+mi}{0} \PYG{k}{else} \PYG{l+m+mi}{30}
    \PYG{n}{total\PYGZus{}time} \PYG{o}{=} \PYG{n+nb}{len}\PYG{p}{(}\PYG{n}{x}\PYG{p}{)} \PYG{o}{/} \PYG{n}{fps}
    \PYG{n}{total\PYGZus{}bins} \PYG{o}{=} \PYG{n+nb}{int}\PYG{p}{(}\PYG{n}{np}\PYG{o}{.}\PYG{n}{floor}\PYG{p}{(}\PYG{n}{total\PYGZus{}time} \PYG{o}{/} \PYG{n}{bin\PYGZus{}size}\PYG{p}{)}\PYG{p}{)}

    \PYG{c+c1}{\PYGZsh{} Calculate velocity per bin}
    \PYG{n}{data\PYGZus{}rate} \PYG{o}{=} \PYG{n+nb}{round}\PYG{p}{(}\PYG{l+m+mf}{1.0} \PYG{o}{/} \PYG{n}{corrected}\PYG{p}{[}\PYG{l+m+mi}{1}\PYG{p}{,} \PYG{l+m+mi}{0}\PYG{p}{]}\PYG{p}{)} \PYG{o}{*} \PYG{n}{bin\PYGZus{}size} \PYG{k}{if} \PYG{n}{corrected}\PYG{p}{[}\PYG{l+m+mi}{1}\PYG{p}{,} \PYG{l+m+mi}{0}\PYG{p}{]} \PYG{o}{\PYGZgt{}} \PYG{l+m+mi}{0} \PYG{k}{else} \PYG{n}{fps} \PYG{o}{*} \PYG{n}{bin\PYGZus{}size}
    \PYG{n}{data\PYGZus{}rate} \PYG{o}{=} \PYG{n+nb}{int}\PYG{p}{(}\PYG{n}{data\PYGZus{}rate}\PYG{p}{)}

    \PYG{k}{if} \PYG{n}{data\PYGZus{}rate} \PYG{o}{\PYGZlt{}} \PYG{l+m+mi}{1}\PYG{p}{:}
        \PYG{k}{raise} \PYG{n+ne}{ValueError}\PYG{p}{(}\PYG{l+s+s1}{\PYGZsq{}}\PYG{l+s+s1}{bin\PYGZus{}size is smaller than the minimum data rate.}\PYG{l+s+s1}{\PYGZsq{}}\PYG{p}{)}

    \PYG{n}{velocity} \PYG{o}{=} \PYG{n}{np}\PYG{o}{.}\PYG{n}{zeros}\PYG{p}{(}\PYG{n}{total\PYGZus{}bins}\PYG{p}{)}
    \PYG{k}{for} \PYG{n}{row} \PYG{o+ow}{in} \PYG{n+nb}{range}\PYG{p}{(}\PYG{l+m+mi}{0}\PYG{p}{,} \PYG{n+nb}{len}\PYG{p}{(}\PYG{n}{corrected}\PYG{p}{)} \PYG{o}{\PYGZhy{}} \PYG{n}{data\PYGZus{}rate}\PYG{p}{,} \PYG{n}{data\PYGZus{}rate}\PYG{p}{)}\PYG{p}{:}
        \PYG{n}{bin\PYGZus{}idx} \PYG{o}{=} \PYG{n}{row} \PYG{o}{/}\PYG{o}{/} \PYG{n}{data\PYGZus{}rate}
        \PYG{k}{if} \PYG{n}{bin\PYGZus{}idx} \PYG{o}{\PYGZgt{}}\PYG{o}{=} \PYG{n}{total\PYGZus{}bins}\PYG{p}{:}
            \PYG{k}{break}
        \PYG{n}{p1} \PYG{o}{=} \PYG{n}{corrected}\PYG{p}{[}\PYG{n}{row}\PYG{p}{,} \PYG{l+m+mi}{1}\PYG{p}{:}\PYG{l+m+mi}{3}\PYG{p}{]}
        \PYG{n}{p2} \PYG{o}{=} \PYG{n}{corrected}\PYG{p}{[}\PYG{n}{row} \PYG{o}{+} \PYG{n}{data\PYGZus{}rate}\PYG{p}{,} \PYG{l+m+mi}{1}\PYG{p}{:}\PYG{l+m+mi}{3}\PYG{p}{]}
        \PYG{k}{if} \PYG{n}{np}\PYG{o}{.}\PYG{n}{any}\PYG{p}{(}\PYG{n}{np}\PYG{o}{.}\PYG{n}{isnan}\PYG{p}{(}\PYG{n}{p1}\PYG{p}{)}\PYG{p}{)} \PYG{o+ow}{or} \PYG{n}{np}\PYG{o}{.}\PYG{n}{any}\PYG{p}{(}\PYG{n}{np}\PYG{o}{.}\PYG{n}{isnan}\PYG{p}{(}\PYG{n}{p2}\PYG{p}{)}\PYG{p}{)}\PYG{p}{:}
            \PYG{n}{velocity}\PYG{p}{[}\PYG{n}{bin\PYGZus{}idx}\PYG{p}{]} \PYG{o}{=} \PYG{n}{np}\PYG{o}{.}\PYG{n}{nan}
        \PYG{k}{else}\PYG{p}{:}
            \PYG{c+c1}{\PYGZsh{} 10x converts cm to mm}
            \PYG{n}{velocity}\PYG{p}{[}\PYG{n}{bin\PYGZus{}idx}\PYG{p}{]} \PYG{o}{=} \PYG{l+m+mf}{10.0} \PYG{o}{*} \PYG{n}{euclidean}\PYG{p}{(}\PYG{n}{p1}\PYG{p}{,} \PYG{n}{p2}\PYG{p}{)}

    \PYG{c+c1}{\PYGZsh{} Convert from mm/bin to mm/s}
    \PYG{n}{velocity} \PYG{o}{=} \PYG{n}{velocity} \PYG{o}{/} \PYG{n}{bin\PYGZus{}size}

    \PYG{n}{all\PYGZus{}velocity}\PYG{o}{.}\PYG{n}{append}\PYG{p}{(}\PYG{n}{velocity}\PYG{p}{)}

    \PYG{c+c1}{\PYGZsh{} Warn about absurd velocities}
    \PYG{k}{if} \PYG{n}{np}\PYG{o}{.}\PYG{n}{nanmax}\PYG{p}{(}\PYG{n}{velocity}\PYG{p}{)} \PYG{o}{\PYGZgt{}} \PYG{l+m+mi}{30}\PYG{p}{:}
        \PYG{n+nb}{print}\PYG{p}{(}\PYG{l+s+sa}{f}\PYG{l+s+s1}{\PYGZsq{}}\PYG{l+s+s1}{WARNING (Video }\PYG{l+s+si}{\PYGZob{}}\PYG{n}{idx}\PYG{+w}{ }\PYG{o}{+}\PYG{+w}{ }\PYG{l+m+mi}{1}\PYG{l+s+si}{\PYGZcb{}}\PYG{l+s+s1}{): Absurdly high velocities detected.}\PYG{l+s+s1}{\PYGZsq{}}\PYG{p}{)}
        \PYG{n+nb}{print}\PYG{p}{(}\PYG{l+s+s1}{\PYGZsq{}}\PYG{l+s+s1}{  Consider changing the ROI or re\PYGZhy{}recording the video.}\PYG{l+s+s1}{\PYGZsq{}}\PYG{p}{)}

    \PYG{n}{mean\PYGZus{}vel} \PYG{o}{=} \PYG{n}{np}\PYG{o}{.}\PYG{n}{nanmean}\PYG{p}{(}\PYG{n}{velocity}\PYG{p}{)}
    \PYG{n}{std\PYGZus{}vel} \PYG{o}{=} \PYG{n}{np}\PYG{o}{.}\PYG{n}{nanstd}\PYG{p}{(}\PYG{n}{velocity}\PYG{p}{)}
    \PYG{n+nb}{print}\PYG{p}{(}\PYG{l+s+sa}{f}\PYG{l+s+s1}{\PYGZsq{}}\PYG{l+s+se}{\PYGZbs{}n}\PYG{l+s+s1}{\PYGZhy{}\PYGZhy{}\PYGZhy{} Video }\PYG{l+s+si}{\PYGZob{}}\PYG{n}{idx}\PYG{+w}{ }\PYG{o}{+}\PYG{+w}{ }\PYG{l+m+mi}{1}\PYG{l+s+si}{\PYGZcb{}}\PYG{l+s+s1}{ (from: }\PYG{l+s+si}{\PYGZob{}}\PYG{n}{os}\PYG{o}{.}\PYG{n}{path}\PYG{o}{.}\PYG{n}{basename}\PYG{p}{(}\PYG{n}{current\PYGZus{}video}\PYG{p}{)}\PYG{l+s+si}{\PYGZcb{}}\PYG{l+s+s1}{) \PYGZhy{}\PYGZhy{}\PYGZhy{}}\PYG{l+s+s1}{\PYGZsq{}}\PYG{p}{)}
    \PYG{n+nb}{print}\PYG{p}{(}\PYG{l+s+sa}{f}\PYG{l+s+s1}{\PYGZsq{}}\PYG{l+s+s1}{  Mean velocity: }\PYG{l+s+si}{\PYGZob{}}\PYG{n}{mean\PYGZus{}vel}\PYG{l+s+si}{:}\PYG{l+s+s1}{.2f}\PYG{l+s+si}{\PYGZcb{}}\PYG{l+s+s1}{ mm/s}\PYG{l+s+s1}{\PYGZsq{}}\PYG{p}{)}
    \PYG{n+nb}{print}\PYG{p}{(}\PYG{l+s+sa}{f}\PYG{l+s+s1}{\PYGZsq{}}\PYG{l+s+s1}{  Std deviation: }\PYG{l+s+si}{\PYGZob{}}\PYG{n}{std\PYGZus{}vel}\PYG{l+s+si}{:}\PYG{l+s+s1}{.2f}\PYG{l+s+si}{\PYGZcb{}}\PYG{l+s+s1}{ mm/s}\PYG{l+s+s1}{\PYGZsq{}}\PYG{p}{)}

    \PYG{c+c1}{\PYGZsh{} Plot velocity over time}
    \PYG{n}{time\PYGZus{}axis} \PYG{o}{=} \PYG{n}{np}\PYG{o}{.}\PYG{n}{arange}\PYG{p}{(}\PYG{n}{total\PYGZus{}bins}\PYG{p}{)} \PYG{o}{*} \PYG{n}{bin\PYGZus{}size}
    \PYG{n}{fig}\PYG{p}{,} \PYG{n}{ax} \PYG{o}{=} \PYG{n}{plt}\PYG{o}{.}\PYG{n}{subplots}\PYG{p}{(}\PYG{n}{figsize}\PYG{o}{=}\PYG{p}{(}\PYG{l+m+mi}{10}\PYG{p}{,} \PYG{l+m+mi}{4}\PYG{p}{)}\PYG{p}{)}
    \PYG{n}{ax}\PYG{o}{.}\PYG{n}{plot}\PYG{p}{(}\PYG{n}{time\PYGZus{}axis}\PYG{p}{,} \PYG{n}{velocity}\PYG{p}{,} \PYG{n}{linewidth}\PYG{o}{=}\PYG{l+m+mf}{1.5}\PYG{p}{)}
    \PYG{n}{ax}\PYG{o}{.}\PYG{n}{set\PYGZus{}xlim}\PYG{p}{(}\PYG{l+m+mi}{0}\PYG{p}{,} \PYG{n}{time\PYGZus{}axis}\PYG{p}{[}\PYG{o}{\PYGZhy{}}\PYG{l+m+mi}{1}\PYG{p}{]} \PYG{o}{+} \PYG{n}{bin\PYGZus{}size} \PYG{k}{if} \PYG{n+nb}{len}\PYG{p}{(}\PYG{n}{time\PYGZus{}axis}\PYG{p}{)} \PYG{o}{\PYGZgt{}} \PYG{l+m+mi}{0} \PYG{k}{else} \PYG{l+m+mi}{1}\PYG{p}{)}
    \PYG{n}{ax}\PYG{o}{.}\PYG{n}{set\PYGZus{}ylim}\PYG{p}{(}\PYG{l+m+mi}{0}\PYG{p}{,} \PYG{n}{np}\PYG{o}{.}\PYG{n}{nanmax}\PYG{p}{(}\PYG{n}{velocity}\PYG{p}{)} \PYG{o}{*} \PYG{l+m+mf}{1.5} \PYG{k}{if} \PYG{n}{np}\PYG{o}{.}\PYG{n}{nanmax}\PYG{p}{(}\PYG{n}{velocity}\PYG{p}{)} \PYG{o}{\PYGZgt{}} \PYG{l+m+mi}{0} \PYG{k}{else} \PYG{l+m+mi}{1}\PYG{p}{)}
    \PYG{n}{ax}\PYG{o}{.}\PYG{n}{set\PYGZus{}xlabel}\PYG{p}{(}\PYG{l+s+s1}{\PYGZsq{}}\PYG{l+s+s1}{Time (s)}\PYG{l+s+s1}{\PYGZsq{}}\PYG{p}{,} \PYG{n}{fontsize}\PYG{o}{=}\PYG{l+m+mi}{11}\PYG{p}{)}
    \PYG{n}{ax}\PYG{o}{.}\PYG{n}{set\PYGZus{}ylabel}\PYG{p}{(}\PYG{l+s+s1}{\PYGZsq{}}\PYG{l+s+s1}{Velocity (mm/s)}\PYG{l+s+s1}{\PYGZsq{}}\PYG{p}{,} \PYG{n}{fontsize}\PYG{o}{=}\PYG{l+m+mi}{11}\PYG{p}{)}
    \PYG{n}{ax}\PYG{o}{.}\PYG{n}{set\PYGZus{}title}\PYG{p}{(}\PYG{l+s+sa}{f}\PYG{l+s+s1}{\PYGZsq{}}\PYG{l+s+s1}{Fly Velocity — Video }\PYG{l+s+si}{\PYGZob{}}\PYG{n}{idx}\PYG{+w}{ }\PYG{o}{+}\PYG{+w}{ }\PYG{l+m+mi}{1}\PYG{l+s+si}{\PYGZcb{}}\PYG{l+s+s1}{\PYGZsq{}}\PYG{p}{)}
    \PYG{n}{plt}\PYG{o}{.}\PYG{n}{tight\PYGZus{}layout}\PYG{p}{(}\PYG{p}{)}
    \PYG{n}{plt}\PYG{o}{.}\PYG{n}{show}\PYG{p}{(}\PYG{p}{)}

\PYG{n+nb}{print}\PYG{p}{(}\PYG{l+s+sa}{f}\PYG{l+s+s1}{\PYGZsq{}}\PYG{l+s+se}{\PYGZbs{}n}\PYG{l+s+s1}{Velocity analysis complete. Total videos analyzed: }\PYG{l+s+si}{\PYGZob{}}\PYG{n+nb}{len}\PYG{p}{(}\PYG{n}{all\PYGZus{}velocity}\PYG{p}{)}\PYG{l+s+si}{\PYGZcb{}}\PYG{l+s+s1}{\PYGZsq{}}\PYG{p}{)}
\end{sphinxVerbatim}

\end{sphinxuseclass}\end{sphinxVerbatimInput}

\end{sphinxuseclass}

\subsection{Summary Across All Videos}
\label{\detokenize{FlyTracking/ColabNotebook:summary-across-all-videos}}
\begin{sphinxuseclass}{cell}\begin{sphinxVerbatimInput}

\begin{sphinxuseclass}{cell_input}
\begin{sphinxVerbatim}[commandchars=\\\{\}]
\PYG{n}{num\PYGZus{}files} \PYG{o}{=} \PYG{n+nb}{len}\PYG{p}{(}\PYG{n}{all\PYGZus{}velocity}\PYG{p}{)}

\PYG{k}{if} \PYG{n}{num\PYGZus{}files} \PYG{o}{\PYGZgt{}} \PYG{l+m+mi}{1}\PYG{p}{:}
    \PYG{c+c1}{\PYGZsh{} Pad velocity arrays to the same length for comparison}
    \PYG{n}{max\PYGZus{}len} \PYG{o}{=} \PYG{n+nb}{max}\PYG{p}{(}\PYG{n+nb}{len}\PYG{p}{(}\PYG{n}{v}\PYG{p}{)} \PYG{k}{for} \PYG{n}{v} \PYG{o+ow}{in} \PYG{n}{all\PYGZus{}velocity}\PYG{p}{)}
    \PYG{n}{velocity\PYGZus{}matrix} \PYG{o}{=} \PYG{n}{np}\PYG{o}{.}\PYG{n}{full}\PYG{p}{(}\PYG{p}{(}\PYG{n}{num\PYGZus{}files}\PYG{p}{,} \PYG{n}{max\PYGZus{}len}\PYG{p}{)}\PYG{p}{,} \PYG{n}{np}\PYG{o}{.}\PYG{n}{nan}\PYG{p}{)}
    \PYG{k}{for} \PYG{n}{i}\PYG{p}{,} \PYG{n}{v} \PYG{o+ow}{in} \PYG{n+nb}{enumerate}\PYG{p}{(}\PYG{n}{all\PYGZus{}velocity}\PYG{p}{)}\PYG{p}{:}
        \PYG{n}{velocity\PYGZus{}matrix}\PYG{p}{[}\PYG{n}{i}\PYG{p}{,} \PYG{p}{:}\PYG{n+nb}{len}\PYG{p}{(}\PYG{n}{v}\PYG{p}{)}\PYG{p}{]} \PYG{o}{=} \PYG{n}{v}

    \PYG{c+c1}{\PYGZsh{} Plot all velocities together}
    \PYG{n}{fig}\PYG{p}{,} \PYG{n}{ax} \PYG{o}{=} \PYG{n}{plt}\PYG{o}{.}\PYG{n}{subplots}\PYG{p}{(}\PYG{n}{figsize}\PYG{o}{=}\PYG{p}{(}\PYG{l+m+mi}{10}\PYG{p}{,} \PYG{l+m+mi}{5}\PYG{p}{)}\PYG{p}{)}
    \PYG{n}{time\PYGZus{}axis} \PYG{o}{=} \PYG{n}{np}\PYG{o}{.}\PYG{n}{arange}\PYG{p}{(}\PYG{n}{max\PYGZus{}len}\PYG{p}{)} \PYG{o}{*} \PYG{n}{bin\PYGZus{}size}
    \PYG{k}{for} \PYG{n}{i} \PYG{o+ow}{in} \PYG{n+nb}{range}\PYG{p}{(}\PYG{n}{num\PYGZus{}files}\PYG{p}{)}\PYG{p}{:}
        \PYG{n}{ax}\PYG{o}{.}\PYG{n}{plot}\PYG{p}{(}\PYG{n}{time\PYGZus{}axis}\PYG{p}{,} \PYG{n}{velocity\PYGZus{}matrix}\PYG{p}{[}\PYG{n}{i}\PYG{p}{]}\PYG{p}{,} \PYG{n}{linewidth}\PYG{o}{=}\PYG{l+m+mi}{2}\PYG{p}{,}
                \PYG{n}{label}\PYG{o}{=}\PYG{l+s+sa}{f}\PYG{l+s+s1}{\PYGZsq{}}\PYG{l+s+s1}{Fly }\PYG{l+s+si}{\PYGZob{}}\PYG{n}{i}\PYG{+w}{ }\PYG{o}{+}\PYG{+w}{ }\PYG{l+m+mi}{1}\PYG{l+s+si}{\PYGZcb{}}\PYG{l+s+s1}{\PYGZsq{}}\PYG{p}{)}
    \PYG{n}{ax}\PYG{o}{.}\PYG{n}{set\PYGZus{}xlabel}\PYG{p}{(}\PYG{l+s+s1}{\PYGZsq{}}\PYG{l+s+s1}{Time (s)}\PYG{l+s+s1}{\PYGZsq{}}\PYG{p}{,} \PYG{n}{fontsize}\PYG{o}{=}\PYG{l+m+mi}{11}\PYG{p}{)}
    \PYG{n}{ax}\PYG{o}{.}\PYG{n}{set\PYGZus{}ylabel}\PYG{p}{(}\PYG{l+s+s1}{\PYGZsq{}}\PYG{l+s+s1}{Velocity (mm/s)}\PYG{l+s+s1}{\PYGZsq{}}\PYG{p}{,} \PYG{n}{fontsize}\PYG{o}{=}\PYG{l+m+mi}{11}\PYG{p}{)}
    \PYG{n}{ax}\PYG{o}{.}\PYG{n}{set\PYGZus{}title}\PYG{p}{(}\PYG{l+s+s1}{\PYGZsq{}}\PYG{l+s+s1}{All Fly Velocities}\PYG{l+s+s1}{\PYGZsq{}}\PYG{p}{)}
    \PYG{n}{ax}\PYG{o}{.}\PYG{n}{legend}\PYG{p}{(}\PYG{p}{)}
    \PYG{n}{plt}\PYG{o}{.}\PYG{n}{tight\PYGZus{}layout}\PYG{p}{(}\PYG{p}{)}
    \PYG{n}{plt}\PYG{o}{.}\PYG{n}{show}\PYG{p}{(}\PYG{p}{)}

    \PYG{c+c1}{\PYGZsh{} Summary stats across videos}
    \PYG{n}{per\PYGZus{}video\PYGZus{}means} \PYG{o}{=} \PYG{p}{[}\PYG{n}{np}\PYG{o}{.}\PYG{n}{nanmean}\PYG{p}{(}\PYG{n}{v}\PYG{p}{)} \PYG{k}{for} \PYG{n}{v} \PYG{o+ow}{in} \PYG{n}{all\PYGZus{}velocity}\PYG{p}{]}
    \PYG{n}{mean\PYGZus{}across} \PYG{o}{=} \PYG{n}{np}\PYG{o}{.}\PYG{n}{mean}\PYG{p}{(}\PYG{n}{per\PYGZus{}video\PYGZus{}means}\PYG{p}{)}
    \PYG{n}{sd\PYGZus{}across} \PYG{o}{=} \PYG{n}{np}\PYG{o}{.}\PYG{n}{std}\PYG{p}{(}\PYG{n}{per\PYGZus{}video\PYGZus{}means}\PYG{p}{)}
    \PYG{n+nb}{print}\PYG{p}{(}\PYG{l+s+sa}{f}\PYG{l+s+s1}{\PYGZsq{}}\PYG{l+s+se}{\PYGZbs{}n}\PYG{l+s+s1}{=== Summary Across }\PYG{l+s+si}{\PYGZob{}}\PYG{n}{num\PYGZus{}files}\PYG{l+s+si}{\PYGZcb{}}\PYG{l+s+s1}{ Videos ===}\PYG{l+s+s1}{\PYGZsq{}}\PYG{p}{)}
    \PYG{n+nb}{print}\PYG{p}{(}\PYG{l+s+sa}{f}\PYG{l+s+s1}{\PYGZsq{}}\PYG{l+s+s1}{Mean velocity across videos: }\PYG{l+s+si}{\PYGZob{}}\PYG{n}{mean\PYGZus{}across}\PYG{l+s+si}{:}\PYG{l+s+s1}{.2f}\PYG{l+s+si}{\PYGZcb{}}\PYG{l+s+s1}{ mm/s}\PYG{l+s+s1}{\PYGZsq{}}\PYG{p}{)}
    \PYG{n+nb}{print}\PYG{p}{(}\PYG{l+s+sa}{f}\PYG{l+s+s1}{\PYGZsq{}}\PYG{l+s+s1}{SD of mean velocity across videos: }\PYG{l+s+si}{\PYGZob{}}\PYG{n}{sd\PYGZus{}across}\PYG{l+s+si}{:}\PYG{l+s+s1}{.2f}\PYG{l+s+si}{\PYGZcb{}}\PYG{l+s+s1}{ mm/s}\PYG{l+s+s1}{\PYGZsq{}}\PYG{p}{)}
\PYG{k}{else}\PYG{p}{:}
    \PYG{n}{mean\PYGZus{}vel} \PYG{o}{=} \PYG{n}{np}\PYG{o}{.}\PYG{n}{nanmean}\PYG{p}{(}\PYG{n}{all\PYGZus{}velocity}\PYG{p}{[}\PYG{l+m+mi}{0}\PYG{p}{]}\PYG{p}{)}
    \PYG{n}{sd\PYGZus{}vel} \PYG{o}{=} \PYG{n}{np}\PYG{o}{.}\PYG{n}{nanstd}\PYG{p}{(}\PYG{n}{all\PYGZus{}velocity}\PYG{p}{[}\PYG{l+m+mi}{0}\PYG{p}{]}\PYG{p}{)}
    \PYG{n+nb}{print}\PYG{p}{(}\PYG{l+s+sa}{f}\PYG{l+s+s1}{\PYGZsq{}}\PYG{l+s+se}{\PYGZbs{}n}\PYG{l+s+s1}{=== Summary (1 Video) ===}\PYG{l+s+s1}{\PYGZsq{}}\PYG{p}{)}
    \PYG{n+nb}{print}\PYG{p}{(}\PYG{l+s+sa}{f}\PYG{l+s+s1}{\PYGZsq{}}\PYG{l+s+s1}{Mean velocity: }\PYG{l+s+si}{\PYGZob{}}\PYG{n}{mean\PYGZus{}vel}\PYG{l+s+si}{:}\PYG{l+s+s1}{.2f}\PYG{l+s+si}{\PYGZcb{}}\PYG{l+s+s1}{ mm/s}\PYG{l+s+s1}{\PYGZsq{}}\PYG{p}{)}
    \PYG{n+nb}{print}\PYG{p}{(}\PYG{l+s+sa}{f}\PYG{l+s+s1}{\PYGZsq{}}\PYG{l+s+s1}{SD of velocity: }\PYG{l+s+si}{\PYGZob{}}\PYG{n}{sd\PYGZus{}vel}\PYG{l+s+si}{:}\PYG{l+s+s1}{.2f}\PYG{l+s+si}{\PYGZcb{}}\PYG{l+s+s1}{ mm/s}\PYG{l+s+s1}{\PYGZsq{}}\PYG{p}{)}
\end{sphinxVerbatim}

\end{sphinxuseclass}\end{sphinxVerbatimInput}

\end{sphinxuseclass}

\subsection{Prepare for Another Video (Optional)}
\label{\detokenize{FlyTracking/ColabNotebook:prepare-for-another-video-optional}}
\sphinxAtStartPar
If you want to process another video with the same ROI:
\begin{enumerate}
\sphinxsetlistlabels{\arabic}{enumi}{enumii}{}{.}%
\item {} 
\sphinxAtStartPar
Go back to the \sphinxstylestrong{Select Video File} cell and update \sphinxcode{\sphinxupquote{current\_video}} to the new file path

\item {} 
\sphinxAtStartPar
Re\sphinxhyphen{}run the video selection cell

\item {} 
\sphinxAtStartPar
Skip the ROI selection (you can reuse the same ROI with \sphinxcode{\sphinxupquote{roi = roi}} or draw a new one)

\item {} 
\sphinxAtStartPar
Run the \sphinxstylestrong{Track Fly} cell again — results automatically append to \sphinxcode{\sphinxupquote{all\_corrected}}

\end{enumerate}

\sphinxAtStartPar
Or, to apply a \sphinxstylestrong{different ROI} for the next video, re\sphinxhyphen{}run the ROI selection cells before tracking.


\subsection{Download Results (Optional)}
\label{\detokenize{FlyTracking/ColabNotebook:download-results-optional}}
\sphinxAtStartPar
Download the tracking data as CSV files.

\begin{sphinxuseclass}{cell}\begin{sphinxVerbatimInput}

\begin{sphinxuseclass}{cell_input}
\begin{sphinxVerbatim}[commandchars=\\\{\}]
\PYG{k}{for} \PYG{n}{idx}\PYG{p}{,} \PYG{p}{(}\PYG{n}{corrected}\PYG{p}{,} \PYG{n}{vel}\PYG{p}{)} \PYG{o+ow}{in} \PYG{n+nb}{enumerate}\PYG{p}{(}\PYG{n+nb}{zip}\PYG{p}{(}\PYG{n}{all\PYGZus{}corrected}\PYG{p}{,} \PYG{n}{all\PYGZus{}velocity}\PYG{p}{)}\PYG{p}{)}\PYG{p}{:}
    \PYG{n}{base} \PYG{o}{=} \PYG{n}{os}\PYG{o}{.}\PYG{n}{path}\PYG{o}{.}\PYG{n}{splitext}\PYG{p}{(}\PYG{n}{os}\PYG{o}{.}\PYG{n}{path}\PYG{o}{.}\PYG{n}{basename}\PYG{p}{(}\PYG{n}{all\PYGZus{}video\PYGZus{}paths}\PYG{p}{[}\PYG{n}{idx}\PYG{p}{]}\PYG{p}{)}\PYG{p}{)}\PYG{p}{[}\PYG{l+m+mi}{0}\PYG{p}{]}

    \PYG{n}{csv\PYGZus{}name} \PYG{o}{=} \PYG{l+s+sa}{f}\PYG{l+s+s1}{\PYGZsq{}}\PYG{l+s+si}{\PYGZob{}}\PYG{n}{base}\PYG{l+s+si}{\PYGZcb{}}\PYG{l+s+s1}{\PYGZus{}tracking.csv}\PYG{l+s+s1}{\PYGZsq{}}
    \PYG{n}{np}\PYG{o}{.}\PYG{n}{savetxt}\PYG{p}{(}\PYG{n}{csv\PYGZus{}name}\PYG{p}{,} \PYG{n}{corrected}\PYG{p}{,} \PYG{n}{delimiter}\PYG{o}{=}\PYG{l+s+s1}{\PYGZsq{}}\PYG{l+s+s1}{,}\PYG{l+s+s1}{\PYGZsq{}}\PYG{p}{,}
               \PYG{n}{header}\PYG{o}{=}\PYG{l+s+s1}{\PYGZsq{}}\PYG{l+s+s1}{Time\PYGZus{}s,X\PYGZus{}cm,Y\PYGZus{}cm}\PYG{l+s+s1}{\PYGZsq{}}\PYG{p}{,} \PYG{n}{comments}\PYG{o}{=}\PYG{l+s+s1}{\PYGZsq{}}\PYG{l+s+s1}{\PYGZsq{}}\PYG{p}{)}
    \PYG{n}{files}\PYG{o}{.}\PYG{n}{download}\PYG{p}{(}\PYG{n}{csv\PYGZus{}name}\PYG{p}{)}
    \PYG{n+nb}{print}\PYG{p}{(}\PYG{l+s+sa}{f}\PYG{l+s+s1}{\PYGZsq{}}\PYG{l+s+s1}{Downloaded }\PYG{l+s+si}{\PYGZob{}}\PYG{n}{csv\PYGZus{}name}\PYG{l+s+si}{\PYGZcb{}}\PYG{l+s+s1}{\PYGZsq{}}\PYG{p}{)}

    \PYG{n}{time\PYGZus{}axis} \PYG{o}{=} \PYG{n}{np}\PYG{o}{.}\PYG{n}{arange}\PYG{p}{(}\PYG{n+nb}{len}\PYG{p}{(}\PYG{n}{vel}\PYG{p}{)}\PYG{p}{)} \PYG{o}{*} \PYG{n}{bin\PYGZus{}size}
    \PYG{n}{vel\PYGZus{}csv} \PYG{o}{=} \PYG{l+s+sa}{f}\PYG{l+s+s1}{\PYGZsq{}}\PYG{l+s+si}{\PYGZob{}}\PYG{n}{base}\PYG{l+s+si}{\PYGZcb{}}\PYG{l+s+s1}{\PYGZus{}velocity.csv}\PYG{l+s+s1}{\PYGZsq{}}
    \PYG{n}{np}\PYG{o}{.}\PYG{n}{savetxt}\PYG{p}{(}\PYG{n}{vel\PYGZus{}csv}\PYG{p}{,} \PYG{n}{np}\PYG{o}{.}\PYG{n}{column\PYGZus{}stack}\PYG{p}{(}\PYG{p}{[}\PYG{n}{time\PYGZus{}axis}\PYG{p}{,} \PYG{n}{vel}\PYG{p}{]}\PYG{p}{)}\PYG{p}{,} \PYG{n}{delimiter}\PYG{o}{=}\PYG{l+s+s1}{\PYGZsq{}}\PYG{l+s+s1}{,}\PYG{l+s+s1}{\PYGZsq{}}\PYG{p}{,}
               \PYG{n}{header}\PYG{o}{=}\PYG{l+s+s1}{\PYGZsq{}}\PYG{l+s+s1}{Time\PYGZus{}s,Velocity\PYGZus{}mm\PYGZus{}per\PYGZus{}s}\PYG{l+s+s1}{\PYGZsq{}}\PYG{p}{,} \PYG{n}{comments}\PYG{o}{=}\PYG{l+s+s1}{\PYGZsq{}}\PYG{l+s+s1}{\PYGZsq{}}\PYG{p}{)}
    \PYG{n}{files}\PYG{o}{.}\PYG{n}{download}\PYG{p}{(}\PYG{n}{vel\PYGZus{}csv}\PYG{p}{)}
    \PYG{n+nb}{print}\PYG{p}{(}\PYG{l+s+sa}{f}\PYG{l+s+s1}{\PYGZsq{}}\PYG{l+s+s1}{Downloaded }\PYG{l+s+si}{\PYGZob{}}\PYG{n}{vel\PYGZus{}csv}\PYG{l+s+si}{\PYGZcb{}}\PYG{l+s+s1}{\PYGZsq{}}\PYG{p}{)}
\end{sphinxVerbatim}

\end{sphinxuseclass}\end{sphinxVerbatimInput}

\end{sphinxuseclass}
\sphinxstepscope


\chapter{Measuring Alpha Wave Activity Using Electroencephalography}
\label{\detokenize{EEG:measuring-alpha-wave-activity-using-electroencephalography}}\label{\detokenize{EEG::doc}}

\section{Background \& Goals}
\label{\detokenize{EEG:background-goals}}
\sphinxAtStartPar
The cerebral cortex contains huge numbers of neurons. Activity of these neurons is — to some extent — synchronized in regular firing rhythms. These are referred to as brain waves. Electrodes placed in pairs on the scalp can pick up variations in electrical potential that derive from this underlying cortical activity. This recording of the electrical activity in the cortex is called an electroencephalogram (EEG). EEG signals are strongly affected by different states of arousal and show characteristic changes in different stages of sleep. EEG signals are also affected by stimulation from the external environment and brain waves can become entrained to external stimuli. Electroencephalography is used, among other things, in the diagnosis of epilepsy and the diagnosis of brain death.


\subsection{Recording EEG signals}
\label{\detokenize{EEG:recording-eeg-signals}}
\sphinxAtStartPar
EEG recording is technically difficult, mainly because of the small size of the voltage signals, which are typically 50 µV \sphinxstylestrong{peak\sphinxhyphen{}to\sphinxhyphen{}peak}. The signals are small because the recording electrodes are separated from the brain’s surface by the scalp, the skull, and a layer of cerebrospinal fluid. A specially designed amplifier, such as the Bio Amp built into the PowerLab, is essential to record EEGs. It is also important to use electrodes made of the right material and to connect them properly. Even with these precautions, recordings may be spoiled by a range of unwanted interfering influences, known as artifacts.

\sphinxAtStartPar
In this lab, you will record EEG activity with an occipital electrode on the scalp at the back of the head. A third (ground or earth) electrode is also attached, to reduce electrical interference. In clinical EEG, it is usual to record many channels of activity from multiple recording electrodes placed in an array over the head.


\subsection{Origins of the EEG signals}
\label{\detokenize{EEG:origins-of-the-eeg-signals}}
\sphinxAtStartPar
The EEG results from slow changes in the membrane potentials of cortical neurons, especially the excitatory and inhibitory post\sphinxhyphen{}synaptic potentials (EPSPs and IPSPs). Very little contribution normally comes from action potentials propagated along nerve axons. The EEG reflects the sum of the electrical potential changes occurring from large populations of cells. Therefore, large amplitude waves require the synchronous activity of a large number of neurons. The rhythmic events that these waves reflect often arise in the thalamus whose activity is in turn affected by a variety of inputs including structures in the brainstem reticular formation.


\subsection{Components of the EEG waveform}
\label{\detokenize{EEG:components-of-the-eeg-waveform}}
\sphinxAtStartPar
The EEG waveform contains component waves of different frequencies. These can be extracted and provide information about different brain activities. The types of brain waves are:
\begin{itemize}
\item {} 
\sphinxAtStartPar
\sphinxstylestrong{alpha} (between 8 to 13 Hz; average amplitudes of 30 to 50 µV peak\sphinxhyphen{}to\sphinxhyphen{}peak), which will be studied in this experiment. Alpha rhythm is seen when the eyes are closed and the volunteer relaxed. It is abolished by eye opening and by mental effort such as doing calculations or concentrating on an idea. It is thus thought to indicate the degree of cortical activation. The greater the activation, the lower the alpha activity. Alpha waves are strongest over the occipital (back of the head) cortex and also over the frontal cortex.

\item {} 
\sphinxAtStartPar
\sphinxstylestrong{beta} (13 to 30 Hz; <20 µV peak\sphinxhyphen{}to\sphinxhyphen{}peak), which are prominent in alert individuals with their eyes open. The beta rhythm may be absent or reduced in areas of cortical damage and can be accentuated by sedative\sphinxhyphen{}hypnotic drugs such as benzodiazepines and barbiturates.

\item {} 
\sphinxAtStartPar
\sphinxstylestrong{theta} (4 to 8 Hz; <30 µV peak\sphinxhyphen{}to\sphinxhyphen{}peak), which are seen in awake children but not adults. The theta rhythm is normal during sleep at all ages. However, some researchers separate this frequency band into two components: low theta (4–5.45 Hz) activity correlated with decreased arousal and increased drowsiness, and high theta (6–7.45 Hz) activity that is claimed to be enhanced during tasks involving working memory.

\item {} 
\sphinxAtStartPar
\sphinxstylestrong{delta} (0.5 to 4 Hz; up to 100–200 µV peak\sphinxhyphen{}to\sphinxhyphen{}peak), which is the dominant rhythm in sleep stages three and four but not seen in conscious adults. The delta rhythm tends to have the highest amplitude of any of the component EEG waves. EEG artifacts caused by movements of jaw and neck muscles can produce waves in the same frequency band.

\item {} 
\sphinxAtStartPar
\sphinxstylestrong{gamma} (30 to 50 Hz). Most people recognize gamma rhythm, but its importance is controversial. It may be associated with higher mental activity, including perception and consciousness, and it disappears under general anesthesia. One suggestion is that the gamma rhythm reflects the mental activity involved in integrating various aspects of an object (color, shape, movement, etc.) to form a coherent picture.

\end{itemize}

\sphinxAtStartPar
\sphinxstyleemphasis{Written by staff of ADInstruments \& A. Juavinett. Copyright © 2015 ADInstruments Pty Ltd. All rights reserved. PowerLab® and LabChart® are registered trademarks of ADInstruments Pty Ltd.}


\bigskip\hrule\bigskip



\section{Lab Protocol}
\label{\detokenize{EEG:lab-protocol}}

\begin{savenotes}\sphinxattablestart
\sphinxthistablewithglobalstyle
\centering
\begin{tabulary}{\linewidth}[t]{T}
\sphinxtoprule
\sphinxstyletheadfamily 
\sphinxAtStartPar
\sphinxstylestrong{Supplies}
\\
\sphinxmidrule
\sphinxtableatstartofbodyhook
\sphinxAtStartPar
LabChart software \& computer
\\
\sphinxhline
\sphinxAtStartPar
PowerLab 26T
\\
\sphinxhline
\sphinxAtStartPar
5 Lead Shielded Bio Amp Cable
\\
\sphinxhline
\sphinxAtStartPar
EEG Flat Electrodes (gold contacts)
\\
\sphinxhline
\sphinxAtStartPar
Electrode Paste
\\
\sphinxhline
\sphinxAtStartPar
Abrasive Gel (NuPrep)
\\
\sphinxhline
\sphinxAtStartPar
Alcohol Swabs
\\
\sphinxhline
\sphinxAtStartPar
Medical tape
\\
\sphinxhline
\sphinxAtStartPar
Elastic bandage
\\
\sphinxhline
\sphinxAtStartPar
Cotton swab
\\
\sphinxbottomrule
\end{tabulary}
\sphinxtableafterendhook\par
\sphinxattableend\end{savenotes}


\bigskip\hrule\bigskip



\subsection{Part I. Spectral Analysis Tutorial}
\label{\detokenize{EEG:part-i-spectral-analysis-tutorial}}
\sphinxAtStartPar
A spectrum is a representation of data based on the frequency distribution of its component sine waves. Spectra indicate the strength of the various frequencies in a time\sphinxhyphen{}varying waveform. Spectrum View allows you to observe the frequency distribution of data that might not otherwise be easily seen. For example, it could be used to break down an EEG waveform into its various components: beta waves, alpha waves, theta waves and delta waves. A mathematical technique known as the Fast Fourier Transform is applied to the raw data. The results of this analysis can be displayed as a plot of the power (vertical axis) of different frequencies (horizontal axis) relative to each other in the input signal. This is called a \sphinxstylestrong{Power Spectrum Density (PSD) plot}. The data can also be displayed as a 3\sphinxhyphen{}dimensional color plot of spectral power, frequency, and time, called a \sphinxstylestrong{Spectrogram}.

\sphinxAtStartPar
\sphinxstylestrong{Note:} The questions here are written to ensure that you are understanding the steps of the tutorial. You do not need to answer them, and they do not directly relate to your lab report. You \sphinxstyleemphasis{should}, however, understand how to interpret a spectrogram.
\begin{enumerate}
\sphinxsetlistlabels{\arabic}{enumi}{enumii}{}{.}%
\item {} 
\sphinxAtStartPar
Open the EEG Spectral Analysis Tutorial settings file (\sphinxurl{http://bit.ly/labchart}).

\item {} 
\sphinxAtStartPar
Examine the Chart View. Use the \sphinxstylestrong{View Buttons} to view each block. You should see five blocks of data. The first record is a slowly oscillating sine wave.

\item {} 
\sphinxAtStartPar
Open Spectrum view by clicking on the Spectrum View button in the Toolbar:

\sphinxAtStartPar
\sphinxincludegraphics{{eeg_spectrum_toolbar}.png}

\item {} 
\sphinxAtStartPar
Click the \sphinxstylestrong{Smart Tile} button in the LabChart Toolbar to display both windows in full screen mode.

\item {} 
\sphinxAtStartPar
In Chart View, \sphinxstylestrong{select} the first record by double\sphinxhyphen{}clicking in the Time axis. This will perform a spectral analysis for this record and display the result in the Spectrum view.

\item {} 
\sphinxAtStartPar
Adjust the horizontal scaling of plots to view the results:
\begin{itemize}
\item {} 
\sphinxAtStartPar
Set the \sphinxstylestrong{horizontal scaling} for the \sphinxstylestrong{Power Spectrum Density (PSD)} plot to 50 Hz.

\item {} 
\sphinxAtStartPar
Use the horizontal scroll bar to display the 0 Hz to 50 Hz region of the plot.

\item {} 
\sphinxAtStartPar
You may need to Autoscale the axes.

\item {} 
\sphinxAtStartPar
Set the \sphinxstylestrong{horizontal scaling} for the Spectrogram to 50:1.

\end{itemize}

\sphinxAtStartPar
\sphinxincludegraphics{{eeg_psd_50hz}.png}

\item {} 
\sphinxAtStartPar
Examine the PSD plot and then the first section of the Spectrogram. Expand the vertical axes if necessary. Use the waveform cursor to identify the frequency in Hertz (Hz) of the peak in the PSD plot and the band in the Spectrogram. Values are displayed at the top of each plot.

\sphinxAtStartPar
\sphinxstylestrong{What is the frequency in Hertz (Hz) of this sine wave?}

\item {} 
\sphinxAtStartPar
Select the second record and again view the result in the Spectrum view.

\sphinxAtStartPar
\sphinxstylestrong{What is the frequency in Hertz (Hz) of this second sine wave?}

\item {} 
\sphinxAtStartPar
Select the third record and again view the result in the Spectrum view. You should now see two prominent peaks (PSD plot) and bands (Spectrogram) in the result.

\sphinxAtStartPar
\sphinxstylestrong{Are these two peaks/bands the same as for the first two records? How is this possible?}

\item {} 
\sphinxAtStartPar
Select the fourth record and again view the result in the Spectrum view.

\sphinxAtStartPar
\sphinxstylestrong{Is there any regular signal within this record?}

\item {} 
\sphinxAtStartPar
In Chart View compare the signal amplitudes of the fourth and fifth records. Note that the fifth record has lower amplitude compared with the fourth record (you may need to Autoscale to see a peak at all).

\item {} 
\sphinxAtStartPar
Select the fourth record again. In the Spectrum view examine the PSD plot. Move the \sphinxstylestrong{Waveform Cursor} to the prominent peak.

\sphinxAtStartPar
\sphinxstylestrong{What is the frequency (Hz) of this signal?}

\sphinxAtStartPar
\sphinxstylestrong{What is the power (mV²) of this signal?}

\item {} 
\sphinxAtStartPar
Select the fifth record and examine the PSD plot. Move the \sphinxstylestrong{Waveform Cursor} to the peak.

\sphinxAtStartPar
\sphinxstylestrong{What is the frequency (Hz) of this signal?}

\sphinxAtStartPar
\sphinxstylestrong{What is the power (mV²) of this signal?}

\item {} 
\sphinxAtStartPar
Examine the Spectrogram. Note that the band corresponding to the signal’s frequency appears to be missing. This is because the power of the signal is small compared with the previous four records. Expand the scale on the right\sphinxhyphen{}hand side of the Spectrogram by setting the scale closer to the power of the signal (e.g., right\sphinxhyphen{}click \sphinxstylestrong{Set Scale} to 0 to 0.004 V²). Note that the band is now visible at the expected frequency.

\sphinxAtStartPar
\sphinxincludegraphics{{eeg_spectrogram_menu}.png} \sphinxincludegraphics{{eeg_setscale}.png}

\item {} 
\sphinxAtStartPar
The fifth record is the same signal as the fourth record, except that the quality of the raw signal has been affected. Compare your features (amplitude, power, frequency) of the fourth and fifth records.

\sphinxAtStartPar
\sphinxstylestrong{How has the quality of the signal affected the wave features?}

\end{enumerate}


\bigskip\hrule\bigskip



\subsection{Part II. Detecting alpha waves}
\label{\detokenize{EEG:part-ii-detecting-alpha-waves}}

\subsubsection{Equipment Setup and Electrode Attachment}
\label{\detokenize{EEG:equipment-setup-and-electrode-attachment}}\begin{enumerate}
\sphinxsetlistlabels{\arabic}{enumi}{enumii}{}{.}%
\item {} 
\sphinxAtStartPar
Make sure the PowerLab is turned \sphinxstylestrong{off} and the USB cable is connected to the computer.

\item {} 
\sphinxAtStartPar
Connect the 5 Lead Shielded Bio Amp Cable to the Bio Amp Connector on the front panel of the PowerLab (Figure 1). The hardware needs to be connected \sphinxstylestrong{before} you open the settings file.

\end{enumerate}

\sphinxAtStartPar
\sphinxincludegraphics{{eeg_setup}.png}

\sphinxAtStartPar
\sphinxstyleemphasis{Figure 1. Equipment setup for PowerLab 26T. Image courtesy of ADInstruments ©}
\begin{enumerate}
\sphinxsetlistlabels{\arabic}{enumi}{enumii}{}{.}%
\setcounter{enumi}{2}
\item {} 
\sphinxAtStartPar
Attach the leads of the EEG Flat Electrodes to the Earth, CH1 NEG and POS pins closest to the labeled side on the Bio Amp Cable. Channel 1 “\sphinxstylestrong{positive}” will lead to the inion (the bump on the back of the head above the neck) and Channel 1 “\sphinxstylestrong{negative}” will lead to the forehead. Channel 2 will be empty and the Earth will lead to the temple. Refer to Figure 1 for proper placement, but do not attach them to the volunteer yet. Follow the color scheme on the Bio Amp Cable.

\item {} 
\sphinxAtStartPar
Remove any jewelry from the volunteer’s face, ears, and neck.

\item {} 
\sphinxAtStartPar
With Figure 1 as a guide, locate where you will place the electrodes. Abrade the skin with the NuPrep Abrasive Gel.

\sphinxAtStartPar
\sphinxstylestrong{Note:} This is important as abrasion helps reduce the skin’s resistance. Your recordings will be much better if you have abraded the skin well (but not too much — that will cause irritation).

\item {} 
\sphinxAtStartPar
After abrasion, clean the area with an alcohol swab to remove the dead skin cells.

\item {} 
\sphinxAtStartPar
While the skin is drying, scoop Electrode Paste into the EEG Flat Electrodes. When the skin is dry, stick the electrodes to the skin (Figure 1). Immediately hold the electrodes and wires in place with the medical tape. Make sure the tape firmly holds the electrodes against the head. Use the elastic bandage to wrap tightly around the head. This will help the electrodes maintain good contact with the skin.

\item {} 
\sphinxAtStartPar
Have the volunteer lie in a comfortable position on his/her back (or sit comfortably), with the head turned so that none of the electrodes are disturbed or compressed.

\item {} 
\sphinxAtStartPar
Check that all three electrodes are properly connected to the volunteer and the Bio Amp Cable before proceeding. Turn on the PowerLab.

\end{enumerate}


\subsubsection{Exercise 1: Recognizing Artifacts}
\label{\detokenize{EEG:exercise-1-recognizing-artifacts}}
\sphinxAtStartPar
In this exercise, you will learn to recognize some of the artifacts that can appear on an EEG.
\begin{enumerate}
\sphinxsetlistlabels{\arabic}{enumi}{enumii}{}{.}%
\item {} 
\sphinxAtStartPar
Launch LabChart with the EEG Settings file (\sphinxurl{http://bit.ly/labchart}).

\item {} 
\sphinxAtStartPar
Select \sphinxstylestrong{Bio Amp} from the EEG Channel Function pop\sphinxhyphen{}up menu. Make sure the settings are as follows: \sphinxstylestrong{Range 200 µV, High Pass 0.5 Hz, and Low Pass 50 Hz}.
\begin{itemize}
\item {} 
\sphinxAtStartPar
💡 \sphinxstylestrong{Why are we applying high\sphinxhyphen{} and low\sphinxhyphen{}pass filters to our data?}

\end{itemize}

\item {} 
\sphinxAtStartPar
\sphinxstylestrong{Start} recording. Add a \sphinxstylestrong{comment} “blinking,” and have the volunteer blink repeatedly. \sphinxstylestrong{Stop} recording after 10 seconds.

\item {} 
\sphinxAtStartPar
Repeat step 3, this time, have the volunteer make eye movements. Add a \sphinxstylestrong{comment} “eye movements.” Have the volunteer gaze up\sphinxhyphen{}and\sphinxhyphen{}down and left\sphinxhyphen{}and\sphinxhyphen{}right in a repeated pattern. Make sure the volunteer’s head is still and only the eyes move.

\item {} 
\sphinxAtStartPar
Repeat step 3, this time, have the volunteer make head movements. Add a \sphinxstylestrong{comment} “head movements.” Have the volunteer gently move his/her head in a repeated pattern.

\item {} 
\sphinxAtStartPar
Examine the vertical scale at the left of the Chart View, and note the positions corresponding to +50 µV and −50 µV. True EEG signals rarely exceed these limits.

\item {} 
\sphinxAtStartPar
Examine the entire data trace and \sphinxstylestrong{Autoscale}, if necessary. There may be some large signals outside the ±75 µV range. \sphinxstylestrong{Such large signals are artifacts.}

\item {} 
\sphinxAtStartPar
\sphinxstylestrong{Save your data}, and open a new file with the same settings.

\end{enumerate}


\subsubsection{★ Exercise 2: Alpha Waves in the EEG}
\label{\detokenize{EEG:exercise-2-alpha-waves-in-the-eeg}}
\sphinxAtStartPar
In this exercise, you will examine the effects of relaxation and eye movement on alpha waves in the EEG.
\begin{enumerate}
\sphinxsetlistlabels{\arabic}{enumi}{enumii}{}{.}%
\item {} 
\sphinxAtStartPar
Make sure the volunteer is relaxed and comfortable. Have the volunteer close his/her eyes and remain quiet. Keep noise to a minimum and keep all distractions away from the volunteer.

\item {} 
\sphinxAtStartPar
\sphinxstylestrong{Start} recording. Prepare a comment with “open”; do not enter it yet.

\item {} 
\sphinxAtStartPar
Tell the volunteer to open both eyes. Immediately press Return/Enter to add the \sphinxstylestrong{comment}. Record with the volunteer’s eyes open for 30 seconds. Do not stop recording.

\item {} 
\sphinxAtStartPar
Prepare a comment with “shut.” When the 30 seconds are complete, tell the volunteer to close both eyes. Immediately press Return/Enter to add the \sphinxstylestrong{comment}.

\item {} 
\sphinxAtStartPar
Repeat steps 3 and 4 twice, to give you three sets of results. \sphinxstylestrong{Save your data}.

\end{enumerate}

\sphinxAtStartPar
\sphinxincludegraphics{{eeg_eyes_shut_open}.png}
\sphinxstyleemphasis{EEG signal with eyes shut and open. Note alpha waves during eyes shut (left).}


\subsubsection{Wave Amplitude Analysis}
\label{\detokenize{EEG:wave-amplitude-analysis}}\begin{enumerate}
\sphinxsetlistlabels{\arabic}{enumi}{enumii}{}{.}%
\item {} 
\sphinxAtStartPar
Use the \sphinxstylestrong{View Buttons} to change the horizontal compression to see data with eyes open and shut. Make a data selection that includes some data from both eyes open and eyes shut conditions. View this selection in \sphinxstylestrong{Zoom View}. This should make it easier to see the alpha wave activity. \sphinxstylestrong{Autoscale}, if necessary.

\item {} 
\sphinxAtStartPar
In Chart View, scroll through the parts of the recording that were made with the volunteer’s eyes shut to look for alpha waves. Use the \sphinxstylestrong{View Buttons} to change the horizontal compression if necessary. The alpha waves can be recognized by their amplitude (usually 30 to 50 µV peak\sphinxhyphen{}to\sphinxhyphen{}peak) and their frequency. Each cycle of an alpha wave lasts \textasciitilde{}0.1 s.

\item {} 
\sphinxAtStartPar
Use the \sphinxstylestrong{Marker} and \sphinxstylestrong{Waveform Cursor} to measure the amplitude of the alpha waves. Place the \sphinxstylestrong{Marker} at the lowest point of the wave and move the \sphinxstylestrong{Waveform Cursor} to the peak of the wave. Measure the amplitudes of five waves from when the volunteer’s eyes were closed. Record the values in \sphinxstylestrong{Table 1} of the Data Notebook.

\item {} 
\sphinxAtStartPar
Now measure wave amplitudes when the volunteer’s eyes were open. Record these values in \sphinxstylestrong{Table 2} of the Data Notebook.

\end{enumerate}


\subsubsection{Spectral Analysis}
\label{\detokenize{EEG:spectral-analysis}}
\sphinxAtStartPar
Now you will use Spectral Analysis to examine the EEG you recorded.
\begin{enumerate}
\sphinxsetlistlabels{\arabic}{enumi}{enumii}{}{.}%
\item {} 
\sphinxAtStartPar
Find the part of the recording when the volunteer had their eyes \sphinxstyleemphasis{\sphinxstylestrong{shut}}. Click\sphinxhyphen{}and\sphinxhyphen{}drag across this part of the data trace to \sphinxstylestrong{select} it.

\item {} 
\sphinxAtStartPar
From the \sphinxstylestrong{Window} menu, select \sphinxstylestrong{Spectrum}. In the Spectrum View, choose \sphinxstylestrong{Selected} (underneath the main toolbar at the top).

\end{enumerate}

\sphinxAtStartPar
\sphinxincludegraphics{{eeg_spectrum}.png}
\sphinxstyleemphasis{Spectrum of an EEG}
\begin{enumerate}
\sphinxsetlistlabels{\arabic}{enumi}{enumii}{}{.}%
\setcounter{enumi}{2}
\item {} 
\sphinxAtStartPar
Next to Display Channels, choose Ch 3: EEG.

\item {} 
\sphinxAtStartPar
Alpha activity should show up in the PSD plot as a clear peak in the 8–12 Hz range.

\item {} 
\sphinxAtStartPar
Alpha activity shows up in the Spectrogram as a band of color in the 8–12 Hz range. If you cannot see the alpha activity as a clear peak in the 8–12 Hz range, manually scale the horizontal and vertical axes. Note that Spectrogram displays all the recorded data and that the selection you have made is highlighted in a darker blue color.

\item {} 
\sphinxAtStartPar
Make a data selection of several seconds from when the volunteer had their eyes \sphinxstyleemphasis{\sphinxstylestrong{open}}. Select \sphinxstylestrong{Spectrum}. Note that in the PSD plot the peak in the alpha activity range of 8–12 Hz is small or absent, and in the Spectrogram the band of color in the alpha activity range of 8–12 Hz is weak or absent.

\item {} 
\sphinxAtStartPar
In the Spectrogram, scale the horizontal axis so that all the data is visible. Note the presence and absence of the band of color in the alpha activity range of 8–12 Hz, which corresponds with the eyes shut and eyes open conditions.

\item {} 
\sphinxAtStartPar
If you’re happy with your spectrograms (e.g., you can clearly see alpha activity in closed conditions but not open conditions), you’ve completed the task for this lab. Save your data in a safe location (we’ll deal with plotting it later).

\end{enumerate}


\bigskip\hrule\bigskip



\subsection{Data Notebook}
\label{\detokenize{EEG:data-notebook}}

\subsubsection{Table 1. Alpha Waves with Eyes Open}
\label{\detokenize{EEG:table-1-alpha-waves-with-eyes-open}}

\begin{savenotes}\sphinxattablestart
\sphinxthistablewithglobalstyle
\centering
\begin{tabulary}{\linewidth}[t]{TT}
\sphinxtoprule

\sphinxAtStartPar

&\sphinxstyletheadfamily 
\sphinxAtStartPar
Amplitude (µV)
\\
\sphinxmidrule
\sphinxtableatstartofbodyhook
\sphinxAtStartPar
\sphinxstylestrong{Wave \#1}
&
\sphinxAtStartPar

\\
\sphinxhline
\sphinxAtStartPar
\sphinxstylestrong{Wave \#2}
&
\sphinxAtStartPar

\\
\sphinxhline
\sphinxAtStartPar
\sphinxstylestrong{Wave \#3}
&
\sphinxAtStartPar

\\
\sphinxhline
\sphinxAtStartPar
\sphinxstylestrong{Wave \#4}
&
\sphinxAtStartPar

\\
\sphinxhline
\sphinxAtStartPar
\sphinxstylestrong{Wave \#5}
&
\sphinxAtStartPar

\\
\sphinxbottomrule
\end{tabulary}
\sphinxtableafterendhook\par
\sphinxattableend\end{savenotes}


\subsubsection{Table 2. Alpha Waves with Eyes Shut}
\label{\detokenize{EEG:table-2-alpha-waves-with-eyes-shut}}

\begin{savenotes}\sphinxattablestart
\sphinxthistablewithglobalstyle
\centering
\begin{tabulary}{\linewidth}[t]{TT}
\sphinxtoprule

\sphinxAtStartPar

&\sphinxstyletheadfamily 
\sphinxAtStartPar
Amplitude (µV)
\\
\sphinxmidrule
\sphinxtableatstartofbodyhook
\sphinxAtStartPar
\sphinxstylestrong{Wave \#1}
&
\sphinxAtStartPar

\\
\sphinxhline
\sphinxAtStartPar
\sphinxstylestrong{Wave \#2}
&
\sphinxAtStartPar

\\
\sphinxhline
\sphinxAtStartPar
\sphinxstylestrong{Wave \#3}
&
\sphinxAtStartPar

\\
\sphinxhline
\sphinxAtStartPar
\sphinxstylestrong{Wave \#4}
&
\sphinxAtStartPar

\\
\sphinxhline
\sphinxAtStartPar
\sphinxstylestrong{Wave \#5}
&
\sphinxAtStartPar

\\
\sphinxbottomrule
\end{tabulary}
\sphinxtableafterendhook\par
\sphinxattableend\end{savenotes}


\subsubsection{Questions for consideration}
\label{\detokenize{EEG:questions-for-consideration}}
\sphinxAtStartPar
\sphinxstylestrong{Under what conditions did you see alpha waves more clearly?}

\sphinxAtStartPar
\sphinxstylestrong{What are alpha waves thought to indicate?}

\sphinxAtStartPar
\sphinxstylestrong{Look at your Spectral Analysis from when the volunteer had their eyes closed. At what range do you have the greatest peak? What type of brain waves does the Spectral Analysis suggest is active?}


\bigskip\hrule\bigskip



\subsection{★ Part III. Run your own experiment}
\label{\detokenize{EEG:part-iii-run-your-own-experiment}}
\sphinxAtStartPar
In the final portion of the lab, you will run your own EEG experiment. Choose two different conditions for your subject — e.g., listening to music (or different types of music), doing something cognitively intensive like math problems, or meditating/napping. Choose a brain wave (e.g., alpha, beta, or gamma) that you’d like to analyze, and use the Spectrum view to see whether it is different across conditions. You should conduct as many repeats as you have time for. Take notes about your methods — you’ll need these for your lab report!

\sphinxAtStartPar
\sphinxstylestrong{Note:} If you do not have time to run this on the first EEG day, you will have time during the subsequent lab period. At minimum, you should decide what you will do for this experiment.


\subsection{★ Analyzing your EEG data}
\label{\detokenize{EEG:analyzing-your-eeg-data}}

\subsubsection{Exporting PSD data from LabChart}
\label{\detokenize{EEG:exporting-psd-data-from-labchart}}
\sphinxAtStartPar
The PSD chart in LabChart will give you the quantitative data you need for your lab report. For example, you can compare the power of alpha between eyes open and closed. You should do the same for your own experiment.
\begin{enumerate}
\sphinxsetlistlabels{\arabic}{enumi}{enumii}{}{.}%
\item {} 
\sphinxAtStartPar
In LabChart, with your data selected in Spectrum View, go to \sphinxstylestrong{File > Export.}

\item {} 
\sphinxAtStartPar
Save your file as a \sphinxstylestrong{Spectrum PSD Text File}.

\item {} 
\sphinxAtStartPar
You can import this txt file into Excel (or MATLAB), similar to the way you import raw recording data (see \sphinxhref{https://docs.google.com/document/d/1nscGwk4ZoLGoE1HhwHiOv8XimzUQ80RTLfwSNeMRO6E/edit}{this document} for a refresher).

\end{enumerate}


\subsubsection{Exporting spectrogram data from LabChart}
\label{\detokenize{EEG:exporting-spectrogram-data-from-labchart}}
\sphinxAtStartPar
\sphinxincludegraphics{{eeg_display_settings}.png}
\begin{enumerate}
\sphinxsetlistlabels{\arabic}{enumi}{enumii}{}{.}%
\item {} 
\sphinxAtStartPar
In LabChart 8, go to \sphinxstylestrong{Setup > Display Settings}. Make sure that “Always seconds” is checked (see example at right).

\item {} 
\sphinxAtStartPar
With Spectrum view open in LabChart, select the data you’d like to export.

\item {} 
\sphinxAtStartPar
Save your file as a \sphinxstylestrong{Spectrogram Text File}. Save your file with a name that \sphinxstyleemphasis{\sphinxstylestrong{does not}} have spaces.

\end{enumerate}

\sphinxAtStartPar
\sphinxincludegraphics{{eeg_spectrogram_export}.png}


\subsubsection{Plotting Spectrogram data in Python (on the DataHub)}
\label{\detokenize{EEG:plotting-spectrogram-data-in-python-on-the-datahub}}
\sphinxAtStartPar
The following instructions will help you plot a spectrogram of your text file output from LabChart.

\sphinxAtStartPar
\sphinxstylestrong{Update your code}
\begin{enumerate}
\sphinxsetlistlabels{\arabic}{enumi}{enumii}{}{.}%
\item {} 
\sphinxAtStartPar
Log on to the DataHub using this link to get the most recent version of our Utilities folder: \sphinxurl{https://datahub.ucsd.edu/hub/user-redirect/git-sync?repo=https://github.com/BIPN145/Utilities}

\item {} 
\sphinxAtStartPar
Open the container for BIPN 145.

\end{enumerate}

\sphinxAtStartPar
\sphinxstylestrong{Upload your data to the DataHub}
\begin{enumerate}
\sphinxsetlistlabels{\arabic}{enumi}{enumii}{}{.}%
\item {} 
\sphinxAtStartPar
In the Jupyter Dashboard, navigate into your BIPN145 folder.

\item {} 
\sphinxAtStartPar
In the upper right\sphinxhyphen{}hand corner, choose the upload button (\sphinxincludegraphics{{datahub_upload_button}.png}) to find your exported .txt file and upload it.

\item {} 
\sphinxAtStartPar
The file will show up on the top of your file list, with a blue upload button (\sphinxincludegraphics{{datahub_blue_upload}.png}). Click the button to upload.

\end{enumerate}


\subsubsection{Open the Notebook to plot your data}
\label{\detokenize{EEG:open-the-notebook-to-plot-your-data}}\begin{enumerate}
\sphinxsetlistlabels{\arabic}{enumi}{enumii}{}{.}%
\item {} 
\sphinxAtStartPar
Find \& open \sphinxcode{\sphinxupquote{plotSpectrogram.ipynb}} in your BIPN145 folder.

\item {} 
\sphinxAtStartPar
Change the file name in the first code block to your file.

\item {} 
\sphinxAtStartPar
Work through the notebook to clean up the data and plot your file.

\item {} 
\sphinxAtStartPar
Right\sphinxhyphen{}click on the image and save.

\end{enumerate}







\renewcommand{\indexname}{Index}
\printindex
\end{document}